\documentclass[11pt]{article}


\usepackage[inline]{asymptote}
\usepackage{asypictureB}

\usepackage{qrcode}
%\usepackage[default]{frcursive}
\usepackage[a4paper,text={17.5cm,25cm},centering,top=2.5cm]{geometry} 
\usepackage{amsfonts}
\usepackage[colorlinks=true,linkcolor=black]{hyperref}
%\usepackage{amsmath}
\usepackage{amssymb,amsmath,stackrel,wasysym,mathrsfs}
\usepackage{multicol}
\usepackage{pdflscape}% girar una pagina, hay que usar \begin{landscape}En este punto inicia la hoja en vertical \end{landscape}
%\usepackage[mathcal]{euscript}
%\usepackage{amssymb}
\usepackage{esvect}%para la flecha de vector 

\usepackage[T1]{fontenc}
\usepackage[spanish,es-noshorthands]{babel}
%\usepackage[latin1,ansinew]{inputenc}\usepackage{calligra}
%\usepackage[pdftex]{graphicx}
%\usepackage{etex}
\usepackage{calligra}
\usepackage{array}
 \usepackage{booktabs}
 \usepackage{ifsym}
\usepackage{pgf,tikz,pgfplots}
\usepackage{tkz-graph}%paquete (di)grafos
\usepackage{tkz-euclide}% Para realizar diagramas que coloreen conjuntos
\usetikzlibrary[decorations.text]% Decorar curvas
%\usetikzlibrary{matrix}
\usetikzlibrary{arrows,matrix,calc}
%\usetkzobj{all}

\pgfplotsset{compat=1.15}
\usepackage{mathrsfs}
\usetikzlibrary{arrows}
\usepackage{makecell} %Paquete para controlar altura de celdas
\usepackage{verbatim} 
\usetikzlibrary{babel}%para no general error al usar babel y tikz simultaneamente.
\usetikzlibrary{patterns,patterns.meta}%%%%%%%%%%%para  pintar de forma rrallada un rectangulo

\usepackage{pictex}
\usepackage{verbatim}%me sirve para usar comentarios con mucha lineas 
\usepackage{amstext}
\usepackage{enumerate}
%\usepackage{enumitem}
\usepackage{amsthm}
\usepackage{xcolor}
\usepackage{tcolorbox}
\usepackage{mathdots}% para tener mas opciones de puntos suspensivos

%\usepackage{pgf,tikz}%,pgfplots
%\pgfplotsset{compat=1.15}
\usepackage{mathrsfs}
%\usetikzlibrary{arrows}
%\usepackage{etex,tkz-euclide} 
%\usetkzobj{all}
\usepackage{color}

\usepackage{cancel}
\newcommand{\Cancelar}[2][black]{{\color{#1}\cancel{\color{black}#2}}}%%% para cancelar con colores recordar usar mayusculas
\usepackage{lscape}%paquete para girar hojas












\newcommand{\R}{\mathbb{R}}
\newcommand{\re}{\mathbb{R}}
\newcommand{\I}{\mathbb{I}}
\newcommand{\Q}{\mathbb{Q}}
\newcommand{\Z}{\mathbb{Z}}
\newcommand{\C}{\mathbb{C}}
\newcommand{\N}{\mathbb{N}}
\newcommand{\K}{\mathbb{K}}
%\newcommand{\B}{\mathcal{B}}
\newcommand{\V}{\mathbb{V}}
\newcommand{\can}{\mathcal{C}}
\newcommand{\M}{\mathbb{M}}
\newcommand{\Pol}{\mathcal{P}}
\newcommand{\moduloprop}{\hookrightarrow_{\hspace{-2.5mm}\mbox{\tiny$^{_\neq}$}}}

\newcommand{\dis}{\displaystyle}



\newtheorem{ejer}{Ejercicio}[]

\newtheorem{teo}{Teorema}%[section]
\newtheorem{lema}[teo]{Lema}
\newtheorem{propo}[teo]{Proposición}
\newtheorem{defi}[teo]{Definición}
\newtheorem{ejem}[teo]{}
\newtheorem*{ejemsinn}{}
\newtheorem{obs}[teo]{Observación}%[section]
\newtheorem{acla}{Aclaración}%[section]
\newtheorem{act}{Actividad}%%%%%%%%%%%%%%%%
\newtheorem*{defsinn}{{\bf Definición}:}
\newtheorem*{ejemplosinn}{{\bf Ejemplo}:}
\newtheorem*{proposinn}{{\bf\sc Proposición}}
\newtheorem*{obssinn}{Observación:}
\newtheorem*{corosinn}{Corolario}
\newtheorem*{propsinn}{Propiedades:}
%\newenvironment{proof} { \noindent {\bf Demostración: }\rm}{$\quad \hfill \blacksquare $}
\usepackage{setspace} % (es una alternativa a \renewcommand{\baselinestretch}{1.3}) usando \doublespacing \onehalfspace \singlespace \spacing{1.5}  se va cambiando segun uno quiera

%\newenvironment{proof}[Demostración:]
\newenvironment{sol}{ \noindent {\bf Solución:} \rm}{%\hfill \protect\tikz[overlay,yshift=-4pt]\protect\node[xshift=8mm, yshift=13mm, blue]() {\huge $\checkmark$ };
}

%%%%%%%%%%%%%%%%%%%%%%%%%%%%%%%%%%%%%%%%


%\renewcommand{\baselinestretch}{1.5}
%%%%%%%%%%%%%%%%%%%%%%%%%%%%%%%%%%%%%%%%%%%%%%%%%%%%%%%%%%%%%%%%%%%%%%%%%%%%%%%%%%%%%%%%%%%%%%%%%%%%%%%%%%%%%%%%%%%%%%%%%%%%%%%%%%%%%%%%%%%%%%%%%%%%%%%%%%%%%%%%%%%%%%%%%%%%%%%%%%%%%%%%%%%%%%%%%%%%%%%%%%%%%%%%%%%%%%%%%%%%%%%%%%%%%%%%%%%%%%%%%%%%%%%%%%%%%%%%%%%%%%%%%%%%%%%%%%%%%%%%
\usepackage[framemethod=TikZ]{mdframed}%Paquete para estilos de los bordes de los ambientes Importante este paquete (mdframed) define un entorno  que trata automáticamente los saltos de página en el texto enmarcado.
            \mdfsetup{skipabove=1\baselineskip, skipbelow=1.5\baselineskip ,             frametitlefont= \sffamily\bfseries , needspace=4\baselineskip , splittopskip=1.5\baselineskip} %
            \mdfsetup{apptotikzsetting={\tikzset{mdfbackground/.append style={draw=none}}}} 
%%%%%%%%%%%%%%%%%%%%%%%%%%%%%%%%%%%%%%%%%%%%%%%%%%%%%%%%%%%%%%
%
%Estilo para Cuentas auxiliares
%
%%%%%%%%%%%%%%%%%%%%%%%%%%%%%%%%%%%%%%%%%%%%%%%%%%%%%%%%%%%%%%
\mdfdefinestyle{Cuentaauxi}{%skipabove=10pt, skipbelow=0,
userdefinedwidth=1.02\textwidth,align=center,
  everyline=true,
  ignorelastdescenders=true,
  middlelinewidth=1pt,
  linecolor=black!20,
  outerlinewidth=1pt,
  %leftline=false,topline=false
  %rightline=false,bottomline=false,
  %backgroundcolor=black!20,
  innerleftmargin=2mm, innerrightmargin=2mm,
  innerbottommargin=3mm,
  innertopmargin=1mm,
  frametitleaboveskip=2mm,
  frametitlebelowskip=2mm,
  %frametitlerule=false,
  %repeatframetitle=false,
 % outerlinewidth=6pt,outerlinecolor=blue!20,
 pstricksappsetting={\addtopsstyle{mdfouterlinestyle}{linestyle=dashed}},
 %firstextra={\node[xshift=5cm,right,draw=White, line width=1.5pt,%star,star points=10,star point ratio=1.2, 
 %minimum size=15mm, fill=white] at (P-|O) {\tikzpicture\draw[draw=black!20,decoration=snake, fill=white] (0,0) --(2.5,0) decorate{--(2.5,1)}--(0,1)decorate{--cycle};%\includegraphics[width=10mm]{.pdf}
 %  \endtikzpicture
 %}; },
 %singleextra={\node[xshift=5cm,right,draw=White, line width=1.5pt,%star,star points=10,star point ratio=1.2, minimum size=15mm,
  %fill=white] at (P-|O) {\tikzpicture\draw[draw=black!20,decoration=snake, fill=white] (0,0) --(2.5,0) decorate{--(2.5,1)}--(0,1)decorate{--cycle};%\includegraphics[width=10mm]{.pdf}    
 %  \endtikzpicture
 %}; }
}
%%%%%%%%%%%%%%%%%%%%%%%%%%%%%%%%%%%%%%%%%%%%%%%%%%%%%%%%%%%%%%
%
%Estilo para Citar
%
%%%%%%%%%%%%%%%%%%%%%%%%%%%%%%%%%%%%%%%%%%%%%%%%%%%%%%%%%%%%%%
\mdfdefinestyle{citar}{
userdefinedwidth=.83\textwidth ,align=center,
skipabove=5pt, skipbelow=5pt,
  everyline=true,
  ignorelastdescenders=true,
  middlelinewidth=1pt,
  linecolor=blue!60,
  outerlinewidth=2pt,
  %leftline=false,
  topline=false,
  rightline=false,
  bottomline=false,
  %backgroundcolor=black!20,
  innerleftmargin=5mm, innerrightmargin=5mm,
  innerbottommargin=3mm,
  innertopmargin=1mm,
  leftmargin= 14mm,
  %frametitleaboveskip=2mm,
  %frametitlebelowskip=2mm,
  %frametitlerule=false,
  %repeatframetitle=false,
 % outerlinewidth=6pt,outerlinecolor=blue!20,
 pstricksappsetting={\addtopsstyle{mdfouterlinestyle}{linestyle=dashed}},
 %firstextra={\node[xshift=5cm,right,draw=White, line width=1.5pt,%star,star points=10,star point ratio=1.2, 
 %minimum size=15mm, fill=white] at (P-|O) {\tikzpicture\draw[draw=black!20,decoration=snake, fill=white] (0,0) --(2.5,0) decorate{--(2.5,1)}--(0,1)decorate{--cycle};%\includegraphics[width=10mm]{.pdf}
 %  \endtikzpicture
 %}; },
 %singleextra={\node[xshift=5cm,right,draw=White, line width=1.5pt,%star,star points=10,star point ratio=1.2, minimum size=15mm,
  %fill=white] at (P-|O) {\tikzpicture\draw[draw=black!20,decoration=snake, fill=white] (0,0) --(2.5,0) decorate{--(2.5,1)}--(0,1)decorate{--cycle};%\includegraphics[width=10mm]{.pdf}    
 %  \endtikzpicture
 %}; }
}
%%%%%%%%%%%%%%%%%%%%%%%%%%%%%%%%%%%%%%%%%%%%%%%%%%%%%%%%%%%%%%%%%%%%%%%%%%%%%%%%%%%%%%%%%%%%%%%%%%%%%%%%%%%%%%%%%%%%%%%%%%%%%%%%%%%%%%%%%%%%%%%%%%%%%%%%%%%%%%%%%%%%%%%%%%%%%%%%%%%%%%%%
%%%%%%%%%%%%%%%%%%%%%%%%%%%%%%%%%%%%%%%%%%%%%%%%%%%%%%%%%%%%%%%%%%%%%%%%%%%%%%%%%%%%%%%%%%%%%%%%%%%%%%%%%%%%%%%%%%%%%%%%%%%%%%%%%%%%%%%%%%%%%%%%%%%%%%%%%%%%%%%%%%%%%%%%%%%%%%%%%%%%%%%%%%%%%%%%%%%%%%%%%%%%%%%%%%%%%%%%%%%%%%%%%%%%%%
%%%%%%%%%%%%%%%%%%%%%%%%%%%%%%%%%%%%%%%%%%%%%%%%%%%%%%%%%%%%%%%%%%%%%%%%%%%%%%%%%%%%%%%%%%%%%%%%%%%%%%%%%%%%%%%%%%%%%%%%%%%%%%%%%%%%%%%%%%%%%%%%%%%%%%%%%%%%%%%%%%%%%%%%%%%%%%%%%%%%%%%%%%%%%%%%%%%%%%%%%%
\newcommand{\Dem}{\protect\tikz[overlay,yshift=-4pt]\protect\draw[line width=1.5pt,line cap=round,color=gray] (0,0)[out=14,in=180] to node [shift=(90:6.pt),color=black] {\large\bf Dem:} (7:1); \qquad \quad }
\DeclareRobustCommand{\michi}{{\mathpalette\irchi\relax}}
\newcommand{\irchi}[1]{\raisebox{\depth}{\large$#1\chi$}} % inner command, used by \rchi



%%%%%%%%%%%%%%%%%%%%%%%%%%%%%%%%%%%%%%%%%%%%%%%%%%%%%%%%%%%%%%%%%%%%%%%%%%%%%%%%%%%%%%%%%%%%%%%%%%%%%%%%%%%%%%%%%%%%%%%%%%%%%%%%%%%%%%%%%%%%%%%%%%%%%%%%%%%%%%%%%%%%%%%%%%%%%%%%%%%%%%%%%%%%%%%%%%%%%%%%%%%%%%%%%%%%%%%%%%%%%%%%%%%%%%%%%%%%%%%%%%%%%%%%%%%%%%%%%%%%%%%%%%%%%%%%%%%%%%%%%%%%%%%%%%%%%%%%%%%%%%%%%%%%%%%%%%%%%%%%%%%%%%%%%%%%%%%%%%%%%%%%%%%%%%%%%%%%%%%%%%%%%%%%%%%%%%%%%%%%%%%%%%%%%%%%%%%%%%%%%%%%%%%%%%%%%%%%%%%%%%%%%%%%%%%%%%%%%%%%%%%%%%%%%%%%%%%%%%%%%%%%%%%%%%%%%%%%%%%%%%%%%%%%%%%%%%%%%%%%%%%%%%%%%%%%
% comando aclaraciones  

\newcommand{\aclaracion}[2]{

\begin{center}
\definecolor{gris}{rgb}{0.5,0.5,0.5}
\begin{tikzpicture}[line cap=round,line join=round,scale=1]
%\draw [color=gris,dash pattern=on 1pt off 1pt, xstep=1cm,ystep=1cm] (0.1,0.1) grid (6.3,6.3);
%\draw[color=black] (0.,0.) -- (\linewidth,0.);
\node at (.23,{.9*#2}) [color=black,above right] {\bf #1};
\foreach \y in {1,2,...,#2}
\draw[shift={(0,{.9*\y})},color=gris,opacity=.5] (0.,0.) -- ({.99*\linewidth},0.);
\end{tikzpicture} 
\end{center}
}





%%%%%%%%%%%%%%%%%%%%%%%%%%%%%%%%%%%%%%%%%%%%%%%%%%%%%%%%%%%%%%%%%%%%%%%%%%%%%%%%%%%%%%%%%%%%%%%%%%%%%%%%%%%%%%%%%%%%%%%%%%%%%%%%%%%%%%%%%%%%%%%%%%%%%%%%%%%%%%%%%%%%%%%%%%%%%%%%%%%%%%%%%%%%%%%%%%%%%%%%%%%%%%%%%%%%%%%%%%%%%%%%%%%%%%%%%%%%%%%%%%%%%%%%%%%%%%%%%%%%%%%%%%%%%%%%%%%%%%%%%%%%%%%%%%%%%%%%%%%%%%%%%%%%%%%%%%%%%%%%%%%%%%%%%%%%%%%%%%%%%%%%%%%%%%%%%%%%%%%%%%%%%%%%%%%%%%%%%%%%%%%%%%%%%%%%%%%%%%%%%%%%%%%%%%%%%%%%%%%%%%%%%%%%%%%%%%%%%%%%%%%%%%%%%%%%%%%%%%%%%%%%%%%%%%%%%%%%%%%%%%%%%%%%%%%%%%%%%%%%%%%%%%%%%%%%
% comando respuesta 

\newcommand{\res}[1]{

\begin{center}
\definecolor{gris}{rgb}{0.5,0.5,0.5}
\begin{tikzpicture}[line cap=round,line join=round,scale=1]
%\draw [color=gris,dash pattern=on 1pt off 1pt, xstep=1cm,ystep=1cm] (0.1,0.1) grid (6.3,6.3);
%\draw[color=black] (0.,0.) -- (\linewidth,0.);
\node at (.23,{.9*#1}) [color=black,above right] {\bf respuesta:};
\foreach \y in {1,2,...,#1}
\draw[shift={(0,{.9*\y})},color=gris,opacity=.5] (0.,0.) -- ({.99*\linewidth},0.);
\end{tikzpicture} 
\end{center}
}

%%%%%%%%%%%%%%%%%%%%%%%%%%%%%%%%%%%%%%%%%%%%%%%%%%%%%%%%%%%%%%%%%%%%%%%%%%%%%%%%%%%%%%%%%%%%%%%%%%%%%%%%%%%%%%%%%%%%%%%%%%%%%%%%%%%%%%%%%%%%%%%%%%%%%%%%%%%%%%%%%%%%%%%%%%%%%%%%%%%%%%%%%%%%%%%%%%%%%%%%%%%%%%%%%%%%%%%%%%%%%%%%%%%%%%%%%%%%%%%%%%%%%%%%%%%%%%%%%%%%%%%%%%%%%%%%%%%%%%%%%%%%%%%%%%%%%%%%%%%%%%%%%%%%%%%%%%%%%%%%%%%%%%%%%%%%%%%%%%%%%%%%%%%%%%%%%%%%%%%%%%%%%%%%%%%%%%%%%%%%%%%%%%%%%%%%%%%%%%%%%%%%%%%%%%%%%%%%%%%%%%%%%%%%%%%%%%%%%%%%%%%%%%%%%%%%%%%%%%%%%%%%%%%%%%%%%%%%%%%%%%%%%%%%%%%%%%%%%%%%%%%%%%%%%%%%
% comando ejemplo para completar 

\newcommand{\ejparacompletar}[1]{

\begin{center}
\definecolor{gris}{rgb}{0.5,0.5,0.5}
\begin{tikzpicture}[line cap=round,line join=round,scale=1]
%\draw [color=gris,dash pattern=on 1pt off 1pt, xstep=1cm,ystep=1cm] (0.1,0.1) grid (6.3,6.3);
%\draw[color=black] (0.,0.) -- (\linewidth,0.);
\node at (.23,{.9*#1}) [color=black,above right] {\bf Ejemplo:};
\foreach \y in {1,2,...,#1}
\draw[shift={(0,{.9*\y})},color=gris,opacity=.5] (0.,0.) -- ({.99*\linewidth},0.);
\end{tikzpicture} 
\end{center}
}


%%%%%%%%%%%%%%%%%%%%%%%%%%%%%%%%%%%%%%%%%%%%%%%%%%%%%%%%%%%%%%%%%%%%%%%%%%%%%%%%%%%%%%%%%%%%%%%%%%%%%%%%%%%%%%%%%%%%%%%%%%%%%%%%%%%%%%%%%%%%%%%%%%%%%%%%%%%%%%%%%%%%%%%%%%%%%%%%%%%%%%%%%%%%%%%%%%%%%%%%%%%%%%%%%%%%%%%%%%%%%%%%%%%%%%%%%%%%%%%%%%%%%%%%%%%%%%%%%%%%%%%%%%%%%%%%%%%%%%%%%%%%%%%%%%%%%%%%%%%%%%%%%%%%%%%%%%%%%%%%%%%%%%%%%%%%%%%%%%%%%%%%%%%%%%%%%%%%%%%%%%%%%%%%%%%%%%%%%%%%%%%%%%%%%%%%%%%%%%%%%%%%%%%%%%%%%%%%%%%%%%%%%%%%%%%%%%%%%%%%%%%%%%%%%%%%%%%%%%%%%%%%%%%%%%%%%%%%%%%%%%%%%%%%%%%%%%%%%%%%%%%%%%%%%%%%
% comando conclusión 

\newcommand{\conclus}[1]{

\begin{center}
\definecolor{gris}{rgb}{0.5,0.5,0.5}
\begin{tikzpicture}[line cap=round,line join=round,scale=1]
%\draw [color=gris,dash pattern=on 1pt off 1pt, xstep=1cm,ystep=1cm] (0.1,0.1) grid (6.3,6.3);
%\draw[color=black] (0.,0.) -- (\linewidth,0.);
\node at (.23,{.9*#1}) [color=black,above right] {\bf Conclusión:};
\foreach \y in {1,2,...,#1}
\draw[shift={(0,{.9*\y})},color=gris,opacity=.5] (0.,0.) -- ({.99*\linewidth},0.);
\end{tikzpicture} 
\end{center}
}

%%%%%%%%%%%%%%%%%%%%%%%%%%%%%%%%%%%%%%%%%%%%%%%%%%%%%%%%%%%%%%%%%%%%%%%%%%%%%%%%%%%%%%%%%%%%%%%%%%%%%%%%%%%%%%%%%%%%%%%%%%%%%%%%%%%%%%%%%%%%%%%%%%%%%%%%%%%%%%%%%%%%%%%%%%%%%%%%%%%%%%%%%%%%%%%%%%%%%%%%%%%%%%%%%%%%%%%%%%%%%%%%%%%%%%%%%%%%%%%%%%%%%%%%%%%%%%%%%%%%%%%%%%%%%%%%%%%%%%%%%%%%%%%%%%%%%%%%%%%%%%%%%%%%%%%%%%%%%%%%%%%%%%%%%%%%%%%%%%%%%%%%%%%%%%%%%%%%%%%%%%%%%%%%%%%%%%%%%%%%%%%%%%%%%%%%%%%%%%%%%%%%%%%%%%%%%%%%%%%%%%%%%%%%%%%%%%%%%%%%%%%%%%%%%%%%%%%%%%%%%%%%%%%%%%%%%%%%%%%%%%%%%%%%%%%%%%%%%%%%%%%%%%%%%%%%
% boton calculadora 

\newcommand{\boton}[2]{
\tikzpicture{  \node[draw,color=#1,fill=#1,text=white,,rounded corners=3mm,text width=7mm,text height =3mm,align=center,midway]  at (0,0) { \large \bf #2 };
}
\endtikzpicture
}

\usepackage[framemethod=TikZ]{mdframed}%Paquete para estilos de los bordes de los ambientes Importante este paquete (mdframed) define un entorno  que trata automáticamente los saltos de página en el texto enmarcado.
            \mdfsetup{skipabove=2\baselineskip,skipbelow=2\baselineskip,frametitlefont=\sffamily\bfseries, needspace=4\baselineskip, splittopskip=1.5\baselineskip} 
            \mdfsetup{apptotikzsetting={\tikzset{mdfbackground/.append style={draw=none}}}} 
            
            
%%%%%%%%%%%%%%%%%%%%%%%%%%%%%%%%%%%%%%%%%%%%%%%%%%%%%%%%%%%%%%
%
%Estilo  para cajas
%
%%%%%%%%%%%%%%%%%%%%%%%%%%%%%%%%%%%%%%%%%%%%%%%%%%%%%%%%%%%%%%

\mdfdefinestyle{teoSty}{backgroundcolor=blue!10, linecolor=blue, linewidth=3pt,
skipabove=4pt,
skipbelow=3pt, 
innerleftmargin=3ex, innerrightmargin=3ex, innertopmargin=3mm, innerbottommargin=3mm, innermargin=15pt, outermargin=-15pt, needspace={0.2\textheight},roundcorner=10pt}

\mdfdefinestyle{ObsSty}{backgroundcolor=red!10, linecolor=red, linewidth=2pt,
skipabove=4pt,
skipbelow=3pt, 
innerleftmargin=3ex, innerrightmargin=3ex, innertopmargin=3mm, innerbottommargin=3mm, innermargin=5pt, outermargin=-5pt, needspace={0.2\textheight},roundcorner=10pt}

\mdfdefinestyle{EjerSty}{backgroundcolor=blue!5!green!5!, linecolor=blue!50!green!50!, linewidth=2pt,
skipabove=4pt,
skipbelow=3pt, 
innerleftmargin=3ex, innerrightmargin=3ex, innertopmargin=3mm, innerbottommargin=3mm, innermargin=5pt, outermargin=-5pt, needspace={0.2\textheight},roundcorner=10pt}

 
            
%%%%%%%%%%%%%%%%%%%%%%%%%%%%%%%%%%%%%%%%%%%%%%%%%%%%%%%%%%%%%%
%
%Estilo  para cajas
%
%%%%%%%%%%%%%%%%%%%%%%%%%%%%%%%%%%%%%%%%%%%%%%%%%%%%%%%%%%%%%%

\mdfdefinestyle{teoSty}{backgroundcolor=blue!10, linecolor=blue, linewidth=3pt,
skipabove=4pt,
skipbelow=3pt, 
innerleftmargin=3ex, innerrightmargin=3ex, innertopmargin=3mm, innerbottommargin=3mm, innermargin=15pt, outermargin=-15pt, needspace={0.2\textheight},roundcorner=10pt}

\mdfdefinestyle{ObsSty}{backgroundcolor=red!10, linecolor=red, linewidth=2pt,
skipabove=4pt,
skipbelow=3pt, 
innerleftmargin=3ex, innerrightmargin=3ex, innertopmargin=3mm, innerbottommargin=3mm, innermargin=5pt, outermargin=-5pt, needspace={0.2\textheight},roundcorner=10pt}

\mdfdefinestyle{EjerSty}{backgroundcolor=blue!5!green!5!, linecolor=blue!50!green!50!, linewidth=2pt,
skipabove=4pt,
skipbelow=3pt, 
innerleftmargin=3ex, innerrightmargin=3ex, innertopmargin=3mm, innerbottommargin=3mm, innermargin=5pt, outermargin=-5pt, needspace={0.2\textheight},roundcorner=10pt}



%%%%%%%%%%%%%%%%%%%%%%%%%%%%%%%%%%%%%%%%%%%%%%%%%%%%%%%%%%%%%%
%
%Estilo  para lineas de relleno usando el paquete patterns y la libreria  patterns.meta
%
%%%%%%%%%%%%%%%%%%%%%%%%%%%%%%%%%%%%%%%%%%%%%%%%%%%%%%%%%%%%%%

\tikzdeclarepattern{
  name=mylines,
  parameters={
      \pgfkeysvalueof{/pgf/pattern keys/size},
      \pgfkeysvalueof{/pgf/pattern keys/angle},
      \pgfkeysvalueof{/pgf/pattern keys/line width},
  },
  bounding box={
    (0,-0.5*\pgfkeysvalueof{/pgf/pattern keys/line width}) and
    (\pgfkeysvalueof{/pgf/pattern keys/size},
0.5*\pgfkeysvalueof{/pgf/pattern keys/line width})},
  tile size={(\pgfkeysvalueof{/pgf/pattern keys/size},
\pgfkeysvalueof{/pgf/pattern keys/size})},
  tile transformation={rotate=\pgfkeysvalueof{/pgf/pattern keys/angle}},
  defaults={
    size/.initial=5pt,
    angle/.initial=45,
    line width/.initial=.4pt,
  },
  code={
      \draw [line width=\pgfkeysvalueof{/pgf/pattern keys/line width}]
        (0,0) -- (\pgfkeysvalueof{/pgf/pattern keys/size},0);
  },
}





\newcommand{\semejante}[1]{\stackrel[#1]{}{\longleftrightarrow}}% \thicksim para decir la Semejanza de matrices la variable #1  indica la operacion elemental 
\newcommand{\OptipoUno}[1]{\stackrel[#1]{}{=}}



\newcommand{\indicador}[4]{
{\tikz[remember picture,overlay]\coordinate (#1) at #2;}{\tikz[remember picture,overlay]\coordinate (#3) at #4;} 
}

\newcommand{\marka}[2]{
{\tikz[remember picture,overlay]\coordinate (#1) at #2;}
}




\newcommand{\cof}{ \mbox{cof}}

\newcommand{\paralela}{\rotatebox[origin=c]{-10}{\Large$\parallel$}}

\usepackage{fancyhdr}
\fancypagestyle{miestilo}{
\fancyhead[L]{}
\fancyhead[C]{}
\fancyhead[R]{}
\fancyfoot[L]{}
\fancyfoot[C]{}
\fancyfoot[R]{pag \thepage}
\renewcommand{\headrulewidth}{0pt}
\renewcommand{\footrulewidth}{0pt}
}
%%%%%%%%%%%%
%%%%%%%%%%%%
\renewcommand{\footrulewidth}{1pt}
\pagestyle{fancy}\markboth{}{}
\renewcommand{\headrulewidth}{.05pt}
\renewcommand{\headrule}{{\color{gray}%
\hrule width\headwidth height\headrulewidth \vskip-\headrulewidth}}
\renewcommand{\footrule}{{\color{gray}%
\vskip-\footruleskip\vskip-\footrulewidth
\hrule width\headwidth height\footrulewidth\vskip\footruleskip}}
\renewcommand{\footrulewidth}{.50pt}
\fancyhead[L]{\color{gray}{\footnotesize \sc ALGEBRA Y GEOMETRIA I} }
\fancyhead[C]{}
\fancyhead[R]{\color{gray}{\sc \footnotesize  Universidad Nacional Del Comahue}}
\fancyfoot[R]{pag. \thepage}
\fancyfoot[C]{\color{gray}{\sc \footnotesize Ingenieria en petroleo}}
\setlength{\headheight}{15pt}

%\usepackage{wrapfig}

\newcommand{\co}{\mathbb{C}}
%\newcommand{\re}{\mathbb{R}}
\newcommand{\q}{\mathbb{Q}}
\newcommand{\z}{\mathbb{Z}}
\newcommand{\n}{\mathbb{N}}
\newcommand{\ve}{\mathbb{V}}
\newcommand{\w}{\mathbb{W}}
\newcommand{\be}{{\cal B}}
\newcommand{\ca}{{\cal C}}
\newcommand{\pqx}{\mathbb{Q}[x]}
\newcommand{\prx}{\mathbb{R}[x]}
\newcommand{\pcx}{\mathbb{C}[x]}
%\newcommand{\izg}{\big\langle}
%\newcommand{\rangle}{\big\rangle}
%%%%%%%%%%%%%%%%%%%%%%%%%%%%%5

%%%%%%%%%%%%%%%%%%%%%%%%%%%%%%%%%%%%5
\newcommand{\materias}{\sc ALGEBRA Y GEOMETRIA I}%separadas con '&'
%%%%%%%%%%%%%%%%%%%%%%%%%%%%%%%%%%%%%%%
% se define el encabezado, 
% el parametro es el texto que va despues de: trabajo practico $N^\circ$... (agregar numero y titulo)
%
%%%%%%%%%%%%%%%%%%%%%%%%%%%%%%%%%\newcommand{\encabezado}[1]{%
\fancyfoot[L]{\vspace*{-3.5mm}\color{gray}{\small T. P. 5: Rectas en el plano.}}
\thispagestyle{miestilo} 
\newcommand{\encabezado}[1]{%
\fancyfoot[L]{\vspace*{-3.5mm}\color{gray}{\small T. P. #1}}
%\thispagestyle{miestilo}%
%--------------------logodpto //texto // logouni------------------
%--------------------logodpto------------------

\vspace*{-18mm}\hspace*{-11mm} \begin{minipage}{2.7cm}
\begin{center} 
     \includegraphics[width=2.7cm, height=2.7cm]{logodpto}
\end{center}  
\end{minipage}\hspace*{-15mm} \begin{minipage}{15.9cm}
\begin{center}
\begin{tabular}{c @{ - } c}
 \materias \\
\carreras
\end{tabular}

\cuatri

\medskip

TRABAJO PRÁCTICO N$^\circ$#1 

\end{center}

\end{minipage}\hspace{-14mm} \begin{minipage}{2.5cm}
  \begin{center} 
    \includegraphics[width=2.5cm, height=2.5cm]{logouni}
  \end{center}  
\end{minipage}
%-----------------------------------------------------

%----------------------------------------------------------------
\hrule
%---------------------------------------------------------------
}%%



\pagestyle{fancy}
\def\identacion{-2.5mm}%mueve la identacion de la enumeracion de los ejercicios.


\usepackage{asypictureB}
\begin{asyheader}
settings.outformat="png";
settings.render=8;
import graph3;
import smoothcontour3;

pen bpen = rgb(0.75, 0.07, 0.01);
//material m = material(diffusepen=0.7bpen, ambientpen=bpen+opacity(0.8), emissivepen=0.3*bpen,specularpen=0.999white, shininess=1.0);      
\end{asyheader}


\usepackage{multicol}
\setlength{\columnsep}{7mm}





\newcommand{\Or}{%
    \tikz[scale=.17%baseline=-0.75ex, 
    ]{ % Cambia "scale" para ajustar el tamaño
        \draw[thick] plot[smooth cycle, tension=1.2] coordinates {(-0.4, 0) (0, -0.6) (0.4, 0) (0, 0.6)};
        \draw[thick,domain=pi:2*pi,samples=200] 
        plot ({0.1*\x * cos(\x r)}, {0.07*\x * sin(\x r)+0.45});
        %\draw[thick,domain=1.5*pi:2*pi,samples=200] plot ({0.08*\x * cos(\x r)-0.9}, {0.09*\x * sin(\x r)-0.1});
    }%
    \,\, 
}

\newcommand{\Noo}{%Es una cruz de que esa solución se descarta
 \tikz[overlay,remember picture] { 
  %%%% cruz en (-, 0)
\draw[red,thick,yshift=2mm] (-.15,\h-.25) -- (+.15,\h+.25); % Línea diagonal ascendente
\draw[red,thick,yshift=2mm] (-.15,\h+.25) -- (+.15,\h-.25); % Línea diagonal descendente
}
}

\begin{document}


\setlength{\headheight}{15pt}

\vspace*{-2.5cm}
%--------------------logodpto------------------
\begin{figure}[htbp]
  \begin{minipage}[r][][t]{0.15\textwidth}
\begin{center} 
 % Universidad Nacional del Comahue.  
 \includegraphics[width=2.2cm, height=2.2cm]{imagen/inglogo.jpg}
\vspace{1.8cm}
\end{center}  
\end{minipage}\quad
% texto central
\begin{minipage}[r][][t]{0.31\textwidth}
\begin{center}
\vspace*{-19mm}
{\bf Universidad Nacional
del Comahue} \\
Facultad de Ingeniería
\end{center}
\end{minipage}\quad\begin{minipage}[r][][t]{0.31\textwidth}
\begin{center}
\vspace*{-20mm}
{\bf  Álgebra y Geometría I}
Ingeniería en Petróleo
Primer Cuatrimestre 2024
\end{center}
\end{minipage}\quad
%----------------------logouni-------------------------
\begin{minipage}[r][][t]{0.15\textwidth}
  \begin{center} 
     \includegraphics[width=2.2cm, height=2.2cm]{imagen/Logo_UNco.jpg}
     \vspace{1.8cm}
  \end{center}  
\end{minipage}
%-----------------------------------------------------
\end{figure}
\vspace*{-22mm}

\begin{center}
{\bf \Large{Resolución Trabajo Práctico 5: Rectas en el plano.}}
\end{center}

\bigskip



\begin{enumerate}[\hspace{-5mm} 1) ]

\item Hallar las ecuaciones, en todas sus formas, de las siguientes rectas definidas gráficamente:

\medskip

\begin{enumerate}[a) ]
   
\begin{minipage}{.45\textwidth}
\item  \, \vspace{-6mm}   
\begin{center}
    



\definecolor{rojo}{rgb}{0.7,0.4,0.4}
\definecolor{azul}{rgb}{.4,.4,0.9}
 
\definecolor{gris}{rgb}{0.5,0.5,0.5}
\begin{tikzpicture}[line cap=round,line join=round,>=triangle 45,scale=.75]

% extremos para los ejes
\xdef\ext{4}
\xdef\yext{3}
\xdef\f{-4}
\xdef\g{-2}
\pgfmathtruncatemacro{\c}{\f+1}
\pgfmathtruncatemacro{\d}{\g+1}
\draw [color=gris!30!,dash pattern=on 1pt off 1pt, xstep=1cm,ystep=1cm] (\f-.3,\g-.3) grid (\ext-.3,\yext+.3);

\draw[->] (\f-.3,0) -- (\ext+.3,0) ;

\draw[->] (0,\g-.3) -- (0,\yext+.3) ;

\foreach \x in {\f,...,-1,1,2,...,\ext}
\draw[shift={(\x,0)},color=black] (0pt,2pt) -- (0pt,-2pt) node[below] {\tiny $\x$};

\foreach \y in {\g,...,-1,1,2,...,\yext }
\draw[shift={(0,\y)},color=black] (2pt,0pt) -- (-2pt,0pt) node[left] {\tiny $\y$};


% 

\coordinate (A) at (3,1);%punto a la derecha
\coordinate (B) at (-2,-1);%punto a la izuierda
\coordinate (M) at ($1.12*(A) -.12*(B)$);%Punto por delente de A
\coordinate (m) at ($-.5*(A)+1.5*(B)$);%Punto por detras de B


%%%%%Grafico de los puntos
%\draw[fill=black] (P) circle (1.5pt) node[shift=(-\ang: 5pt),right] {\scriptsize $P=(-2,1)$};


%%%% Recta que pasa por A y B
\draw[blue,fill=blue] (A) circle (2pt) ;

\draw[blue,fill=blue] (B) circle (2pt);
 
% Dibujo de la línea entre A y B
\draw[thick] (m) -- (M);


\end{tikzpicture}
  
\end{center}  




\end{minipage}\quad\begin{minipage}{.45\textwidth}


\item \, \vspace{-6mm}   
\begin{center}
    



\definecolor{rojo}{rgb}{0.7,0.4,0.4}
\definecolor{azul}{rgb}{.4,.4,0.9}
 
\definecolor{gris}{rgb}{0.5,0.5,0.5}
\begin{tikzpicture}[line cap=round,line join=round,>=triangle 45,scale=.75]

% extremos para los ejes
\xdef\ext{4}
\xdef\yext{1}
\xdef\f{-3}
\xdef\g{-4} 
\pgfmathtruncatemacro{\c}{\f+1}
\pgfmathtruncatemacro{\d}{\g+1}
\draw [color=gris!30!,dash pattern=on 1pt off 1pt, xstep=1cm,ystep=1cm] (\f-.3,\g-.3) grid (\ext+.3,\yext+.3);
\draw[->] (\f-.3,0) -- (\ext+.3,0) ;

\draw[->] (0,\g-.3) -- (0,\yext+.3) ;


\foreach \x in {\f,...,-1,1,2,...,\ext}
\draw[shift={(\x,0)},color=black] (0pt,2pt) -- (0pt,-2pt) node[below] {\tiny $\x$};

\foreach \y in {\g,\d,...,-1,1,...,\yext }
\draw[shift={(0,\y)},color=black] (2pt,0pt) -- (-2pt,0pt) node[left] {\tiny $\y$};


% 

\coordinate (A) at (4,-4);%punto a la derecha
\coordinate (B) at (-1,-1);%punto a la izuierda
\coordinate (M) at ($1.1*(A) -.1*(B)$);%Punto por delente de A
\coordinate (m) at ($-.49*(A)+1.49*(B)$);%Punto por detras de B


%%%%%Grafico de los puntos
%\draw[fill=black] (P) circle (1.5pt) node[shift=(-\ang: 5pt),right] {\scriptsize $P=(-2,1)$};


%%%% Recta que pasa por A y B
\draw[blue,fill=blue] (A) circle (2pt) ;

\draw[blue,fill=blue] (B) circle (2pt);

% Dibujo de la línea entre A y B
\draw[thick] (m) -- (M);


\end{tikzpicture}
  
\end{center}  




\end{minipage}


\end{enumerate}









\textbf{Solución:}

\begin{enumerate}
    \item \,


\hspace{-13mm}
\begin{tikzpicture}[xscale=8,yscale=5.2,>=triangle 45,rounded corners=3mm]
%\clip(0,0) rectangle (1,1);
%\begin{scriptsize}
\path (0,.3) node[draw=orange,below right](a)  {
\definecolor{rojo}{rgb}{0.7,0.4,0.4}
\definecolor{azul}{rgb}{.4,.4,0.9}
 
\definecolor{gris}{rgb}{0.5,0.5,0.5}
\begin{tikzpicture}[line cap=round,line join=round,>=triangle 45,scale=.5]

% extremos para los ejes
\xdef\ext{4}
\xdef\yext{3}
\xdef\f{-4}
\xdef\g{-2}
\pgfmathtruncatemacro{\c}{\f+1}
\pgfmathtruncatemacro{\d}{\g+1}
\draw [color=gris!30!,dash pattern=on 1pt off 1pt, xstep=1cm,ystep=1cm] (\f-.3,\g-.3) grid (\ext-.3,\yext+.3);

\draw[->] (\f-.3,0) -- (\ext+.3,0) ;

\draw[->] (0,\g-.3) -- (0,\yext+.3) ;

\foreach \x in {\f,...,-1,1,2,...,\ext}
\draw[shift={(\x,0)},color=black] (0pt,2pt) -- (0pt,-2pt) node[below] {\tiny $\x$};

\foreach \y in {\g,...,-1,1,2,...,\yext }
\draw[shift={(0,\y)},color=black] (2pt,0pt) -- (-2pt,0pt) node[left] {\tiny $\y$};


% 

\coordinate (A) at (-2,-1);%punto a la izquierda
\coordinate (B) at (3,1);%punto a la derecha
\coordinate (M) at ($1.5*(A) -.5*(B)$);%Punto por detras de A
\coordinate (m) at ($-.2*(A)+1.2*(B)$);%Punto por delante de B

% vector director
    \draw[->, thick, purple,line width=2pt] (A) -- (B) node[midway, above left,pos=.75] {$\vec{d}$};
    
%%%%%Grafico de los puntos
%\draw[fill=black] (P) circle (1.5pt) node[shift=(-\ang: 5pt),right] {\scriptsize $P=(-2,1)$};


%%%% Recta que pasa por A y B
\draw[blue,fill=blue] (A) circle (2pt) node[below]{$A$} ;

\draw[blue,fill=blue] (B) circle (2pt)node[below]{$B$} ;
 
% Dibujo de la línea entre A y B
\draw[thick] (m) -- (M);



\end{tikzpicture}
  
} 
(.8,.3) node[draw=orange,text width=70mm,below right] (bb) {
Se tiene que la recta pasa por los puntos $A=(-2,-1)$ y $B= (3,1)$. 
}
(1.3,-.3) node[draw=orange,text width=68mm,below right] (c) { 
  Tenemos que la recta  pasa por \linebreak  $A=(-2\indicador{a}{(-.1,-.15)}{a1}{(-.2,-.4)}
  ,-1\indicador{b}{(-.1,-.1)}{b1}{(.1,-.4)}
  )$ y tiene director $\vec{d} = (5\indicador{u}{(-.1,-.1)}{u1}{(-.2,-.4)} 
  ,2 \indicador{uu}{(-.1,-.1)}{uu1}{(.1,-.4)} 
  )$.
  \medskip
\medskip
}
(1.4,-1.5) node[draw=red,line width=2pt,text width=70mm,below right] (EcuVectorial) {
\quad\textbf{Ecuación vectorial:} 
\[
L: \ \ (x,y) =  (-2,-1) + \lambda (5,2 ) , \quad \lambda \in \R 
\]
}
(0,-1.3) node[draw=red,line width=2pt,text width=56mm,below right] (EcuParaM) {
\quad\textbf{Ecuaciones paramétricas}
\[
L : \left\{
\begin{array}{rl}
    x &= -2+5\lambda
    \\[1mm]
    y &=-1 + 2\lambda
\end{array}
\right. , \quad \lambda\in \R.
\]
}
(.8,-1.5) node[draw=red,line width=2pt,text width=38mm,below right] (EcuCart) {
\textbf{Ecuación cartesiana}\\[3mm]

$\dfrac{x +2}{5} 
=%%%%%
\dfrac{y+1}{2}
$  
}
(1.7,-2) node[draw=red,line width=2pt,text width=39mm,below right] (EcExpli) {
    \textbf{Ecuación explícita}
\[
 y = \frac{2}{5}x - \frac{1}{5}
\]
}
(.5,-2) node[draw=cyan,text width=70mm,below right] (d) { 
De la ecuación cartesiana  $y$:
\begin{align*}
 \dfrac{x +2}{5} 
&=%%%%%
\dfrac{y+1}{2}
\\[2mm]
\dfrac{2}{5} (x +2)
&=%%%%%
y+1
\\[2mm]
\tfrac{2}{5}x +\tfrac{4}{5} -1
&=%%%%%
y
\\[2mm]
 \tfrac{2}{5}x -\tfrac{1}{5} 
&=y   
\end{align*}
}
;%--------------%------------------%----------------------------%------------------------------------%-----------------------%-----------------------------------------------------%--------------------------------%-----------------------------%----------------------------%----------------------------%-------------------------------------------------------
\draw[thick,->,dashed,color=teal] (bb) to[out=-90,in=165] node [pos=0.65,text width=43mm ,  left,xshift=4mm ,yshift=-4mm
] {\footnotesize
Podemos tomar como  vector director a  $\overrightarrow{AB}$:
\begin{align*}
  \overrightarrow{AB} &=  (3,1)- (-2,-1)   
  \\
  &= \left(3-(-2) ,\, 1-(-1) \right)
  \\
  &= \left(5,2 \right)
\end{align*}
}(c);
%----------------------------------------------------------------------------------------------------------------------------------------------------------------------------
\draw[thick,->,dashed,color=blue] (c) to  node [pos=0.6,text width=40mm ,  right,xshift=-14mm,fill=white 
] { \footnotesize 
\qquad formula:
$
\textcolor{purple}{
L: (x, y) = (x_0, y_0) + \lambda (u_x, u_y), \quad \lambda \in \mathbb{R}
}
$
    } (EcuVectorial) ;
\draw[thick,->,dashed,color=blue] (c) to[out=185,in=60]  node [pos=0.3,text width=40mm ,  below ,xshift=-2mm, fill=white %,yshift=-8mm
] { \footnotesize 
\quad formula: \\
$
\textcolor{purple}{
L: 
\begin{cases}
x = x_0 + u_x \, \lambda \\
y = y_0 + u_y \, \lambda
\end{cases},
\lambda \in \mathbb{R}
}$
%%%%%%%%%%%%%%%%%%%%%%%%%%%%%%%%%%%%%%%%%%%%%%%%%%%%%%%%%%%%%%%%%%%%%%%%%%%%%%%%%%%%%%%%%%%%%%%%%%%%%%%%%%%%%%%%%%
    } (EcuParaM) ;
\draw[thick,->,dashed,color=blue] (c) to  node [pos=0.3,text width=25mm ,  below ,xshift=-2mm,fill=white 
] { \footnotesize 
\quad formula: \\
\textcolor{purple}{$\dfrac{x -x_0}{u_x} 
=%%%%%
\dfrac{y-y_0}{u_y}
$ }
    } (EcuCart) ;
\draw[thick,->,dashed,color=blue] (EcuCart) to (d);
\draw[thick,->,dashed,color=blue] (d) to (EcExpli);
\end{tikzpicture}%%%%%%%%%%%%%%%%%%%%%%%%%%%%%%%%%%%%%%%%%%%%%%%%%%%%%%%%%%%%%%%%%%%%%%%%%%%%%%%%%%%%%%%%%%%%%%%%%%%%%%%%%%%%%%%%%%%%%%%%%%%%%%%%%%%%%%%%%%%%%%%%%%%%%%%%%%%%%%%%%%%%%%%%%%%%%%%%%%%%%%%%%%%%%%%%%%%%%%%%%%%%%%%%%%%%%%%%%%%%%%%%%%%%%%%%%%%%%%%%%%%%%%%%%%%%%%%%%%%%%%%%%%%%%%%%%%%%%%%%%%%%%%%%%%%%%%%%%%%%%%%%%%%%%%%%%%%%%%%%%%%%%%%%%%%%%%%%%%%%%%%%%%%%%%%%%%%%%%%%%%%%%%%%%%%%%%%%%%%%%%%%%%%%%%%%%%%%%%%%%%%%%%%%%%%%%%%%%%%%%%%%%%%%%%%%%%%%%%%%%%%%%%%%%%%%%%%%
\begin{tikzpicture}[overlay,remember picture]
\node at (a1) [ blue,below  ]{ 
$x_0$
};
\draw [overlay,<-,very thick,blue,opacity=.5]
(a) |-(a1);  %%%%%%%%%%%%%%%%%%%%%%%%%%%%%%%%%%%%%%%%%%%%%%%%%%%%%%%%%%%%%%%%%%%%%%%%%%%%%%%%%%%%%%%%%%%%%%%%%%%%%%%%%%%%%%%%%%%%%%%%%%%%%%%%%%%%%%%%%%%%%%%%
\node at (b1) [ blue,below]{ 
$y_0$
};
\draw [overlay,<-,very thick,blue,opacity=.5]
(b) |-(b1); 
%%%%%%%%%%%%%%%%%%%%%%%%%%%%%%%%%%%%%%%%%%%%%%%%%%%%%%%%%%%%%%%%%%%%%%%%%%%%%%%%%%%%%%%%%%%%%%%%%%%%%%%%%%%%%%%%%%%%%%%%%%%%%%%%%%%%%%%%%%%%%%%%
\draw [overlay,<-,very thick,blue,opacity=.5]
(u) |-(u1); 
\node at (u1) [ blue,below,xshift=2mm ]{ 
$u_x$
}; %%%%%%%%%%%%%%%%%%%%%%%%%%%%%%%%%%%%%%%%%%%%%%%%%%%%%%%%%%%%%%%%%%%%%%%%%%%%%%%%%%%%%%%%%%%%%%%%%%%%%%%%%%%%%%%%%%%%%%%%%%%%%%%%%%%%%%%%%%%%%%%%
\draw [overlay,<-,very thick,blue,opacity=.5]
(uu) |-(uu1); 
\node at (uu1) [ blue,below,xshift=2mm ]{ 
$u_y$
}; 
\end{tikzpicture} 















\begin{tikzpicture}[xscale=8,yscale=4.3,>=triangle 45,rounded corners=3mm]
%\clip(0,0) rectangle (1,1);
%\begin{scriptsize}
\path (0,0) node[draw=orange,below right](a){



\definecolor{rojo}{rgb}{0.7,0.4,0.4}
\definecolor{azul}{rgb}{.4,.4,0.9}
 
\definecolor{gris}{rgb}{0.5,0.5,0.5}
\begin{tikzpicture}[line cap=round,line join=round,>=triangle 45, scale=.5]

% extremos para los ejes
\xdef\ext{4}
\xdef\yext{3}
\xdef\f{-4}
\xdef\g{-2}
\pgfmathtruncatemacro{\c}{\f+1}
\pgfmathtruncatemacro{\d}{\g+1}
\draw [color=gris!30!,dash pattern=on 1pt off 1pt, xstep=1cm,ystep=1cm] (\f-.3,\g-.3) grid (\ext-.3,\yext+.3);

\draw[->] (\f-.3,0) -- (\ext+.3,0) ;

\draw[->] (0,\g-.3) -- (0,\yext+.3) ;

\foreach \x in {\f,...,-1,1,2,...,\ext}
\draw[shift={(\x,0)},color=black] (0pt,2pt) -- (0pt,-2pt) node[below] {\tiny $\x$};

\foreach \y in {\g,...,-1,1,2,...,\yext }
\draw[shift={(0,\y)},color=black] (2pt,0pt) -- (-2pt,0pt) node[left] {\tiny $\y$};


% 

\coordinate (A) at (-2,-1);%punto a la izquierda
\coordinate (B) at (3,1);%punto a la derecha
\coordinate (M) at ($1.5*(A) -.5*(B)$);%Punto por detras de A
\coordinate (m) at ($-.2*(A)+1.2*(B)$);%Punto por delante de B
% vector director
    \draw[->, thick, purple,line width=2pt] (A) -- (B) node[midway, above left,pos=.75] {$\vec{d}$};
    
     % Calcular el punto rotado 90° alrededor de A 
    

     

%%%%%Grafico de los puntos
%\draw[fill=black] (P) circle (1.5pt) node[shift=(-\ang: 5pt),right] {\scriptsize $P=(-2,1)$};


%%%% Recta que pasa por A y B
\draw[blue,fill=blue] (A) circle (2pt) node[below]{$A$} ;

\draw[blue,fill=blue] (B) circle (2pt)node[below]{$B$} ;
 
% Dibujo de la línea entre A y B
\draw[thick] (m) -- (M);

    \coordinate (C) at ($(B)-(A)$);
    \coordinate (Cperp) at ([rotate around={90:(0,0)}]C);
    \coordinate (CC) at ($(Cperp)+(A)$);
    % Dibujar el vector perpendicular

\draw[rotate=90,->, thick, red] (A) -- (CC) node[midway, left] {\small $\vec{n}$};

    \xdef\r{.1}
    \coordinate (E) at ($\r*(Cperp)+(A)$);
    \coordinate (F) at ($\r*(Cperp)+\r*(C)+(A)$);
    \coordinate (G) at ($\r*(C)+(A)$);      
  \draw[red,dashed] (E) --(F)--(G) ;
  \fill[red,opacity=.3,,rounded corners=0mm] (A)--(E) --(F)--(G)--cycle ;
\end{tikzpicture}
  
}
(.65,0) node[draw=orange,text width=110mm,below right](bb){

Por un lado, el vector director de la recta es $\overrightarrow{AB} = (5,2)$, lo que implica que un vector normal a la recta es $\vec{n} = (-2,5)$.  

Luego la recta pasa por el punto 
$ A = (-2\indicador{a}{(-.1,-.15)}{a1}{(-.2,-.4)} 
,-1 \indicador{b}{(-.1,-.1)}{b1}{(.1,-.4)} 
)$ y es perpendicular al vector normal $\vec{n} = (-2\indicador{n}{(-.1,-.1)}{n1}{(-.2,-.4)} 
,\, 5 
\indicador{nn}{(-.1,-.1)}{nn1}{(.1,-.4)} 
)$.

\begin{tikzpicture}[overlay,remember picture, >=angle 90]
\node at (a1) [ blue,below  ]{ 
$x_0$
};
\draw [overlay,<-,very thick,blue,opacity=.5,rounded corners=0mm]
(a) |-(a1);  %%%%%%%%%%%%%%%%%%%%%%%%%%%%%%%%%%%%%%%%%%%%%%%%%%%%%%%%%%%%%%%%%%%%%%%%%%%%%%%%%%%%%%%%%%%%%%%%%%%%%%%%%%%%%%%%%%%%%%%%%%%%%%%%%%%%%%%%%%%%%%%%
\node at (b1) [ blue,below]{ 
$y_0$
};
\draw [overlay,<-,very thick,blue,opacity=.5,rounded corners=0mm]
(b) |-(b1); 
%%%%%%%%%%%%%%%%%%%%%%%%%%%%%%%%%%%%%%%%%%%%%%%%%%%%%%%%%%%%%%%%%%%%%%%%%%%%%%%%%%%%%%%%%%%%%%%%%%%%%%%%%%%%%%%%%%%%%%%%%%%%%%%%%%%%%%%%%%%%%%%%
\draw [overlay,<-,very thick,blue,opacity=.5,rounded corners=0mm]
(n) |-(n1); 
\node at (n1) [ blue,below,xshift=2mm ]{ 
$n_1$
}; %%%%%%%%%%%%%%%%%%%%%%%%%%%%%%%%%%%%%%%%%%%%%%%%%%%%%%%%%%%%%%%%%%%%%%%%%%%%%%%%%%%%%%%%%%%%%%%%%%%%%%%%%%%%%%%%%%%%%%%%%%%%%%%%%%%%%%%%%%%%%%%%
\draw [overlay,<-,very thick,blue,opacity=.5,rounded corners=0mm]
(nn) |-(nn1); 
\node at (nn1) [ blue,below,xshift=2mm ]{ 
$n_2$
}; 
\end{tikzpicture} 

\medskip
\medskip

%con esos datos podemos calcular su ecuación usando la formula:
%\[ (x - x_0, y - y_0) \cdot (n_1, n_2) = 0\]

}
(0.7,-.7)  node[draw=orange,text width=63mm,below right](c){
\small
\vspace{-4mm}\begin{align*}
    \left(x -(-2), y - (-1) \rule{0mm}{4mm}\right) \cdot (-2, 5 ) &= 0
    \\[2mm]
    (x+2, y +1) \cdot (-2, 5 ) &= 0
    \\[2mm]
   -2 (x+2)  +  5(y +1)&= 0
   \\[2mm]
   -2 x-4   +  5y + 5  &= 0
   \\[2mm]
   -2 x   +  5y + 1&= 0
\end{align*}

}
(0,-1.3)  node[draw=red,text width=37mm,line width=2pt,below right](EcuImpli){
\textbf{Ecuación implicita} 
\[
 -2x + 5y + 1 = 0
\] 

}
(.46,-1.65)  node[draw=orange,text width=75mm ,below right](d){
\footnotesize A partir de la ecuación implícita, realizamos manipulaciones algebraicas para obtener la ecuación segmentaria.
\begin{align*}
 -2x + 5y + 1 &= 0
 \\
  -2x + 5y  &= -1
  \\
  2x - 5y  &= 1
  \\
 \dfrac{\, x\,}{\tfrac{1}{2}}  + \dfrac{y}{-\tfrac{1}{5}}   &= 1 
\end{align*}


}
(1.5,-2)  node[draw=red,text width=43mm,line width=2pt,below right](EcuSeg){
\textbf{Ecuación segmentaria}  
\[
\dfrac{\ x\ }{\tfrac{1}{2}}  + \dfrac{\ y \ }{-\tfrac{1}{5}\  \,}   = 1
\]

};%--------------%------------------%----------------------------%------------------------------------%-----------------------%-----------------------------------------------------%--------------------------------%-----------------------------%----------------------------%----------------------------%-------------------------------------------------------
\draw[thick,->,dashed,color=teal] (bb) to[out=-45,in=20]  node [pos=0.5,text width=43mm ,below right,xshift=-4mm
] {\footnotesize
\qquad Usaremos la formula:
\begin{align*}
     (x - x_0, y - y_0) \cdot (n_1, n_2) &= 0
\end{align*}
}(c);
\draw[thick,->,dashed,color=teal] (c) to[out=180, in=60] (EcuImpli);
\draw[thick,->,dashed,color=teal] (EcuImpli) to[out=-125, in=180] (d);
\draw[thick,->,dashed,color=teal] (d) to (EcuSeg);
\end{tikzpicture}

 



\begin{comment}
    
{\bf Ecuación implícita o general de la recta en $\mathbb{R}^2$}

Sean un punto $P_0=(x_0, y_0)$ del plano y un vector $\vec{n} = (n_1, n_2)$ distinto del vector nulo, la ecuación implicita es: 

 \[
(x - x_0, y - y_0) \cdot (n_1, n_2) = 0
\]

Al resolver el producto escalar entre estos vectores, resulta

\[
n_1 (x - x_0) + n_2 (y - y_0) = 0
\]

Luego, operando en $\mathbb{R}$,

\[
n_1 x + n_2 y + (-n_1 x_0 - n_2 y_0) = 0
\]

Se tiene que la ecuación de la recta es:

\[
r : Ax + By + C = 0
\]

Esta expresión de la recta en $\mathbb{R}^2$ se denomina ecuación implícita, o general, de la recta, donde los coeficientes $A$ y $B$ no son simultáneamente nulos.

\end{comment}








    \item \,


\hspace{-13mm}
\begin{tikzpicture}[xscale=8,yscale=4.4,>=triangle 45,rounded corners=3mm]
%\clip(0,0) rectangle (1,1);
%\begin{scriptsize}
\path (0,0) node[draw=orange,below right](a)  {

\definecolor{rojo}{rgb}{0.7,0.4,0.4}
\definecolor{azul}{rgb}{.4,.4,0.9}
 
\definecolor{gris}{rgb}{0.5,0.5,0.5}
\begin{tikzpicture}[line cap=round,line join=round,>=triangle 45,scale=.5]

% extremos para los ejes
\xdef\ext{4}
\xdef\yext{1}
\xdef\f{-3}
\xdef\g{-4} 
\pgfmathtruncatemacro{\c}{\f+1}
\pgfmathtruncatemacro{\d}{\g+1}
\draw [color=gris!30!,dash pattern=on 1pt off 1pt, xstep=1cm,ystep=1cm] (\f-.3,\g-.3) grid (\ext+.3,\yext+.3);
\draw[->] (\f-.3,0) -- (\ext+.3,0) ;

\draw[->] (0,\g-.3) -- (0,\yext+.3) ;


\foreach \x in {\f,...,-1,1,2,...,\ext}
\draw[shift={(\x,0)},color=black] (0pt,2pt) -- (0pt,-2pt) node[below] {\tiny $\x$};

\foreach \y in {\g,\d,...,-1,1,...,\yext }
\draw[shift={(0,\y)},color=black] (2pt,0pt) -- (-2pt,0pt) node[left] {\tiny $\y$};


% 

\coordinate (A) at (-1,-1);%punto a la derecha
\coordinate (B) at (4,-4);%punto a la izuierda
\coordinate (M) at ($1.49*(A) -.49*(B)$);%Punto por delente de A
\coordinate (m) at ($-.1*(A)+1.1*(B)$);%Punto por detras de B


% vector director
    \draw[->, thick, purple,line width=2pt] (A) -- (B) node[midway, above left,pos=.75] {$\vec{d}$};
    
%%%%%Grafico de los puntos
%\draw[fill=black] (P) circle (1.5pt) node[shift=(-\ang: 5pt),right] {\scriptsize $P=(-2,1)$};


%%%% Recta que pasa por A y B
\draw[blue,fill=blue] (A) circle (2pt) node[below]{$A$} ;

\draw[blue,fill=blue] (B) circle (2pt)node[below]{$B$} ;

% Dibujo de la línea entre A y B
\draw[thick] (m) -- (M);


\end{tikzpicture}
  
} 
(.65,0) node[draw=orange,text width=70mm,below right] (bb) {
Se tiene que la recta pasa por los puntos $A=(-1,-1)$ y $B= (4,-4)$. 
}
(1.4,-.5) node[draw=orange,text width=71mm,below right] (c) { 
  Tenemos que la recta  pasa por \linebreak  $A=(-1\indicador{a}{(-.1,-.15)}{a1}{(-.2,-.4)}
  ,-1\indicador{b}{(-.1,-.1)}{b1}{(.1,-.4)}
  )$ y tiene director $\vec{d} = (5\indicador{u}{(-.1,-.1)}{u1}{(-.2,-.4)} 
  , -3 \indicador{uu}{(-.1,-.1)}{uu1}{(.1,-.4)} 
  )$.
  \medskip
\medskip
}
(1.4,-1.5) node[draw=red,line width=2pt,text width=70mm,below right] (EcuVectorial) {
\quad\textbf{Ecuación vectorial:} 
\[
L: \ \ (x,y) =  (-1,-1) + \lambda (5,-3) , \quad \lambda \in \R 
\]
}
(0,-1.3) node[draw=red,line width=2pt,text width=56mm,below right] (EcuParaM) {
\quad\textbf{Ecuaciones paramétricas}
\[
L : \left\{
\begin{array}{rl}
    x &= -1+5\lambda
    \\[1mm]
    y &=-1 - 3\lambda
\end{array}
\right. , \quad \lambda\in \R.
\]
}
(.8,-1.5) node[draw=red,line width=2pt,text width=38mm,below right] (EcuCart) {
\textbf{Ecuación cartesiana}\\[3mm]

$\dfrac{x +1}{5} 
=%%%%%
\dfrac{y+1}{-3}
$  
}
(1.7,-2.4) node[draw=red,line width=2pt,text width=39mm,below right] (EcExpli) {
    \textbf{Ecuación explícita}
\[
 y =- \tfrac{3}{5}x -\tfrac{8}{5} 
\]
}
(.5,-2.1) node[draw=cyan,text width=70mm,below right] (d) { 
\small
De la ecuación cartesiana despejamos $y$:
\begin{align*}
 \dfrac{x+1}{5} 
&=%%%%%
\dfrac{y+1}{-3}
\\ 
\dfrac{-3}{5} (x +1)
&=%%%%%
y+1
\\[2mm]
\tfrac{-3}{5}x +\tfrac{-3}{5} -1
&=%%%%%
y
\\[2mm]
- \tfrac{3}{5}x -\tfrac{8}{5} 
&=y   
\end{align*}
}
;%--------------%------------------%----------------------------%------------------------------------%-----------------------%-----------------------------------------------------%--------------------------------%-----------------------------%----------------------------%----------------------------%-------------------------------------------------------
\draw[thick,->,dashed,color=teal] (bb) to[out=-90,in=165] node [pos=0.85,text width=43mm ,  left,xshift=2mm ,yshift=-6mm
] {\footnotesize
$\vec{n}=\overrightarrow{AB}$:
\begin{align*}
  \overrightarrow{AB} &=  (4,-4)- (-1,-1)
  \\
  &= \left(4-(-1) ,\, -4-(-1) \right)
  \\
  &= \left(5,-3 \right)
\end{align*}
}(c);
%----------------------------------------------------------------------------------------------------------------------------------------------------------------------------
\draw[thick,->,dashed,color=blue] (c) to  node [pos=0.4,text width=40mm ,  right, xshift=-14mm , fill=white
] { \footnotesize 
\qquad\qquad formula:\\ 
$
\textcolor{purple}{
L: (x, y) = (x_0, y_0) + \lambda (u_x, u_y), \quad \lambda \in \mathbb{R}
}
$
} (EcuVectorial) ;
\draw[thick,->,dashed,color=blue] (c) to[out=185,in=60]  node [pos=0.3,text width=40mm , below ,xshift=-4mm, fill=white %,yshift=-8mm
] { \footnotesize 
\qquad \qquad formula:\\
$
\textcolor{purple}{
L: 
\begin{cases}
x = x_0 + u_x \, \lambda \\
y = y_0 + u_y \, \lambda
\end{cases},
\lambda \in \mathbb{R}
}$
    } (EcuParaM);
%%%%%%%%%%%%%%%%%%%%%%%%%%%%%%%%%%%%%%%%%%%%%%%%%%%%%%%%%%%%%%%%%%%%%%%%%%%%%%%%%%%%%%%%%%%%%%%%%%%%%%%%%%%%%%%%%%
\draw[thick,->,dashed,color=blue] (c) to  node [pos=0.1,text width=25mm ,  below ,xshift=-10mm, fill=white
] { \footnotesize 
\qquad formula: \\
\textcolor{purple}{$\dfrac{x -x_0}{u_x} 
=%%%%%
\dfrac{y-y_0}{u_y}
$ }
    } (EcuCart) ;
\draw[thick,->,dashed,color=blue] (EcuCart) to (d);
\draw[thick,->,dashed,color=blue] (d) to (EcExpli);
\end{tikzpicture}%%%%%%%%%%%%%%%%%%%%%%%%%%%%%%%%%%%%%%%%%%%%%%%%%%%%%%%%%%%%%%%%%%%%%%%%%%%%%%%%%%%%%%%%%%%%%%%%%%%%%%%%%%%%%%%%%%%%%%%%%%%%%%%%%%%%%%%%%%%%%%%%%%%%%%%%%%%%%%%%%%%%%%%%%%%%%%%%%%%%%%%%%%%%%%%%%%%%%%%%%%%%%%%%%%%%%%%%%%%%%%%%%%%%%%%%%%%%%%%%%%%%%%%%%%%%%%%%%%%%%%%%%%%%%%%%%%%%%%%%%%%%%%%%%%%%%%%%%%%%%%%%%%%%%%%%%%%%%%%%%%%%%%%%%%%%%%%%%%%%%%%%%%%%%%%%%%%%%%%%%%%%%%%%%%%%%%%%%%%%%%%%%%%%%%%%%%%%%%%%%%%%%%%%%%%%%%%%%%%%%%%%%%%%%%%%%%%%%%%%%%%%%%%%%%%%%%%%
\begin{tikzpicture}[overlay,remember picture]
\node at (a1) [ blue,below  ]{ 
$x_0$
};
\draw [overlay,<-,very thick,blue,opacity=.5]
(a) |-(a1);  %%%%%%%%%%%%%%%%%%%%%%%%%%%%%%%%%%%%%%%%%%%%%%%%%%%%%%%%%%%%%%%%%%%%%%%%%%%%%%%%%%%%%%%%%%%%%%%%%%%%%%%%%%%%%%%%%%%%%%%%%%%%%%%%%%%%%%%%%%%%%%%%
\node at (b1) [ blue,below]{ 
$y_0$
};
\draw [overlay,<-,very thick,blue,opacity=.5]
(b) |-(b1); 
%%%%%%%%%%%%%%%%%%%%%%%%%%%%%%%%%%%%%%%%%%%%%%%%%%%%%%%%%%%%%%%%%%%%%%%%%%%%%%%%%%%%%%%%%%%%%%%%%%%%%%%%%%%%%%%%%%%%%%%%%%%%%%%%%%%%%%%%%%%%%%%%
\draw [overlay,<-,very thick,blue,opacity=.5]
(u) |-(u1); 
\node at (u1) [ blue,below,xshift=2mm ]{ 
$u_x$
}; %%%%%%%%%%%%%%%%%%%%%%%%%%%%%%%%%%%%%%%%%%%%%%%%%%%%%%%%%%%%%%%%%%%%%%%%%%%%%%%%%%%%%%%%%%%%%%%%%%%%%%%%%%%%%%%%%%%%%%%%%%%%%%%%%%%%%%%%%%%%%%%%
\draw [overlay,<-,very thick,blue,opacity=.5]
(uu) |-(uu1); 
\node at (uu1) [ blue,below,xshift=.5mm ]{ 
$u_y$
}; 
\end{tikzpicture} 















\begin{tikzpicture}[xscale=8,yscale=5.4,>=triangle 45,rounded corners=3mm]
%\clip(0,0) rectangle (1,1);
%\begin{scriptsize}
\path (0,0) node[draw=orange,below right](a){



\definecolor{rojo}{rgb}{0.7,0.4,0.4}
\definecolor{azul}{rgb}{.4,.4,0.9}
 
\definecolor{gris}{rgb}{0.5,0.5,0.5}
\begin{tikzpicture}[line cap=round,line join=round,>=triangle 45,scale=.5]

% extremos para los ejes
\xdef\ext{4}
\xdef\yext{4}
\xdef\f{-3}
\xdef\g{-4} 
\pgfmathtruncatemacro{\c}{\f+1}
\pgfmathtruncatemacro{\d}{\g+1}
\draw [color=gris!30!,dash pattern=on 1pt off 1pt, xstep=1cm,ystep=1cm] (\f-.3,\g-.3) grid (\ext+.3,\yext+.3);
\draw[->] (\f-.3,0) -- (\ext+.3,0) ;

\draw[->] (0,\g-.3) -- (0,\yext+.3) ;


\foreach \x in {\f,...,-1,1,2,...,\ext}
\draw[shift={(\x,0)},color=black] (0pt,2pt) -- (0pt,-2pt) node[below] {\tiny $\x$};

\foreach \y in {\g,\d,...,-1,1,...,\yext }
\draw[shift={(0,\y)},color=black] (2pt,0pt) -- (-2pt,0pt) node[left] {\tiny $\y$};


% 

\coordinate (A) at (-1,-1);%punto a la derecha
\coordinate (B) at (4,-4);%punto a la izuierda
\coordinate (M) at ($1.46*(A) -.46*(B)$);%Punto por delente de A
\coordinate (m) at ($-.1*(A)+1.1*(B)$);%Punto por detras de B

% vector director
    \draw[->, thick, purple,line width=2pt] (A) -- (B) node[midway, above left,pos=.75] {$\vec{d}$};
    
     % Calcular el punto rotado 90° alrededor de A 
    

     

%%%%%Grafico de los puntos
%\draw[fill=black] (P) circle (1.5pt) node[shift=(-\ang: 5pt),right] {\scriptsize $P=(-2,1)$};


%%%% Recta que pasa por A y B
\draw[blue,fill=blue] (A) circle (2pt) node[below]{$A$} ;

\draw[blue,fill=blue] (B) circle (2pt)node[below]{$B$} ;
%%%%%Grafico de los puntos
%\draw[fill=black] (P) circle (1.5pt) node[shift=(-\ang: 5pt),right] {\scriptsize $P=(-2,1)$};


%%%% Recta que pasa por A y B
\draw[blue,fill=blue] (A) circle (2pt) ;

\draw[blue,fill=blue] (B) circle (2pt);

% Dibujo de la línea entre A y B
\draw[thick] (m) -- (M);

\coordinate (C) at ($(B)-(A)$);
    \coordinate (Cperp) at ([rotate around={90:(0,0)}]C);
    \coordinate (CC) at ($(Cperp)+(A)$);
    % Dibujar el vector perpendicular

\draw[rotate=90,->, thick, red] (A) -- (CC) node[midway, left,pos=.7] {\small $\vec{n}$};

    \xdef\r{.1}
    \coordinate (E) at ($\r*(Cperp)+(A)$);
    \coordinate (F) at ($\r*(Cperp)+\r*(C)+(A)$);
    \coordinate (G) at ($\r*(C)+(A)$);      
  \draw[red,dashed] (E) --(F)--(G) ;
  \fill[red,opacity=.3,,rounded corners=0mm] (A)--(E) --(F)--(G)--cycle ;
\end{tikzpicture}

}
(.65,0) node[draw=orange,text width=110mm,below right](bb){

Por un lado, el vector un director de la recta es $\overrightarrow{AB} = (5,-3)$, lo que implica que un vector normal  es $\vec{n} = (3,5)$.  

Luego la recta pasa por el punto 
$ A = (-1\indicador{a}{(-.1,-.15)}{a1}{(-.2,-.4)} 
,-1 \indicador{b}{(-.1,-.1)}{b1}{(.1,-.4)} 
)$ y es perpendicular al vector normal $\vec{n} = (3\indicador{n}{(-.1,-.1)}{n1}{(-.2,-.4)} 
,\, 5 
\indicador{nn}{(-.1,-.1)}{nn1}{(.1,-.4)} 
)$.

\begin{tikzpicture}[overlay,remember picture, >=angle 90]
\node at (a1) [ blue,below  ]{ 
$x_0$
};
\draw [overlay,<-,very thick,blue,opacity=.5,rounded corners=0mm]
(a) |-(a1);  %%%%%%%%%%%%%%%%%%%%%%%%%%%%%%%%%%%%%%%%%%%%%%%%%%%%%%%%%%%%%%%%%%%%%%%%%%%%%%%%%%%%%%%%%%%%%%%%%%%%%%%%%%%%%%%%%%%%%%%%%%%%%%%%%%%%%%%%%%%%%%%%
\node at (b1) [ blue,below]{ 
$y_0$
};
\draw [overlay,<-,very thick,blue,opacity=.5,rounded corners=0mm]
(b) |-(b1); 
%%%%%%%%%%%%%%%%%%%%%%%%%%%%%%%%%%%%%%%%%%%%%%%%%%%%%%%%%%%%%%%%%%%%%%%%%%%%%%%%%%%%%%%%%%%%%%%%%%%%%%%%%%%%%%%%%%%%%%%%%%%%%%%%%%%%%%%%%%%%%%%%
\draw [overlay,<-,very thick,blue,opacity=.5,rounded corners=0mm]
(n) |-(n1); 
\node at (n1) [ blue,below,xshift=2mm ]{ 
$n_1$
}; %%%%%%%%%%%%%%%%%%%%%%%%%%%%%%%%%%%%%%%%%%%%%%%%%%%%%%%%%%%%%%%%%%%%%%%%%%%%%%%%%%%%%%%%%%%%%%%%%%%%%%%%%%%%%%%%%%%%%%%%%%%%%%%%%%%%%%%%%%%%%%%%
\draw [overlay,<-,very thick,blue,opacity=.5,rounded corners=0mm]
(nn) |-(nn1); 
\node at (nn1) [ blue,below,xshift=2mm ]{ 
$n_2$
}; 
\end{tikzpicture} 

\medskip
\medskip

%con esos datos podemos calcular su ecuación usando la formula:
%\[ (x - x_0, y - y_0) \cdot (n_1, n_2) = 0\]

}
(0.8,-.7)  node[draw=orange,text width=63mm,below right](c){
\begin{align*}
    \left(x -(-1), y - (-1) \rule{0mm}{4mm}\right) \cdot (3, 5 ) &= 0
    \\[2mm]
    (x+1, y +1) \cdot (3, 5 ) &= 0
    \\[2mm]
   3 (x+1)  +  5(y +1)&= 0
   \\[2mm]
   3 x +3   +  5y + 5  &= 0
   \\[2mm]
   3 x   +  5y + 8 &= 0
\end{align*}

}
(0,-1.5)  node[draw=red,text width=37mm,line width=2pt,below right](EcuImpli){
\textbf{Ecuación implicita} 
\[
 3x + 5y + 8 = 0
\] 

}
(.3,-1.9)  node[draw=orange,text width=85mm ,below right](d){
\small A partir de la ecuación implícita, realizamos manipulaciones algebraicas para obtener la ecuación segmentaria.
\begin{align*}
 3x + 5y + 8&= 0
 \\
  3x + 5y  &= -8
  \\
  -3x - 5y  &= 8
   \\
  \frac{-3x}{8} - \frac{ 5y}{8}  &= 1 
  \\
 \dfrac{\, x\,}{-\tfrac{8}{3}}  + \dfrac{y}{-\tfrac{8}{5}}   &= 1 
\end{align*}


}
(1.5,-2.3)  node[draw=red,text width=43mm,line width=2pt,below right](EcuSeg){
\textbf{Ecuación segmentaria}  
\[
\dfrac{\, x\,}{-\tfrac{8}{3}}  + \dfrac{y}{-\tfrac{8}{5}}   = 1
\]

};%--------------%------------------%----------------------------%------------------------------------%-----------------------%-----------------------------------------------------%--------------------------------%-----------------------------%----------------------------%----------------------------%-------------------------------------------------------
\draw[thick,->,dashed,color=teal] (bb) to[out=-30,in=20]  node [pos=0.7,text width=43mm , right,xshift=-1mm
] {\footnotesize
\qquad Usaremos la formula:
\begin{align*}
     (x - x_0, y - y_0) \cdot (n_1, n_2) &= 0
\end{align*}
}(c);
\draw[thick,->,dashed,color=teal] (c) to[out=180, in=60] (EcuImpli);
\draw[thick,->,dashed,color=teal] (EcuImpli) to[out=-125, in=180] (d);
\draw[thick,->,dashed,color=teal] (d) to (EcuSeg);
\end{tikzpicture}

 


\end{enumerate}












\item  Hallar la ecuación de la recta que cumple con las siguientes condiciones y graficarlas:

\begin{enumerate}[(a)]
    \item Pasa por el origen de coordenadas y el punto $P = (-3,-1)$.
    \item  Pasa por el punto $A = (-1,-2)$ y es paralela a la recta $L_1 : y - 3x - 9 = 0$.
    \item  Pasa por el punto $B = (3,-2)$ y tiene un ángulo de inclinación $\hat{\theta}$, con $\hat{\theta} = 135^\circ$.
    \item  Forma sobre el semieje positivo $x$ un segmento de longitud 7 y pasa por el punto de abscisa 4 de la recta $M : 5x + 2y = 30$.
    \item  Es perpendicular al segmento que forman los puntos $P = (2,1)$ y $Q = (-1,5)$ y pasa por el punto $R = (0,-3)$.
\end{enumerate}





\section*{Solución}

\begin{enumerate}[(a)]
\item Pasa por el origen y por $P = (-3,-1)$. Con lo cual tenemos que pasa por: $O=(0\indicador{a}{(-.1,-.1)}{a1}{(0,-.4)} ,0\indicador{b}{(-.1,-.1)}{b1}{(.1,-.4)})$ y $P = (-3\indicador{x}{(-.1,-.1)}{x1}{(.1,-.4)},-1\indicador{y}{(-.1,-.1)}{y1}{(.1,-.4)})$.

\begin{tikzpicture}[overlay,remember picture] 
\node at (a1) [below, blue,xshift=1.5mm]{ 
$x_0$
};
\draw [overlay,<-,very thick,blue,opacity=.5]
(a) --(a1);  %%%%%%%%%%%%%%%%%%%%%%%%%%%%%%%%%%%%%%%%%%%%%%%%%%%%%%%%%%%%%%%%%%%%%%%%%%%%%%%%%%%%%%%%%%%%%%%%%%%%%%%%%%%%%%%%%%%%%%%%%%%%%%%%%%%%%%%%%%%%%%%%
\node at (b1) [below, blue,xshift=1.5mm]{ 
$y_0$
};
\draw [overlay,<-,very thick,blue,opacity=.5]
(b) --(b1); 
%%%%%%%%%%%%%%%%%%%%%-----------------------------------------------------------------------------------------------------------------------%%%%%%%%%%%%%%%%%-----------------------------------------------------------%%%%%%%%%%%%%%%%%%%%%%%-----------------------------------------------------------%%%%%%%%%%%%%%%%%%%%----------------------------------------------------------------------------------------------------------------------%%%%%%%%%%%%%%%%%%%%%%%%%%%%%%-----------------------------------------------------------%%%%%%%%%%%----------------------------------------------------------------------------------------------------------------------%%%%%%%%%%%%%%%%%%%%%%
\node at (x1) [below, blue,xshift=1.5mm]{ 
$x_1$
};
\draw [overlay,<-,very thick,blue,opacity=.5]
(x) --(x1);  %%%%%%%%%%%%%%%%%%%%%%%%%%%%%%%%%%%%%%%%%%%%%%%%%%%%%%%%%%%%%%%%%%%%%%%%%%%%%%%%%%%%%%%%%%%%%%%%%%%%%%%%%%%%%%%%%%%%%%%%%%%%%%%%%%%%%%%%%%%%%%%%
\node at (y1) [below, blue,xshift=1.5mm]{ 
$y_1$
};
\draw [overlay,<-,very thick,blue,opacity=.5]
(y) --(y1); 
\end{tikzpicture}










La ecuación de la recta que pasa por dos puntos es: $y = \frac{y_1 - y_0}{x_1-x_0} \,  (x -x_0 ) + y_0 $
    \[
         y = \frac{y_1 - y_0}{x_1-x_0} \,  (x -x_0 ) + y_0  
         \quad
         \stackrel[\substack{\text{Sustituyendo}\\ \text{los valores}}]{}{\Longleftrightarrow}
         \quad
         y =  \frac{-1 - 0}{-3 - 0}    \,  (x - 0 ) + 0 
         \quad
         \stackrel[\text{resolviendo}]{}{\Longleftrightarrow} 
         \quad 
          y = \tfrac{1}{3}x
    \]
      

 \begin{center}
     
\definecolor{rojo}{rgb}{0.7,0.4,0.4}
\definecolor{azul}{rgb}{.4,.4,0.9}
 
\definecolor{gris}{rgb}{0.5,0.5,0.5}
\begin{tikzpicture}[scale=0.7]
% extremos para los ejes
\xdef\ext{4}
\xdef\yext{3}
\xdef\f{-4}
\xdef\g{-2}
\pgfmathtruncatemacro{\c}{\f+1}
\pgfmathtruncatemacro{\d}{\g+1}
\draw [color=gris!30!,dash pattern=on 1pt off 1pt, xstep=1cm,ystep=1cm] (\f-.3,\g-.3) grid (\ext+.3,\yext+.3);

\draw[->] (\f-.3,0) -- (\ext+.3,0) ;

\draw[->] (0,\g-.3) -- (0,\yext+.3) ;

\foreach \x in {\f,...,-1,1,2,...,\ext}
\draw[shift={(\x,0)},color=black] (0pt,2pt) -- (0pt,-2pt) node[below] {\tiny $\x$};

\foreach \y in {\g,...,-1,1,2,...,\yext }
\draw[shift={(0,\y)},color=black] (2pt,0pt) -- (-2pt,0pt) node[left] {\tiny $\y$};


    \coordinate (P) at (-3,-1);%punto a la derecha
\coordinate (O) at (0,0);%punto a la izuierda
\coordinate (M) at ($1.45*(P) -.45*(O)$); 
\coordinate (m) at ($-1.43*(P)+2.43*(O)$);%Punto por detras de B



%%%% Recta que pasa por A y B
\filldraw[red] (P) circle (2pt) node[above left] {$P(-3,-1)$};
    \filldraw[red] (O) circle (2pt) node[above right] {$O$};
 
% Dibujo de la línea entre A y B
\draw[thick,blue] (m) -- (M);
    
\end{tikzpicture}

 \end{center}
\hfill{\bf Rta: } $y = \tfrac{1}{3} x $.


\medskip 

\item Pasa por $A = (-1,-2)$ y es paralela a $y - 3x - 9 = 0$

\begin{minipage}{.4\textwidth}
{\bf Solución 1:}

Reescribimos la ecuación de la recta dada:
\[ y = 3x + 9 \]
Como las rectas paralelas tienen la misma pendiente, la ecuación de la recta que buscamos es:
\[ y = 3x + b \]
Usamos el punto $A=(-1,-2)$ para determinar $b$:
\[ -2 = 3(-1) + b \Rightarrow b = 1 \]
Entonces, la ecuación es:
\[ y = 3x + 1 \]

\hfill{\bf Rta: } $y = 3x + 1 $.
\end{minipage}\qquad\begin{minipage}{.5\textwidth}
{\bf Solución 2:}

La recta $\ \indicador{aa}{(-.1,-.15)}{aa1}{(-.2,-.4)}y - 3\indicador{ab}{(-.1,-.15)}{ab1}{(-.2,-.4)} \,x - 9 = 0$ tiene normal $\vec{n}= (-3,1)$.

\bigskip   
\bigskip  

Como las rectas paralelas tienen el mismo normal, tenemos que la recta que buscamos pasa por \linebreak $A=(-1\indicador{a}{(-.1,-.15)}{a1}{(-.2,-.4)}
  ,-2\indicador{b}{(-.1,-.1)}{b1}{(.1,-.4)}
  )$ y tiene normal $\vec{n} = (-3\indicador{u}{(-.1,-.1)}{u1}{(-.2,-.4)} ,1 \indicador{uu}{(-.1,-.1)}{uu1}{(.1,-.4)} 
  )$.

\begin{tikzpicture}[overlay,remember picture]
\node at (aa1) [ blue,below  ]{ 
$n_y$
};
\draw [overlay,<-,very thick,blue,opacity=.5]
(aa) |-(aa1);  %%%%%%%%%%%%%%%%%%%%%%%%%%%%%%%%%%%%%%%%%%%%%%%%%%%%%%%%%%%%%%%%%%%%%%%%%%%%%%%%%%%%%%%%%%%%%%%%%%%%%%%%%%%%%%%%%%%%%%%%%%%%%%%%%%%%%%%%%%%%%%%%
\node at (ab1) [ blue,below]{ 
$n_x$
};
\draw [overlay,<-,very thick,blue,opacity=.5]
(ab) |-(ab1); 
%%%%%%%%%%%%%%%%%%%%%%%%%%%%%%%%%%%%%%%%%%%%%%%%%%%%%%%%%%%%%%%%%%%%%%%%%%%%%%%%%%%%%%%%%%%%%%%%%%%%%%%%%%%%%%%%%%%%%%%%%%%%%%%%%%%%%%%%%%%%%%%%%%%%%%%%%%%%%%%%%%%%%%%%%%%%%%%%%%%%%%%%%%%%%%%%%%%%%%%%%%%%%%%%%%%%%%%%%%%%%%%%%%%%%%%%%%%%%%%%%%%%%%%%%%%%%%%%%%%%%%%%%%%%%%%%%%%%%%%%%%%%%%%%
\node at (a1) [ blue,below  ]{ 
$x_0$
};
\draw [overlay,<-,very thick,blue,opacity=.5]
(a) |-(a1);  %%%%%%%%%%%%%%%%%%%%%%%%%%%%%%%%%%%%%%%%%%%%%%%%%%%%%%%%%%%%%%%%%%%%%%%%%%%%%%%%%%%%%%%%%%%%%%%%%%%%%%%%%%%%%%%%%%%%%%%%%%%%%%%%%%%%%%%%%%%%%%%%
\node at (b1) [ blue,below]{ 
$y_0$
};
\draw [overlay,<-,very thick,blue,opacity=.5]
(b) |-(b1); 
%%%%%%%%%%%%%%%%%%%%%%%%%%%%%%%%%%%%%%%%%%%%%%%%%%%%%%%%%%%%%%%%%%%%%%%%%%%%%%%%%%%%%%%%%%%%%%%%%%%%%%%%%%%%%%%%%%%%%%%%%%%%%%%%%%%%%%%%%%%%%%%%
\draw [overlay,<-,very thick,blue,opacity=.5]
(u) |-(u1); 
\node at (u1) [ blue,below,xshift=2mm ]{ 
$n_x$
}; %%%%%%%%%%%%%%%%%%%%%%%%%%%%%%%%%%%%%%%%%%%%%%%%%%%%%%%%%%%%%%%%%%%%%%%%%%%%%%%%%%%%%%%%%%%%%%%%%%%%%%%%%%%%%%%%%%%%%%%%%%%%%%%%%%%%%%%%%%%%%%%%
\draw [overlay,<-,very thick,blue,opacity=.5]
(uu) |-(uu1); 
\node at (uu1) [ blue,below,xshift=2mm ]{ 
$n_y$
}; 
\end{tikzpicture} 

\bigskip  



Usaremos la formula: \quad $(x - x_0, y - y_0) \cdot (n_x, n_y) = 0$

\begin{align*}
    \left(x -(-1), y - (-2) \rule{0mm}{4mm}\right) \cdot (-3,1 ) &= 0
    \\[2mm]
    (x+1, y +2) \cdot (-3,1) &= 0
    \\[2mm]
   -3 (x+1)  +  1(y +2)&= 0
   \\[2mm]
   -3 x-3   +  y + 2 &= 0
   \\[2mm]
   -3x   +  y - 1&= 0
\end{align*}

  \hfill{\bf Rta: } $-3x   +  y - 1 \, =\, 0$.
\end{minipage}






\begin{center}
    

\definecolor{gris}{rgb}{0.5,0.5,0.5}
\begin{tikzpicture}[scale=0.5]
% extremos para los ejes
\xdef\ext{4}
\xdef\yext{11}
\xdef\f{-4}
\xdef\g{-3}
\pgfmathtruncatemacro{\c}{\f+1}
\pgfmathtruncatemacro{\d}{\g+1}
\draw [color=gris!30!,dash pattern=on 1pt off 1pt, xstep=1cm,ystep=1cm] (\f-.3,\g-.3) grid (\ext+.3,\yext+.3);

\draw[->] (\f-.3,0) -- (\ext+.3,0) ;

\draw[->] (0,\g-.3) -- (0,\yext+.3) ;

\foreach \x in {\f,...,-1,1,2,...,\ext}
\draw[shift={(\x,0)},color=black] (0pt,2pt) -- (0pt,-2pt) node[below] {\tiny $\x$};

\foreach \y in {\g,...,-1,1,2,...,\yext }
\draw[shift={(0,\y)},color=black] (2pt,0pt) -- (-2pt,0pt) node[left] {\tiny $\y$};


    \draw[thick,blue,domain=-4:.8] plot ({\x}, {3*\x+9}) node[above left] {$y-3x-9=0$};
    \draw[thick,blue,domain=-1.33:3.4] plot ({\x}, {3*\x+1}) node[above right] {$-3x+y-1=0$};
    \filldraw[red] (-1,-2) circle (2pt) node[ right] {$A(-1,-2)$};
\end{tikzpicture}

\end{center}













\item Pasa por $B = (3,-2)$ y tiene inclinación $135^\circ$

\begin{minipage}{.5\textwidth}
La pendiente es:
\[ m = \tan(135^\circ) = -1 \]
Así tenemos que la recta pasa por $B = (3,-2)$ y tiene pendiente $m=-1$, entonces usaremos la formula pendiente punto:\quad $ y = m(x - x_0) + y_0 .$
\begin{align*}
    y &= -1 \, (x - 3) + (-2)
    \\
    y &= -x +3 -2
    \\
    y &= -x +1
\end{align*}

\end{minipage}\qquad\begin{minipage}{.3\textwidth}


\definecolor{gris}{rgb}{0.5,0.5,0.5}
\begin{tikzpicture}[scale=0.6]
% extremos para los ejes
\xdef\ext{4}
\xdef\yext{4}
\xdef\f{-4}
\xdef\g{-3}
\pgfmathtruncatemacro{\c}{\f+1}
\pgfmathtruncatemacro{\d}{\g+1}
\draw [color=gris!30!,dash pattern=on 1pt off 1pt, xstep=1cm,ystep=1cm] (\f-.3,\g-.3) grid (\ext+.3,\yext+.3);

\draw[->] (\f-.3,0) -- (\ext+.3,0) ;

\draw[->] (0,\g-.3) -- (0,\yext+.3) ;

\foreach \x in {\f,...,-1,1,2,...,\ext}
\draw[shift={(\x,0)},color=black] (0pt,2pt) -- (0pt,-2pt) node[below] {\tiny $\x$};

\foreach \y in {\g,...,-1,1,2,...,\yext }
\draw[shift={(0,\y)},color=black] (2pt,0pt) -- (-2pt,0pt) node[left] {\tiny $\y$};

    \draw[thick,blue,domain=4.3:-3.3] plot ({\x}, {-\x+1}) node[above left] {$y=-x+1$};
    \filldraw[red] (3,-2) circle (2pt) node[ right] {$B(3,-2)$};
\end{tikzpicture}
\end{minipage}

\hfill{\bf Rta: } $y = -x + 1 $.

\medskip 

\item Forma sobre el semieje positivo $x$ un segmento de longitud 7 y pasa por el punto de abscisa 4 de la recta $M : 5x + 2y = 30$.

\begin{itemize}
\begin{multicols}{2}
    \item Como forma sobre el semieje positivo $x$ un segmento de longitud 7, tenemos que la recta pasa por el punto $(7,0)$

    \columnbreak % Forzar cambio de columna si es necesario

    \item Para determinar el punto de abscisa 4 en la recta $M$, sustituimos $x=4$ en su ecuación y despejamos $y$:
\begin{align*}
    5(4) + 2\, y &= 30
    \\
    20 + 2\, y &= 30
    \\
     2\, y &= 30 -20
     \\
     y &= \frac{10}{2}
     \\
     y &= 5.
\end{align*}
con la cual el punto de abscisa 4 tiene ordenada 5, es decir la recta que buscamos pasa por el punto $(4,5)$



\end{multicols}
\end{itemize}

La recta que buscamos pasa por los puntos $(7,0)$ y $(4,5)$, usamos la ecuacion de la recta que pasa por dos puntos:
    \[
         y = \frac{y_1 - y_0}{x_1-x_0} \,  (x -x_0 ) + y_0  
         \quad
         \stackrel[\substack{\text{Sustituyendo}\\ \text{los valores}}]{}{\longleftrightarrow}
         \quad
         y =  \frac{5-0}{4-7}    \,  (x - 7 ) +0 
         \quad
         \stackrel[\text{resolviendo}]{}{\longleftrightarrow} 
         \quad 
          y = -\tfrac{5}{3} x + \frac{35}{3}
    \]


\definecolor{gris}{rgb}{0.5,0.5,0.5}
\begin{tikzpicture}[scale=0.6]
% extremos para los ejes
\xdef\ext{9}
\xdef\yext{7}
\xdef\f{-1}
\xdef\g{-2}
\pgfmathtruncatemacro{\c}{\f+1}
\pgfmathtruncatemacro{\d}{\g+1}
\draw [color=gris!30!,dash pattern=on 1pt off 1pt, xstep=1cm,ystep=1cm] (\f-.3,\g-.3) grid (\ext+.3,\yext+.3);

\draw[->] (\f-.3,0) -- (\ext+.3,0) ;

\draw[->] (0,\g-.3) -- (0,\yext+.3) ;

\foreach \x in {\f,...,-1,1,2,...,\ext}
\draw[shift={(\x,0)},color=black] (0pt,2pt) -- (0pt,-2pt) node[below] {\tiny $\x$};

\foreach \y in {\g,...,-1,1,2,...,\yext }
\draw[shift={(0,\y)},color=black] (2pt,0pt) -- (-2pt,0pt) node[left] {\tiny $\y$}; 


    \draw[thick,blue,domain=8.66:2.66] plot ({\x}, {-5/3*(\x-7)}) node[above right] {$y=-\frac{5}{3}x+\frac{35}{3}$};
    \filldraw[red] (7,0) circle (2pt) node[below] {$7$};
    \filldraw[red] (4,5) circle (2pt) node[below left] {$(4,5)$};
\end{tikzpicture}

\vspace{-9mm}

\hfill{\bf Rta: } $y = - \tfrac{5}{3} x + \frac{35}{3}$.


\bigskip

\item Es perpendicular al segmento que forman los puntos $P = (2,1)$ y $Q = (-1,5)$,  pasando por $R=(0,-3)$

\begin{itemize}
\begin{minipage}{.3 \textwidth}
    

    \item El vector normal es $\overrightarrow{PQ}$:
\begin{align*}
     \overrightarrow{PQ}&= (-1,5)- (2,1) 
     \\
     &= \left(-1-2  ,\  5 -1 \right) 
    \\
     &= \left( -3 ,\, 4\right) 
\end{align*}

\, \\
\, \\
 \, \\
\, \\  
    \end{minipage}\qquad\begin{minipage}{.6 \textwidth}
    


\item Tenemos que la recta que buscamos pasa por \linebreak $R=(0\indicador{a}{(-.1,-.15)}{a1}{(-.2,-.4)}
  ,-3\indicador{b}{(-.1,-.1)}{b1}{(.1,-.4)}
  )$ y tiene normal $\overrightarrow{PQ} = (-3\indicador{u}{(-.1,-.1)}{u1}{(-.2,-.4)} , 4\indicador{uu}{(-.1,-.1)}{uu1}{(.1,-.4)} 
  )$.

\begin{tikzpicture}[overlay,remember picture] 
%%%%%%%%%%%%%%%%%%%%%%%%%%%%%%%%%%%%%%%%%%%%%%%%%%%%%%%%%%%%%%%%%%%%%%%%%%%%%%%%%%%%%%%%%%%%%%%%%%%%%%%%%%%%%%%%%%%%%%%%%%%%%%%%%%%%%%%%%%%%%%%%%%%%%%%%%%%%%%%%%%%%%%%%%%%%%%%%%%%%%%%%%%%%%%%%%%%%%%%%%%%%%%%%%%%%%%%%%%%%%%%%%%%%%%%%%%%%%%%%%%%%%%%%%%%%%%%%%%%%%%%%%%%%%%%%%%%%%%%%%%%%%%%%
\node at (a1) [ blue,below  ]{ 
$x_0$
};
\draw [overlay,<-,very thick,blue,opacity=.5]
(a) |-(a1);  %%%%%%%%%%%%%%%%%%%%%%%%%%%%%%%%%%%%%%%%%%%%%%%%%%%%%%%%%%%%%%%%%%%%%%%%%%%%%%%%%%%%%%%%%%%%%%%%%%%%%%%%%%%%%%%%%%%%%%%%%%%%%%%%%%%%%%%%%%%%%%%%
\node at (b1) [ blue,below]{ 
$y_0$
};
\draw [overlay,<-,very thick,blue,opacity=.5]
(b) |-(b1); 
%%%%%%%%%%%%%%%%%%%%%%%%%%%%%%%%%%%%%%%%%%%%%%%%%%%%%%%%%%%%%%%%%%%%%%%%%%%%%%%%%%%%%%%%%%%%%%%%%%%%%%%%%%%%%%%%%%%%%%%%%%%%%%%%%%%%%%%%%%%%%%%%
\draw [overlay,<-,very thick,blue,opacity=.5]
(u) |-(u1); 
\node at (u1) [ blue,below,xshift=2mm ]{ 
$n_x$
}; %%%%%%%%%%%%%%%%%%%%%%%%%%%%%%%%%%%%%%%%%%%%%%%%%%%%%%%%%%%%%%%%%%%%%%%%%%%%%%%%%%%%%%%%%%%%%%%%%%%%%%%%%%%%%%%%%%%%%%%%%%%%%%%%%%%%%%%%%%%%%%%%
\draw [overlay,<-,very thick,blue,opacity=.5]
(uu) |-(uu1); 
\node at (uu1) [ blue,below,xshift=2mm ]{ 
$n_y$
}; 
\end{tikzpicture} 

\bigskip  



Usaremos la formula: \quad $(x - x_0, y - y_0) \cdot (n_x, n_y) = 0$

\begin{align*}
    \left(x -0, y - (-3)\rule{0mm}{3.5mm}\right) \cdot (-3, 4 ) &= 0
    \\[2mm]
    (x, y + 3 ) \cdot (-3,4) &= 0
    \\[2mm]
   -3 x  +  4 (y +3)&= 0
   \\[2mm]
   -3 x + 4y +12 &= 0 
\end{align*} 


\end{minipage}
\end{itemize}

\medskip

  \hfill{\bf Rta: } $-3 x + 4y +12 \, =\, 0$.


\vspace{-18mm}
\begin{tikzpicture}[scale=0.4]
% extremos para los ejes
\xdef\ext{5}
\xdef\yext{6}
\xdef\f{-5}
\xdef\g{-5}
\pgfmathtruncatemacro{\c}{\f+1}
\pgfmathtruncatemacro{\d}{\g+1}
\draw [color=gris!30!,dash pattern=on 1pt off 1pt, xstep=1cm,ystep=1cm] (\f-.3,\g-.3) grid (\ext+.3,\yext+.3);

\draw[->] (\f-.3,0) -- (\ext+.3,0) ;

\draw[->] (0,\g-.3) -- (0,\yext+.3) ;

\foreach \x in {\f,...,-1,1,2,...,\ext}
\draw[shift={(\x,0)},color=black] (0pt,2pt) -- (0pt,-2pt) node[below] {\tiny $\x$};

\foreach \y in {\g,...,-1,1,2,...,\yext }
\draw[shift={(0,\y)},color=black] (2pt,0pt) -- (-2pt,0pt) node[left] {\tiny $\y$}; 



% 

\coordinate (P) at (2,1);%punto a la derecha
\coordinate (Q) at (-1,5);%punto a la izuierda 
% Dibujo de la línea entre P y Q 
\draw[blue,fill=blue] (P) circle (2pt)  node[above right] {$P$};

\draw[blue,fill=blue] (Q) circle (2pt) node[below left] {$Q$};
\draw[thick,teal] (P) -- (Q);

    \draw[thick,blue,domain=-3:5.3] plot ({\x}, {3/4*\x-3}) node[above right] {$-3 x + 4y +12= 0$};
    \filldraw[red] (0,-3) circle (2pt) node[below right] {$R=(0,-3)$};
\end{tikzpicture}








\end{enumerate}






    
\item  Determinar:

\begin{itemize}
    \item[(a)] Dos puntos de la recta $L : 3x - 2y - 3 = 0$.
    \item[(b)] Si los puntos $P = (-3,-1)$ y $Q = \left(-\frac{1}{2}, 3\right)$ pertenecen a la recta $M : -2x + y = 4$.
    \item[(c)] Si los puntos $A = (6,3)$, $B = (-3,3)$ y $C = (0,-1)$ están alineados. (Resolver este inciso de dos maneras diferentes).
\end{itemize}




\textbf{Solución:}

\begin{enumerate}[(a)]
   
\item  
Dada la ecuación de la recta $L: 3x - 2y - 3 = 0$, determinamos dos puntos:

\begin{itemize}
\begin{multicols}{2} 
    
    \item Para $x = 0$:
\begin{align*}
    3(0) - 2y - 3 &= 0 \\
    -2y &= 3 \\
    y &= -\frac{3}{2}
\end{align*}
Por lo que el primer punto es $P_1 = (0, -\frac{3}{2})$.

\item Para $y = 0$:
\begin{align*}
    3x - 2(0) - 3 &= 0 \\
    3x &= 3 \\
    x &= 1
\end{align*}
Por lo que el segundo punto es $P_2 = (1,0)$.

\end{multicols}
\end{itemize}


\begin{center}
    



\definecolor{rojo}{rgb}{0.7,0.4,0.4}
\definecolor{azul}{rgb}{.4,.4,0.9}
 
\definecolor{gris}{rgb}{0.5,0.5,0.5}
\begin{tikzpicture}[line cap=round,line join=round,>=triangle 45,scale=.75]

% extremos para los ejes
\xdef\ext{4}
\xdef\yext{3}
\xdef\f{-2}
\xdef\g{-2}
\pgfmathtruncatemacro{\c}{\f+1}
\pgfmathtruncatemacro{\d}{\g+1}
\draw [color=gris!30!,dash pattern=on 1pt off 1pt, xstep=1cm,ystep=1cm] (\f-.3,\g-.3) grid (\ext-.3,\yext+.3);

\draw[->] (\f-.3,0) -- (\ext+.3,0) ;

\draw[->] (0,\g-.3) -- (0,\yext+.3) ;

\foreach \x in {\f,...,-1,1,2,...,\ext}
\draw[shift={(\x,0)},color=black] (0pt,2pt) -- (0pt,-2pt) node[below] {\tiny $\x$};

\foreach \y in {\g,...,-1,1,2,...,\yext }
\draw[shift={(0,\y)},color=black] (2pt,0pt) -- (-2pt,0pt) node[left] {\tiny $\y$};


% 

\coordinate (A) at (0,-3/2);%punto a la izuierda
\coordinate (B) at (1,0);%punto a la derecha
\coordinate (M) at ($1.5*(A) -.5*(B)$);%Punto por atras de A
\coordinate (m) at ($-2.2*(A)+3.2*(B)$);%Punto por adelante de B


%%%%%Grafico de los puntos
%\draw[fill=black] (P) circle (1.5pt) node[shift=(-\ang: 5pt),right] {\scriptsize $P=(-2,1)$};


%%%% Recta que pasa por A y B
\draw[blue,fill=blue] (A) circle (2pt) node[ left] {$P_1$} ;

\draw[blue,fill=blue] (B) circle (2pt) node[above left] {$P_2$};
 
% Dibujo de la línea entre A y B
\draw[thick,blue] (M) -- (m) node[right] {$ L:\ 3x - 2y - 3 = 0 $};


\end{tikzpicture}
  
\end{center}  




\item 
La ecuación de la recta es $M: -2x + y = 4$. Verificamos si los puntos pertenecen a ella.

\begin{itemize}
\begin{multicols}{2}
    
    \item Para $P(-3,-1)$:
\begin{align*}
    -2(-3) + (-1) &= 6 - 1 = 5 \neq 4
\end{align*}
Por lo tanto, $P$ no pertenece a la recta.

\item Para $Q\left(-\frac{1}{2}, 3\right)$:
\begin{align*}
    -2\left(-\frac{1}{2}\right) + 3 &= 1 + 3 = 4
\end{align*}
Por lo tanto, $Q$ sí pertenece a la recta.

\end{multicols}
\end{itemize}

\begin{center}
\definecolor{gris}{rgb}{0.5,0.5,0.5}
\begin{tikzpicture}[line cap=round,line join=round,>=triangle 45,scale=.75]

% extremos para los ejes
\xdef\ext{2}
\xdef\yext{4}
\xdef\f{-4}
\xdef\g{-2}
\pgfmathtruncatemacro{\c}{\f+1}
\pgfmathtruncatemacro{\d}{\g+1}
\draw [color=gris!30!,dash pattern=on 1pt off 1pt, xstep=1cm,ystep=1cm] (\f-.3,\g-.3) grid (\ext-.3,\yext+.3);

\draw[->] (\f-.3,0) -- (\ext+.3,0) ;

\draw[->] (0,\g-.3) -- (0,\yext+.3) ;

\foreach \x in {\f,...,-1,1,2,...,\ext}
\draw[shift={(\x,0)},color=black] (0pt,2pt) -- (0pt,-2pt) node[below] {\tiny $\x$};

\foreach \y in {\g,...,-1,1,2,...,\yext }
\draw[shift={(0,\y)},color=black] (2pt,0pt) -- (-2pt,0pt) node[left] {\tiny $\y$};


%grafico de la recta
\draw[thick,purple,domain=-3.33:.15] plot ({\x}, {4+2*\x}) node[above right] {$-2 x+ y = 4$};
        \fill[red] (-3,-1) circle (2pt) node[below left] {$P$};
        \fill[blue] (-0.5,3) circle (2pt) node[left] {$Q$};
    \end{tikzpicture}
\end{center}

\item  Los puntos dados son $A = (6,3)$, $B = (-3,3)$ y $C = (0,-1)$. 


\begin{itemize}
    \item Con los tres puntos determinamos dos vectores que tengan un mismo punto inicial, si los vectores son paralelos, tenemos que los puntos están alineados.
\begin{align*}
    \overrightarrow{AB} &= (-3 - 6, 3 - 3) = (-9,0) \\
    \overrightarrow{AC} &=  (0 - 6, -1- 3) = (-6,-4)
\end{align*}

planteamos que para que resulten paralelos uno de ellos es multiplo escalar del otro:
\begin{align*}
    \overrightarrow{AB}  &= \alpha \overrightarrow{AC} \\
    (-9,0) &= \alpha  (-6,-4)\\
    &= (-6\alpha,\, -4 \alpha) , 
\end{align*}
de lo que se obtiene un absurdo, y resulta que $\alpha$ no existe, con lo cual los vectores no son paralelos  y lo puntos no están alineados.

Dado que el producto escalar no es cero, los vectores no son paralelos, lo que implica que los puntos no están alineados.




    \item  {\bf Otra forma de comprobarlo}  Están alineados si la pendiente entre dos de ellos es la misma que la pendiente con el tercero.

\begin{itemize}
\begin{multicols}{2}
    \item Pendiente entre $A$ y $B$:
\begin{align*}
    m_{AB} &= \frac{3 - 3}{6 - (-3)} = \frac{0}{9} = 0
\end{align*}

\item Pendiente entre $B$ y $C$:
\begin{align*}
    m_{BC} &= \frac{-1 - 3}{0 - (-3)} = \frac{-4}{3}
\end{align*}
\end{multicols}
\end{itemize}


Dado que $m_{AB} \neq m_{BC}$, los puntos no están alineados.

\bigskip

\noindent
\item {\bf Otra forma de comprobarlo} es sustituyendo los puntos en la ecuación de la recta que pasa por dos de ellos y verificando si el tercero la satisface. La ecuación de la recta que pasa por $B$ y $C$ es:
\begin{align*}
    y - 3 &= \frac{-4}{3} (x + 3)
\end{align*}
Verificamos si $A = (6,3)$ la satisface:
\begin{align*}
    3 - 3 &= \frac{-4}{3} (6 + 3) \\
    0 &= -12 \quad \text{(falso)}
\end{align*}
Por lo tanto, los puntos no están alineados.

\end{itemize}
\begin{center}
\definecolor{gris}{rgb}{0.5,0.5,0.5}
\begin{tikzpicture}[line cap=round,line join=round,>=triangle 45,scale=.5]

% extremos para los ejes
\xdef\ext{7}
\xdef\yext{3}
\xdef\f{-4}
\xdef\g{-2}
\pgfmathtruncatemacro{\c}{\f+1}
\pgfmathtruncatemacro{\d}{\g+1}
\draw [color=gris!30!,dash pattern=on 1pt off 1pt, xstep=1cm,ystep=1cm] (\f-.3,\g-.3) grid (\ext-.3,\yext+.3);

\draw[->] (\f-.3,0) -- (\ext+.3,0) ;

\draw[->] (0,\g-.3) -- (0,\yext+.3) ;

\foreach \x in {\f,...,-1,1,2,...,\ext}
\draw[shift={(\x,0)},color=black] (0pt,2pt) -- (0pt,-2pt) node[below] {\tiny $\x$};

\foreach \y in {\g,...,-1,1,2,...,\yext }
\draw[shift={(0,\y)},color=black] (2pt,0pt) -- (-2pt,0pt) node[left] {\tiny $\y$};


%
        \fill[red] (6,3) circle (3pt) node[above right] {$A$};
        \fill[purple] (-3,3) circle (3pt) node[above left] {$B$};
        \fill[violet] (0,-1) circle (3pt) node[below right] {$C$};
    \end{tikzpicture}
\end{center}

\end{enumerate}

\item  Dadas las ecuaciones $ax + (2 - b)y - 23 = 0$ y $(a - 1)x + by + 15 = 0$, hallar los valores de $a$ y $b$ para que representen rectas que se intersecan en el punto $R = (2,-3)$.

\textbf{Solución:}  

Las ecuaciones de las rectas son: 


\vspace{-18mm}
\begin{align*}
    ax + (2 - b)y - 23 &= 0, \\
    (a - 1)x + by + 15 &= 0.
\end{align*}
Dado que estas rectas se intersecan en el punto $R = (2,-3)$, sus coordenadas deben satisfacer ambas ecuaciones.

Sustituyendo $x = 2$ y $y = -3$ en las  ecuaciónes, tenemos lo siguiente:

\[
\left\{
\begin{array}{rl}
 a 2+ (2 - b)(-3) - 23 &= 0 \\
    (a - 1)2 + b(-3) + 15 &= 0
\end{array}
\right.
\quad
\Longleftrightarrow 
\quad 
\left\{
\begin{array}{rl}
 2a - 3(2 - b) - 23 &= 0 \\
2a - 2 - 3b + 15 &= 0
\end{array}
\right.
\quad
\Longleftrightarrow 
\quad 
\left\{
\begin{array}{rl}
2a - 6 + 3b - 23 &= 0 \\
 2a - 3b + 13 &= 0
\end{array}
\right.
\]


\bigskip 


$\Longleftrightarrow 
\quad 
\left\{
\begin{array}{rl}
2a + 3b  &= 29   \marka{a}{(.1,.25)}
\\[5mm]
 2a - 3b  &= -13 \indicador{b}{(.1,.15)}{c}{(.9,.45)} \marka{d}{(1.6,.45)}
\end{array}
\right.
$


\tikz[overlay,remember picture] {
\draw [overlay,->,very thick,red,opacity=.5](b)-|(c)--(d);  
\draw [overlay,very thick,red,opacity=.5](a)++(0,-.1)-|(c);
\node  at (d) [text width=55mm, below right , yshift= 7mm ]{
{\scriptsize Sumamos ambas ecuaciones:}

\vspace{-8mm}
\begin{align*}
    (2a + 3b) + (2a - 3b) &= 29 - 13, \\
    4a &= 16, \\
   \marka{a1}{(-.1,-.2)} 
   a &= 4.
   \marka{a2}{(.1,.45)} \indicador{b}{(.13,.15)}{c}{(1.6,1.45)} \marka{d}{(2.4,1.45)}
\end{align*}
};
\draw[color=black!50!,rounded corners=1mm] (a1) rectangle (a2);
%%%%%%%%%%%%%%%%%%%%%%%%%%%%%%%%%%%%%%%%%%%%%%%%%%%%%%%%%%%%%%%%%%%%%%%%%%%%%%%%%%%%%%%%%%%%%%%%%%%%%%%%%%%%%%%%%%%%%%%%%%%%%%%%%%%%%%%%%%%%%%%%%%%%%%%%%%%%%%%%%%%%%%%%%%%%%%%%%%%%
\draw [overlay,->,very thick,blue,opacity=.5](b)-|(c)--(d);  
\draw [overlay,very thick,blue,opacity=.5](a)-|(c);
\node  at (d) [text width=50mm, below right , yshift= 13mm ]{
{\scriptsize Sustituyendo $a = 4$ en la primera ecuación:}

\vspace{-8mm}
\begin{align*}
    2(4) + 3b &= 29, \qquad\\
    8 + 3b &= 29, \\
    3b &= 21, \\
     \marka{a1}{(-.1,-.2)}
    b &= 7.
   \marka{a2}{(.1,.45)} 
\end{align*}
};
\draw[color=black!50!,rounded corners=1mm] (a1) rectangle (a2);
}
%%%%%%%%%%%%%%%%%%%%%%%%%%%%%%%%%%%%%%%%%%%%%%%%%%%%%%%%%%%%%%%%%%%%%%%%%%%%%%%%%%%%%%%%%%%%%%%%%%%%%%%%%%%%%%%%%%%%%%%%%%%%%%%%%%%%%%%%%%%%%%%%%%%%%%%%%%%%%%%%%%%%%%%%%%%%%%%%%%%%



\vspace{11mm}


\hfill Por lo tanto, los valores de $a$ y $b$ son $a = 4$ y $b = 7$.

\textbf{Gráfico de las rectas: } al sustituir $a = 4$ y $b = 7$, nos queda las rectas: $4x - 5y - 23 = 0$ y  $3x + 7y + 15 = 0$.

\begin{center}
\definecolor{gris}{rgb}{0.5,0.5,0.5}
\begin{tikzpicture}[scale=0.8]
% extremos para los ejes
\xdef\ext{6}
\xdef\yext{2}4
\xdef\f{-2}
\xdef\g{-5}
\pgfmathtruncatemacro{\c}{\f+1}
\pgfmathtruncatemacro{\d}{\g+1}
\draw [color=gris!30!,dash pattern=on 1pt off 1pt, xstep=1cm,ystep=1cm] (\f-.3,\g-.3) grid (\ext+.3,\yext+.3);

\draw[->] (\f-.3,0) -- (\ext+.3,0) ;

\draw[->] (0,\g-.3) -- (0,\yext+.3) ;

\foreach \x in {\f,...,-1,1,2,...,\ext}
\draw[shift={(\x,0)},color=black] (0pt,2pt) -- (0pt,-2pt) node[below] {\tiny $\x$};

\foreach \y in {\g,...,-1,1,2,...,\yext }
\draw[shift={(0,\y)},color=black] (2pt,0pt) -- (-2pt,0pt) node[left] {\tiny $\y$};
    
    % Primera recta: 4x - 5y - 23 = 0 
    \draw[thick,red,domain=-.99:6.4] plot ({\x}, {(4*\x-23)/5})  node[above right] {$4x - 5y - 23 = 0$};
    
    % Segunda recta: 3x + 7y + 15 = 0
    \draw[thick,blue,domain=-2.3:6.3] plot ({\x}, {(3*\x+15)/(-7)}) node[right] {$3x + 7y + 15 = 0$};
    
    % Punto de intersección R = (2,-3)
    \filldraw[black] (2,-3) circle (2pt) node[right] {$R=(2,-3)$};
\end{tikzpicture}
\end{center}

 










\item Calcular, si es posible, los valores de $h$ y $k$ para que las rectas de ecuaciones $3x + hy + 4 = 0$ y $2x - 4y + k = 0$ sean:

\begin{enumerate}[(a)]
\begin{multicols}{3}
    \item  Paralelas no coincidentes.
    \item Perpendiculares.
    \item Coincidentes.
\end{multicols}
\end{enumerate}


\textbf{Solución:}

\begin{comment}
Consideremos el siguiente sistema de ecuación: \quad $\left\{ \begin{array}{rl}
    3x + hy + 4 &= 0
    \\
    2x - 4y + k &= 0
\end{array}
\right.$
    
\end{comment}
 






\begin{enumerate}[(a)]
    
    \item \textbf{Rectas paralelas no coincidentes}\\
    Dos rectas son paralelas si sus normales son paralelos.
    \begin{itemize}
        \item El normal de la primera recta es: $\Vec{n}_1= (3,h)$
        \item  El normal de la segunda recta es: $\Vec{n}_2 = (2,-4)$
    \end{itemize}
    Veamos la condición de paralelismo en este caso:    
    \begin{align*}
    \Vec{n}_1 = t \Vec{n}_2, \text{ para algún } t\in \R,
    \quad
    &\Longleftrightarrow
    \quad 
     (3,h) = t  (2,-4) \hspace{55mm}\,
     \\[4mm]
    &\Longleftrightarrow
    \quad 
    \left\{
    \begin{array}{rl}
         3=& 2t  \indicador{a}{(.1,.15)}{b}{(1.1,.15)}
         \\[3mm]
        h & = -4 t \marka{c}{(.1,.15)} 
    \end{array}
    \right.
    \end{align*}

\tikz[overlay,remember picture] {
\draw [overlay,->,very thick,red,opacity=.5](a)--(b);
\node  at (b) [ right ]{
$\marka{a1}{(-.1,-.2)} 
t
=  %%%%%
\tfrac{3}{2}  \marka{a2}{(.1,.45)}
$ \indicador{x}{(.1,.15)}{y}{(.8,-.35)}\marka{z}{(1.3,-.35)}
};
%%%%%%%%%%%%%%%%%%%%%%%%
\draw[color=black!50!,rounded corners=1mm] (a1) rectangle (a2);
\draw [overlay,->,very thick,red,opacity=.5](x)-|(y)--(z);
\draw [overlay,very thick,red,opacity=.5](c)-|(y);
\node  at (z) [text width=20mm, below right , yshift=9mm ]{
\[
\begin{array}{rl}
  h
 &=%%%%%%%%%%%%%%%%%%%%%%
 -4 \cdot \tfrac{3}{2}
 \\[2mm]%%%%%%%%%%%%%%%%%%%%%%%%%%%%%%%%%%%%%%%%%%%%%%%%%%%%%%%%%%%%%%%%%%
  \marka{a1}{(-.1,-.2)} 
   h
  &= %%%%%%%%%%%%%%%%%%%%%%
   -6\marka{a2}{(.1,.45)}
\end{array}
\]
};
%%%%%%%%%%%%%%%%%%%%%%%%
\draw[color=black!50!,rounded corners=1mm] (a1) rectangle (a2);
}

    
    Para que no sean coincidentes, sus ecuaciones   no deben ser proporcionales. Ya tenemos que los primeros coeficientes son proporcionales para $t= \frac{3}{2}$, entonces la proporción entre los términos independientes deve ser distinta a $\tfrac{3}{2}$ :
    \begin{equation*}
        \frac{4}{k} \neq \frac{3}{2} \Rightarrow k \neq \frac{8}{3}
    \end{equation*}
    
\hfill {\bf Rta:} para que las rectas sean paralelas, $h=-6$ y $ k \neq \frac{8}{3}$

    \item \textbf{Rectas perpendiculares}\\
    Dos rectas son perpendiculares si sus normales son perpendiculares:
   
    
    Veamos la condición de paralelismo en este caso:   

    \[
    \Vec{n}_1\cdot \Vec{n}_2 = 0
    \quad   
    \Longleftrightarrow   
    \quad    
    (3,h) \cdot  (2,-4)= 0
     \quad   
    \Longleftrightarrow   
    \quad    
    6 - 4 h = 0
    \quad   
    \Longleftrightarrow   
    \quad    
    6 = 4 h 
    \quad   
    \Longleftrightarrow   
    \quad    
    h = \frac{6}{4} = \frac{3}{2}
    \]
\hfill {\bf Rta:} Para que las rectas sean perpendiculares, $h  = \tfrac{3}{2}$.
    \item \textbf{Rectas coincidentes}\\
    Dos rectas son coincidentes si sus ecuaciones son proporcionales:
    \begin{equation*}
        \frac{3}{2} = \frac{h}{-4} = \frac{4}{k}
    \end{equation*}
    De aquí, despejamos:
    \begin{equation*}
        h = -6, \quad k = \frac{8}{3}
    \end{equation*}
     
\hfill {\bf Rta:} para que las rectas sean coincidentes, $h=-6$ y $ k = \frac{8}{3}$

\end{enumerate}







\item Sean $A = (2,-2)$, $B = (6,-1)$ y $C = (9,3)$ tres puntos de $\mathbb{R}^2$:

\begin{itemize}
    \item[(a)] Hallar la ecuación de las rectas que contienen los lados del paralelogramo $ABCD$ y determinar las coordenadas del punto $D$ (tener en cuenta que se nombra el paralelogramo en sentido anti-horario).
    \item[(b)] Escribir la ecuación vectorial de la recta que contiene a la diagonal $AC$.
\end{itemize}




\textbf{Recordemos:}


{\bf Definición de Paralelogramo: } Un \textbf{paralelogramo} es un cuadrilátero cuyos lados opuestos son \textbf{paralelos}. Es decir, si un cuadrilátero tiene sus lados opuestos paralelos, entonces es un paralelogramo.

\begin{multicols}{2}
Matemáticamente, si el cuadrilátero $ABCD$ es un paralelogramo, entonces:

\[
\overleftrightarrow{AB} \paralela \overleftrightarrow{CD} \quad \text{y} \quad \overleftrightarrow{BC} \paralela \overleftrightarrow{DA}
\]


\bigskip

\begin{center}
\begin{tikzpicture}[scale=.8]
    % Dibujar el paralelogramo
    \draw[thick] (0,0) -- (4,1) -- (5,3) -- (1,2) -- cycle;
    
    % Etiquetas de los vértices
    \node[left] at (0,0) {\textbf{A}};
    \node[right] at (4,1) {\textbf{B}};
    \node[right] at (5,3) {\textbf{C}};
    \node[left] at (1,2) {\textbf{D}};

    
    % Marcas de igualdad de lados
    \draw[dashed] (0,0) -- (4,1);
    \draw[dashed] (1,2) -- (5,3);
    \draw[dashed] (0,0) -- (1,2);
    \draw[dashed] (4,1) -- (5,3);
    
\end{tikzpicture}
\end{center}

\end{multicols}







{\bf Solución:} Sean los puntos dados:
\[ A = (2,-2), \quad B = (6,-1), \quad C = (9,3) \]

\begin{enumerate}[a) ]
    \item {\bf Ecuaciones de las rectas que contienen los lados del paralelogramo ${\bf ABCD}$:}










{\bf Recordar que } la ecuación de la recta que pasa por dos puntos $(x_1, y_1)$ y $(x_2, y_2)$ se obtiene con la ecuación:
\[
y =  \frac{y_2 - y_1}{x_2 - x_1} (x - x_1) + y_1 
\]

Calculamos la pendiente y ecuación de cada recta:

\begin{itemize}
\begin{multicols}{2}
    \item  {\bf Recta $AB$:}

  \begin{align*}
      y &=  \frac{-1 - (-2)}{6 - 2} (x - 2) -2 
      \\
      y  &=  \frac{1}{4} (x - 2) - 2 
      \\
      y  &=  \frac{1}{4} x - \frac{1}{4} \cdot 2  -2 
      \\
      y &= \tfrac{1}{4}x - \tfrac{5}{2}.
  \end{align*} 

    \item {\bf Recta $BC$:} 
  \begin{align*}
     y  &=  \frac{3 - (-1)}{9 - 6} (x - 6) -1
     \\
     y  &=  \frac{4}{3} (x - 6) -1
     \\
     y  &=  \frac{4}{3} x -  \frac{4}{3} \cdot 6 -1
     \\
     y &= \tfrac{4}{3}x - 9.
  \end{align*}
\end{multicols}
\end{itemize}


\setlength{\columnsep}{9mm}
\begin{multicols}{2}

Por ahora tenemos dos rectas y son las que aparecen en el siguiente gráfico: 

\vspace{-6mm}
\begin{center}
\definecolor{gris}{rgb}{0.5,0.5,0.5}
\begin{tikzpicture}[scale=0.5]
% extremos para los ejes
\xdef\ext{10}
\xdef\yext{4}
\xdef\f{-1}
\xdef\g{-4}
\pgfmathtruncatemacro{\c}{\f+1}
\pgfmathtruncatemacro{\d}{\g+1}
\draw [color=gris!30!,dash pattern=on 1pt off 1pt, xstep=1cm,ystep=1cm] (\f-.3,\g-.3) grid (\ext+.3,\yext+.3);

\draw[->] (\f-.3,0) -- (\ext+.3,0) ;

\draw[->] (0,\g-.3) -- (0,\yext+.3) ;

\foreach \x in {\f,...,-1,1,2,...,\ext}
\draw[shift={(\x,0)},color=black] (0pt,2pt) -- (0pt,-2pt) node[below] {\tiny $\x$};

\foreach \y in {\g,...,-1,1,2,...,\yext }
\draw[shift={(0,\y)},color=black] (2pt,0pt) -- (-2pt,0pt) node[left] {\tiny $\y$};
    
    
    
    % Puntos
    \coordinate (A) at (2,-2);%punto a la izquierda
\coordinate (B) at (6,-1);%punto a la derecha
\coordinate (C) at (9,3);%Punto por detras de A
\coordinate (D) at ($(A)+(C)-(B)$);%Punto por delante de B




%\draw[blue,fill=blue] (D) circle (3pt)node[below]{$D$} ;

    % Lados del paralelogramo
    \draw[thick, line width=3pt,orange,opacity=.5] (A) -- (B) -- (C);

    % Rectas extendidas
    \coordinate (M) at ($1.9*(A) -.9*(B)$);%Punto por detras de A
    \coordinate (m) at ($-1.2*(A)+2.2*(B)$);%Punto por delante de B
     \draw[thick, line width=2pt, dash pattern=on 2pt off 2pt, opacity=.6] (M) -- (m) node[opacity=.9,right] {$y = \tfrac{1}{4}x - \tfrac{5}{2}.$}; % AB 
    
    \coordinate (L) at ($1.9*(B) -.9*(C)$);%Punto por detras de B
    \coordinate (l) at ($-.3*(B)+1.3*(C)$);%Punto por delante de C
     \draw[thick, line width=2pt, dash pattern=on 2pt off 2pt, opacity=.6] (L) -- (l)  node[opacity=.9,right] {$y = \tfrac{4}{3}x - 9$}; %    
   

%puntos
    \draw[blue,fill=blue] (A) circle (3pt) node[below]{$A$} ;

\draw[blue,fill=blue] (B) circle (3pt)node[below]{$B$} ;

\draw[blue,fill=blue] (C) circle (3pt)node[below]{$C$} ;

\end{tikzpicture}
\end{center}


Como en el paralelogramo los lados opuestos son \textbf{paralelos}, 
Las otras dos rectas que faltan son las siguientes 

\begin{itemize}
    \item La recta que pasa por $C$ y es paralela a $\overleftrightarrow{AB} $.

    \item La recta que pasa por $A$ y es paralela a $\overleftrightarrow{BC}$.
\end{itemize}






\end{multicols}



Para resolver usaremos la fórmula de la recta que pasa por los puntos $(x_0,y_0)$ y tiene pendiente $M$: \[y = m \,  (x -x_0 ) + y_0\]


\begin{itemize}
\setlength{\columnsep}{9mm}
\begin{multicols}{2}
    \item La recta que pasa por el punto $ C = (9,3) $ y es paralela a la recta $ \overleftrightarrow{AB}$ tiene la misma pendiente que esta, es decir, su pendiente es $ \tfrac{1}{4} $.

    \begin{align*}
        y &= \tfrac{1}{4}  \,  (x -9 ) + 3 
        \\[2mm]
        y &= \tfrac{1}{4} x - \tfrac{1}{4} \cdot 9  + 3 
        \\[2mm]
        y &= \tfrac{1}{4} x +\tfrac{3}{4}  
    \end{align*}

    \item La recta que pasa por $A=(2,-2)$ y es paralela a la recta  $\overleftrightarrow{BC}$ tiene la misma pendiente que esta, es decir, su pendiente es $\tfrac{4}{3}$.

    \begin{align*}
        y &= \tfrac{4}{3}  \,  (x -2 ) -2
        \\[2mm]
        y &= \tfrac{4}{3} x - \tfrac{4}{3} \cdot 2  -2 
        \\[2mm]
        y &= \tfrac{4}{3} x -\tfrac{14}{3}  
    \end{align*}
\end{multicols}
\end{itemize}

\begin{center}
\definecolor{gris}{rgb}{0.5,0.5,0.5}
\begin{tikzpicture}[scale=0.5]
% extremos para los ejes
\xdef\ext{10}
\xdef\yext{4}
\xdef\f{-1}
\xdef\g{-4}
\pgfmathtruncatemacro{\c}{\f+1}
\pgfmathtruncatemacro{\d}{\g+1}
\draw [color=gris!30!,dash pattern=on 1pt off 1pt, xstep=1cm,ystep=1cm] (\f-.3,\g-.3) grid (\ext+.3,\yext+.3);

\draw[->] (\f-.3,0) -- (\ext+.3,0) ;

\draw[->] (0,\g-.3) -- (0,\yext+.3) ;

\foreach \x in {\f,...,-1,1,2,...,\ext}
\draw[shift={(\x,0)},color=black] (0pt,2pt) -- (0pt,-2pt) node[below] {\tiny $\x$};

\foreach \y in {\g,...,-1,1,2,...,\yext }
\draw[shift={(0,\y)},color=black] (2pt,0pt) -- (-2pt,0pt) node[left] {\tiny $\y$};
    
    
    
    % Puntos
    \coordinate (A) at (2,-2);%punto a la izquierda
\coordinate (B) at (6,-1);%punto a la derecha
\coordinate (C) at (9,3);%Punto por detras de A
\coordinate (D) at ($(A)+(C)-(B)$);%Punto por delante de B




%\draw[blue,fill=blue] (D) circle (3pt)node[below]{$D$} ;

    % Lados del paralelogramo
    \draw[thick, line width=3pt,orange,opacity=.5] (A) -- (B) -- (C)--(D)--cycle;

    % Rectas extendidas
    \coordinate (M) at ($1.9*(A) -.9*(B)$);%Punto por detras de A
    \coordinate (m) at ($-1.2*(A)+2.2*(B)$);%Punto por delante de B
     \draw[thick, line width=2pt, dash pattern=on 2pt off 2pt, opacity=.6] (M) -- (m) node[opacity=.9,right] {$y = \tfrac{1}{4}x - \tfrac{5}{2}.$}; % AB 
    
    \coordinate (L) at ($1.9*(B) -.9*(C)$);%Punto por detras de B
    \coordinate (l) at ($-.3*(B)+1.3*(C)$);%Punto por delante de C
     \draw[thick, line width=2pt, dash pattern=on 2pt off 2pt, opacity=.6] (L) -- (l)  node[opacity=.9, above right] {$y = \tfrac{4}{3}x - 9$}; %   

    \coordinate (M) at ($2.9*(D) -1.9*(C)$);%Punto por detras de A
    \coordinate (m) at ($-.3*(D)+1.3*(C)$);%Punto por delante de B
     \draw[thick, line width=2pt, dash pattern=on 2pt off 2pt, opacity=.6] (M) -- (m) node[opacity=.9,right] {$y = \tfrac{1}{4}x +\tfrac{3}{4}.$}; 

     \coordinate (L) at ($1.55*(A) -.55*(D)$);%Punto por detras de B
    \coordinate (l) at ($-.6*(A)+1.6*(D)$);%Punto por delante de C
     \draw[thick, line width=2pt, dash pattern=on 2pt off 2pt, opacity=.6] (L) -- (l)  node[opacity=.9, above ] {$y = \tfrac{4}{3}x - \tfrac{14}{3}$}; % 
   

%puntos
    \draw[blue,fill=blue] (A) circle (3pt) node[below]{$A$} ;

\draw[blue,fill=blue] (B) circle (3pt)node[below]{$B$} ;

\draw[blue,fill=blue] (C) circle (3pt)node[below right]{$C$} ;

%\draw[blue,fill=blue] (D) circle (3pt)node[above left]{$D$} ;

\end{tikzpicture}
\end{center}



\begin{itemize}
\begin{minipage}{.6\textwidth}
    

    \item 
Para hallar $D$, usamos que en un paralelogramo $D = A + \overrightarrow{BC}$:
  \[ D = (2,-2) + ((9,3) - (6,-1)) = (2 + 3, -2 + 4) = (5,2). \]
\, \\
\, \\
\, \\
  \end{minipage}\qquad\begin{minipage}{.3\textwidth}
    
\begin{center}
\definecolor{gris}{rgb}{0.5,0.5,0.5}
\begin{tikzpicture}[scale=0.4]
% extremos para los ejes
\xdef\ext{10}
\xdef\yext{4}
\xdef\f{-1}
\xdef\g{-4}
\pgfmathtruncatemacro{\c}{\f+1}
\pgfmathtruncatemacro{\d}{\g+1}
\draw [color=gris!30!,dash pattern=on 1pt off 1pt, xstep=1cm,ystep=1cm] (\f-.3,\g-.3) grid (\ext+.3,\yext+.3);

\draw[->] (\f-.3,0) -- (\ext+.3,0) ;

\draw[->] (0,\g-.3) -- (0,\yext+.3) ;

\foreach \x in {\f,...,-1,1,2,...,\ext}
\draw[shift={(\x,0)},color=black] (0pt,2pt) -- (0pt,-2pt) node[below] {\tiny $\x$};

\foreach \y in {\g,...,-1,1,2,...,\yext }
\draw[shift={(0,\y)},color=black] (2pt,0pt) -- (-2pt,0pt) node[left] {\tiny $\y$};
    
    
    
    % Puntos
    \coordinate (A) at (2,-2);%punto a la izquierda
\coordinate (B) at (6,-1);%punto a la derecha
\coordinate (C) at (9,3);%Punto por detras de A
\coordinate (D) at ($(A)+(C)-(B)$);%Punto por delante de B

  
    % Lados del paralelogramo
    \draw[thick, line width=1pt,orange,opacity=.5] (A) -- (B) -- (C)--(D)--cycle;
    
% Diagonal AC
% Rectas extendidas
    
    % vector director
    \draw[->, thick, purple,line width=2pt,>=triangle 45] (B) -- (C) node[midway, right ] {$\overrightarrow{BC}$};
    \draw[->, thick, purple,line width=2pt,>=triangle 45] (A) -- (D);

\draw[blue,fill=violet , color=blue] (A) circle (2pt) node[below left]{$A$} ;

\draw[blue,fill=blue] (B) circle (2pt)node[below]{$B$} ;

\draw[blue,fill=blue] (C) circle (2pt)node[above left]{$C$} ;

\draw[blue,fill=blue] (D) circle (2pt)node[above]{$D$} ;
\end{tikzpicture}
\end{center}

\end{minipage}
\end{itemize}


\item {\bf  Ecuación vectorial de la diagonal $AC$}

La ecuación vectorial de la recta que pasa por $A$ en la dirección de $\overrightarrow{AC}$ es:
\[ L:\ \  (x,y) = (2,-2) + \lambda(7,5), \quad \lambda \in \mathbb{R}. \]


 {\bf Gráfico}

\begin{center}
\definecolor{gris}{rgb}{0.5,0.5,0.5}
\begin{tikzpicture}[scale=0.5]
% extremos para los ejes
\xdef\ext{10}
\xdef\yext{4}
\xdef\f{-1}
\xdef\g{-4}
\pgfmathtruncatemacro{\c}{\f+1}
\pgfmathtruncatemacro{\d}{\g+1}
\draw [color=gris!30!,dash pattern=on 1pt off 1pt, xstep=1cm,ystep=1cm] (\f-.3,\g-.3) grid (\ext+.3,\yext+.3);

\draw[->] (\f-.3,0) -- (\ext+.3,0) ;

\draw[->] (0,\g-.3) -- (0,\yext+.3) ;

\foreach \x in {\f,...,-1,1,2,...,\ext}
\draw[shift={(\x,0)},color=black] (0pt,2pt) -- (0pt,-2pt) node[below] {\tiny $\x$};

\foreach \y in {\g,...,-1,1,2,...,\yext }
\draw[shift={(0,\y)},color=black] (2pt,0pt) -- (-2pt,0pt) node[left] {\tiny $\y$};
    
    
    
    % Puntos
    \coordinate (A) at (2,-2);%punto a la izquierda
\coordinate (B) at (6,-1);%punto a la derecha
\coordinate (C) at (9,3);%Punto por detras de A
\coordinate (D) at ($(A)+(C)-(B)$);%Punto por delante de B

  
    % Lados del paralelogramo
    \draw[thick, line width=1pt,orange,opacity=.5] (A) -- (B) -- (C)--(D)--cycle;
    
% Diagonal AC
% Rectas extendidas
    \coordinate (M) at ($1.47*(A) -.47*(C)$);%Punto por detras de A
    \coordinate (m) at ($-.2*(A)+1.2*(C)$);%Punto por delante de B
     \draw[thick, line width=1pt, red] (M) -- (m) node[opacity=.9,right] {$L :\ (x,y) = (2,-2) + \lambda(7,5), \quad \lambda \in \mathbb{R}. $}; % AB 
    \draw[thick, red] (A) -- (C);
    

\draw[blue,fill=violet , color=blue] (A) circle (2pt) node[below]{$A$} ;

\draw[blue,fill=blue] (B) circle (2pt)node[below]{$B$} ;

\draw[blue,fill=blue] (C) circle (2pt)node[below]{$C$} ;

\draw[blue,fill=blue] (D) circle (2pt)node[above]{$D$} ;
\end{tikzpicture}
\end{center}
 




\end{enumerate}











\item Sean los puntos $A = (-1,-2)$ y $B = (3,6)$:

\begin{enumerate}[(a)]
    \item Encontrar la ecuación de la recta mediatriz del segmento $AB$.
    \item  Hallar un punto sobre el eje $x$ que equidiste de los puntos $A$ y $B$.
    \item  Hallar, si es posible, un punto sobre la recta $L : 2x - y + 1 = 0$ que equidiste de los puntos $A$ y $B$.
    \item Hallar, si es posible, un punto sobre la recta $L : 2x + 4y + 4 = 0$ que equidiste de los puntos $A$ y $B$.
\end{enumerate}



{\bf Recordemos:}

\begin{itemize}
    \item {\bf Definición de Mediatriz:}  La {\bf mediatriz} de un segmento es la recta perpendicular al segmento en su punto medio. Es decir, si tenemos un segmento $ AB $, la mediatriz es la recta que pasa por el punto medio de $ AB $ y es perpendicular a este segmento.


\item La principal propiedad de la recta mediatriz a un segmento $\overline{AB}$ es que sus puntos equidistan a los extremos $A$ y $B$.


\item {\bf punto medio} de un segmento cuyos extremos son los puntos $ A=(x_1, y_1) $ y $ B=(x_2, y_2) $ se calcula con la siguiente fórmula:

\[
M= \dfrac{A+B}{2} = \left( \frac{x_1 + x_2}{2}, \frac{y_1 + y_2}{2} \right)
\]

\item Si una recta tiene pendiente $ m $, entonces la pendiente de la recta {\bf perpendicular} es:

\[
m_{\perp} = -\frac{1}{m}
\]
 
\end{itemize}



{\bf Solución:}

\begin{enumerate}[(a)]
    \item {\bf Ecuación de la recta mediatriz del segmento $\bf AB$}
El punto medio del segmento $AB$ es:
\begin{equation*}
    M= \dfrac{\ (-1,2) +(3,6)\ }{2} = \dfrac{\ (-1+3,2+6) \ }{2} = \dfrac{\ (2,8) \ }{2}   = (1,2)
\end{equation*}

\begin{itemize} 
\setlength{\columnsep}{9mm}
\begin{multicols}{2}
\item La pendiente de $AB$ es:
\begin{equation*}
    m_{AB} = \frac{6 - (-2)}{3 - (-1)} = \frac{8}{4} = 2
\end{equation*}

\item La mediatriz es perpendicular a $AB$, por lo que su pendiente es:
\begin{equation*}
    m = -\frac{1}{2}
\end{equation*}
\end{multicols}
 
\end{itemize}

Tenemos que la recta mediatriz  pasa por $M= (1,2)$ y tiene pendiente $-\tfrac{1}{2}$,  utilizando la ecuación punto-pendiente $Y= m (x-x_0) + y_0$ tenemos:
\[
    y  = -\frac{1}{2} (x - 1) +2 
\qquad 
\stackrel[\text{Despejando:}]{}{\Longleftrightarrow}
\qquad   
 y = -\frac{1}{2}x + \frac{5}{2}
\]



\begin{center}
\definecolor{gris}{rgb}{0.5,0.5,0.5}
\begin{tikzpicture}[scale=0.5]
% extremos para los ejes
\xdef\ext{5}
\xdef\yext{7}
\xdef\f{-2}
\xdef\g{-3}
\pgfmathtruncatemacro{\c}{\f+1}
\pgfmathtruncatemacro{\d}{\g+1}
\draw [color=gris!30!,dash pattern=on 1pt off 1pt, xstep=1cm,ystep=1cm] (\f-.3,\g-.3) grid (\ext+.3,\yext+.3);

\draw[->] (\f-.3,0) -- (\ext+.3,0) ;

\draw[->] (0,\g-.3) -- (0,\yext+.3) ;

\foreach \x in {\f,...,-1,1,2,...,\ext}
\draw[shift={(\x,0)},color=black] (0pt,2pt) -- (0pt,-2pt) node[below] {\tiny $\x$};

\foreach \y in {\g,...,-1,1,2,...,\yext }
\draw[shift={(0,\y)},color=black] (2pt,0pt) -- (-2pt,0pt) node[left] {\tiny $\y$};
    
    
    \draw[thick,blue] (-1,-2) -- (3,6) node[above] {$AB$}; 
    \draw[thick,red, dashed,domain=-2.33:5.3] plot ({\x}, {-.5*\x+5/2}) node[above right] {Mediatriz};
    \node[red,right] at (5.3,-.3) {$y = -\tfrac{1}{2}x + \tfrac{5}{2}$};
    \filldraw[black] (-1,-2) circle (2pt) node[left] {$A$};
    \filldraw[black] (3,6) circle (2pt) node[right] {$B$};
    \filldraw[black] (1,2) circle (2pt) node[above] {$M$};
\end{tikzpicture} 

\end{center}

\item {\bf Punto sobre el eje $x$ equidistante de $A$ y $B$}
Buscamos $P = (x,0)$ tal que $PA = PB$:
\begin{align*}
    d\left(\rule{0mm}{4mm} (x,0),\, (-1,-2) \right) &= 
    d\left(\rule{0mm}{4mm}  (x,0),\, (3,6) \right)
    \\
    \sqrt{(x + 1)^2 + (0 + 2)^2} &= \sqrt{(x - 3)^2 + (0 - 6)^2}
    \\
    \sqrt{(x + 1)^2 + 4} &= \sqrt{(x - 3)^2 + 36}
    \intertext{Elevando al cuadrado}
   (x + 1)^2 +  4 &=  (x - 3)^2 + 36
   \\
   x^2+2x +1 +  4 &=  x^2 -6x +9+ 36
   \\
  \cancel{ x^2}+2x +1 +  4 &= \cancel{ x^2} -6x +9+ 36
   \\
  2x +6x &=  9+ 36 -1-4
  \\
  8x &= 40
  \\
  x &= \frac{40}{8}
  \\
  x &= 5
\end{align*}
\hfill {\bf Rta:} el punto es: $P = (5,0)$.


\medskip 

\item {\bf Punto sobre $\bf 2x - y + 1 = 0$ equidistante de $A$ y $B$}:
Un punto equidistante de $A$ y $B$ está en la mediatriz. Verificamos si la intersección de la mediatriz con $2x - y + 1 = 0$ es posible.

\bigskip

$
\left\{
\begin{array}{rl}
y= -\frac{1}{2}x + &\!  \frac{5}{2}   \marka{a}{(.1,.25)}
\\[5mm]
 2x - y + 1  &= 0 \indicador{b}{(.1,.15)}{c}{(.6,.45)} \marka{d}{(1.3,.45)}
\end{array}
\right.
$


\tikz[overlay,remember picture] {
\draw [overlay,->,very thick,red,opacity=.5](b)-|(c)--(d);  
\draw [overlay,very thick,red,opacity=.5](a)++(0,-.1)-|(c);
\node  at (d) [text width=55mm, below right , yshift= 7mm ]{
{\scriptsize Sustituyendo:}


\vspace{-8mm}
\begin{align*}
     2x - \left(-\tfrac{1}{2}x + \tfrac{5}{2}\right) + 1 &= 0  
     \\
     2x + \tfrac{1}{2}x - \tfrac{5}{2} + 1 &= 0  
     \\ 
     \tfrac{5}{2}x - \tfrac{3}{2} &= 0  
     \\ 
     \tfrac{5}{2}x \ \  &= \tfrac{3}{2}
     \\
     x \ \  &= \tfrac{3}{2} : \tfrac{5}{2}
     \\[2mm]
   \marka{a1}{(-.1,-.2)} 
   x &= \tfrac{3}{5} .
   \marka{a2}{(.1,.45)} \indicador{b}{(.13,.15)}{c}{(.9,1.95)} \marka{d}{(1.7,1.95)}
\end{align*}
};
\draw[color=black!50!,rounded corners=1mm] (a1) rectangle (a2);
%%%%%%%%%%%%%%%%%%%%%%%%%%%%%%%%%%%%%%%%%%%%%%%%%%%%%%%%%%%%%%%%%%%%%%%%%%%%%%%%%%%%%%%%%%%%%%%%%%%%%%%%%%%%%%%%%%%%%%%%%%%%%%%%%%%%%%%%%%%%%%%%%%%%%%%%%%%%%%%%%%%%%%%%%%%%%%%%%%%%
\draw [overlay,->,very thick,blue,opacity=.5](b)-|(c)--(d);  
\draw [overlay,very thick,blue,opacity=.5](a)-|(c);
\node  at (d) [text width=60mm, below right , yshift= 9mm ]{
{\scriptsize Sustituyendo $ x= \tfrac{3}{5}$ en la primera ecuación:}

\vspace{-8mm}
\begin{align*}
   y&= -\tfrac{1}{2} \cdot \tfrac{3}{5} +  \tfrac{5}{2}  \qquad\qquad \qquad\\
    y &= - \tfrac{3}{10} +  \tfrac{5}{2}, \\[2mm] 
     \marka{a1}{(-.1,-.2)}
    y &= \tfrac{11}{5}.
   \marka{a2}{(.1,.45)} 
\end{align*}
};
\draw[color=black!50!,rounded corners=1mm] (a1) rectangle (a2);
}
%%%%%%%%%%%%%%%%%%%%%%%%%%%%%%%%%%%%%%%%%%%%%%%%%%%%%%%%%%%%%%%%%%%%%%%%%%%%%%%%%%%%%%%%%%%%%%%%%%%%%%%%%%%%%%%%%%%%%%%%%%%%%%%%%%%%%%%%%%%%%%%%%%%%%%%%%%%%%%%%%%%%%%%%%%%%%%%%%%%%




\vspace{29mm}







Resolviendo, obtenemos $x = \frac{1}{5}$ y $y = \frac{12}{5}$. Luego: \quad $\displaystyle  P = \left(\frac{1}{5}, \frac{12}{5}\right) $.  


\bigskip

\item {\bf Punto sobre $\bf 2x + 4y + 4 = 0$ equidistante de $A$ y $B$:}
 
 

$
\left\{
\begin{array}{rl}
y= -\frac{1}{2}x + &\!  \frac{5}{2}   \marka{a}{(.1,.25)}
\\[5mm]
 2x + 4y + 4 &= 0 \indicador{b}{(.1,.15)}{c}{(.6,.65)} \marka{d}{(1.3,.65)}
\end{array}
\right.
$


\tikz[overlay,remember picture] {
\draw [overlay,->,very thick,red,opacity=.5](b)-|(c)--(d);  
\draw [overlay,very thick,red,opacity=.5](a)++(0,-.1)-|(c);
\node  at (d) [text width=55mm, below right , yshift= 7mm ]{
{\scriptsize Sustituyendo:}


\vspace{-8mm}
\begin{align*}
     2x +4  \left(-\tfrac{1}{2}x + \tfrac{5}{2}\right) + 4 &= 0  
     \\
     2x - 2x + 10 + 4 &= 0  
     \\ 
     14 &= 0   \quad \text{ Absurdo!}
\end{align*}
};
}
%%%%%%%%%%%%%%%%%%%%%%%%%%%%%%%%%%%%%%%%%%%%%%%%%%%%%%%%%%%%%%%%%%%%%%%%%%%%%%%%%%%%%%%%%%%%%%%%%%%%%%%%%%%%%%%%%%%%%%%%%%%%%%%%%%%%%%%%%%%%%%%%%%%%%%%%%%%%%%%%%%%%%%%%%%%%%%%%%%%%

\vspace{9mm}
 Por lo tanto, no existe un punto sobre esta recta que equidiste de $A$ y $B$.
 

\end{enumerate}












  \item  Hallar, si es posible, un punto del eje $y$ que equidiste de la recta $L : 2x + y - 7 = 0$ y del punto $P = (-2,2)$.


\textbf{Solución:}

Sea el punto buscado $Q = (0, y_0)$, ya que está sobre el eje $y$. La condición de equidistancia implica que la distancia de $Q$ a la recta $L$ es igual a la distancia de $Q$ al punto $P$.

\textbf{Paso 1: Distancia del punto $Q$ a la recta $L$}
La distancia de un punto $(x_0, y_0)$ a la recta $Ax + By + C = 0$ está dada por la fórmula:
\[
    d = \frac{|Ax_0 + By_0 + C|}{\sqrt{A^2 + B^2}}.
\]

Sustituyendo los valores de la recta $L: 2x + y - 7 = 0$ y el punto $Q = (0, y_0)$:
\[
    d_1 = \frac{|2(0) + y_0 - 7|}{\sqrt{2^2 + 1^2}} = \frac{|y_0 - 7|}{\sqrt{5}}.
\]

\textbf{Paso 2: Distancia del punto $Q$ al punto $P$}
La distancia entre dos puntos $(x_1, y_1)$ y $(x_2, y_2)$ es:
\[
    d = \sqrt{(x_2 - x_1)^2 + (y_2 - y_1)^2}.
\]

Sustituyendo $P = (-2,2)$ y $Q = (0, y_0)$: \quad $\displaystyle d = \sqrt{(0 + 2)^2 + (y_0 - 2)^2} = \sqrt{4 + (y_0 - 2)^2}. $ 

\textbf{Paso 3: Igualando ambas distancias y despejamos:}

\begin{align*}
    \frac{|y_0 - 7|}{\sqrt{5}} &= \sqrt{4 + (y_0 - 2)^2} &\ \hspace{.1\textwidth}
    \\ 
    \left(\frac{|y_0 - 7|}{\sqrt{5}}\right)^2 &= \left(\sqrt{4 + (y_0 - 2)^2} \right)^2
    \\
     \frac{(y_0 - 7)^2}{ 5 } \ \   &=  \ \ \ 4 + (y_0 - 2 )`^2
    \\ 
    (y_0 - 7)^2  \ \ \ &= \ \ \ 5\cdot \left( 4 + (y_0 - 2)^2  \rule{0mm}{4mm} \right)
    \\ 
    (y_0 - 7)^2 \quad &= \quad 20 + 5 \, (y_0 - 2)^2  
    \\[3mm]
    \text{Expandimos ambos lados de la igualdad:}\hspace{34mm}
     y_0^2 - 14y_0 + 49 \quad &= \quad  20 + 5(y_0^2 - 4y_0 + 4)
     \\[3mm]
     y_0^2 - 14y_0 + 49 \quad &= \quad  20 + 5y_0^2 - 20y_0 + 20
     \\
     y_0^2 - 14y_0 + 49 - 5y_0^2 + 20y_0 - 40 \quad &= \quad  0
     \\
      -4y_0^2 + 6y_0 + 9 \quad &= \quad  0
      \\
       4y_0^2 - 6y_0 - 9 \quad &= \quad  0
\end{align*}

Resolvemos con la fórmula cuadrática:
\[
    y_0 = \frac{-(-6) \pm \sqrt{(-6)^2 - 4(4)(-9)}}{2(4)} = \frac{6 \pm \sqrt{36 + 144}}{8} = \frac{6 \pm \sqrt{180}}{8}.
\]
\[
    y_0 = \frac{6 \pm 6\sqrt{5}}{8} = \frac{3 \pm 3\sqrt{5}}{4}.
\]

\textbf{Rta:} Los puntos del eje $y$ que cumplen la condición de equidistancia son:
\[
    Q_1 = \left(0, \frac{3 + 3\sqrt{5}}{4}\right), \quad Q_2 = \left(0, \frac{3 - 3\sqrt{5}}{4}\right).
\]

\textbf{Gráfico}


 \begin{center}
     
\definecolor{rojo}{rgb}{0.7,0.4,0.4}
\definecolor{azul}{rgb}{.4,.4,0.9}
 
\definecolor{gris}{rgb}{0.5,0.5,0.5}
\begin{tikzpicture}[scale=0.7]
% extremos para los ejes
\xdef\ext{4}
\xdef\yext{7}
\xdef\f{-4}
\xdef\g{-2}
\pgfmathtruncatemacro{\c}{\f+1}
\pgfmathtruncatemacro{\d}{\g+1}
\draw [color=gris!30!,dash pattern=on 1pt off 1pt, xstep=1cm,ystep=1cm] (\f-.3,\g-.3) grid (\ext+.3,\yext+.3);

\draw[->] (\f-.3,0) -- (\ext+.3,0) ;

\draw[->] (0,\g-.3) -- (0,\yext+.3) ;

\foreach \x in {\f,...,-1,1,2,...,\ext}
\draw[shift={(\x,0)},color=black] (0pt,2pt) -- (0pt,-2pt) node[below] {\tiny $\x$};

\foreach \y in {\g,...,-1,1,2,...,\yext }
\draw[shift={(0,\y)},color=black] (2pt,0pt) -- (-2pt,0pt) node[left] {\tiny $\y$};


    \coordinate (P) at (-2,2);%punto 
%%%% Recta que pasa por A y B
\filldraw[red] (P) circle (2pt) node[left] {$P=(-2,2)$};
 
% la recta
    \draw[thick,blue,domain=-.3:4.3] plot ({\x}, {7-2*\x}) node[above right] {$2x+y-7=0$};

% Puntos equidistantes en el eje y
        \pgfmathsetmacro{\yy}{(3 + 3*sqrt(5))/4}
        \pgfmathsetmacro{\yyy}{(3 - 3*sqrt(5))/4}


        \filldraw[teal] (0,\yy) circle (2pt) node[right] {$Q_1$};
        \filldraw[violet] (0,\yyy) circle (2pt) node[right] {$Q_2$};

          % Circunferencias
          \begin{scope}
        \clip (-2.1,1.2) --(3,2.5)--(3,5)-- (-2.1,5)--cycle;
        \draw[dashed] (0,\yy) circle (2.045);
        \end{scope}
        \begin{scope}
         \clip (-2.4,-.9) rectangle (4,5);
        \draw[dashed]  (0,\yyy)  circle (3.545);     
        \end{scope}

        \draw[thick, teal,opacity=.2] (P) -- (0,\yy);
        \draw[thick, violet,opacity=.2] (P) -- (0,\yyy);

%otros radios a ojimetro
        \draw[thick, teal,opacity=.2] (1.9,3.2) -- (0,\yy);
        \draw[thick, violet,opacity=.2] (3.2,0.6) -- (0,\yyy);
        
\end{tikzpicture}

 \end{center}













   
\item  Calcular el ángulo que forman los siguientes pares de rectas:

\begin{enumerate}[(a) ]
\begin{multicols}{2}
    \item $\begin{cases}
    L_1 : 3x - y + 2 = 0 \\[6mm]
    L_2 : 2x + y = 2
    \end{cases} $


    \item $\left\{
    \begin{array}{l}
        L_1 : \begin{cases}
        x = 3 + \lambda \\
        y = 2 + 2\lambda
        \end{cases},\quad \lambda \in \mathbb{R},
        \\[4mm]
        L_2 : -2x + y + 4 = 0.
    \end{array}\right.$
    
\end{multicols}
\end{enumerate}


\textbf{Solución:}



\textbf{  (a)}

\begin{multicols}{2}
\begin{center}
    

\definecolor{gris}{rgb}{0.5,0.5,0.5}
\begin{tikzpicture}[scale=0.5]
% extremos para los ejes
\xdef\ext{4}
\xdef\yext{10}
\xdef\f{-4}
\xdef\g{-1}
\pgfmathtruncatemacro{\c}{\f+1}
\pgfmathtruncatemacro{\d}{\g+1}
\draw [color=gris!30!,dash pattern=on 1pt off 1pt, xstep=1cm,ystep=1cm] (\f-.3,\g-.3) grid (\ext+.3,\yext+.3);

\draw[->] (\f-.3,0) -- (\ext+.3,0) ;

\draw[->] (0,\g-.3) -- (0,\yext+.3) ;

\foreach \x in {\f,...,-1,1,2,...,\ext}
\draw[shift={(\x,0)},color=black] (0pt,2pt) -- (0pt,-2pt) node[below] {\tiny $\x$};

\foreach \y in {\g,...,-1,1,2,...,\yext }
\draw[shift={(0,\y)},color=black] (2pt,0pt) -- (-2pt,0pt) node[left] {\tiny $\y$};


    \draw[thick,blue,domain=-3.6:.6] plot ({\x}, {3*\x+9}) node[above right] {$y-3x-9=0$};
    \draw[thick,blue,domain=-4.33:1.5] plot ({\x}, {-2*\x+2}) node[below right] {$L_2 : \ \ 2x + y = 2$};
    
\end{tikzpicture}

\end{center}



Los vectores normales de las rectas son:
\begin{align*}
    \vec{n_1} &= (3, -1), \\
    \vec{n_2} &= (2, 1).
\end{align*}


El ángulo $ \theta $ entre las rectas se obtiene mediante la fórmula:
\begin{equation*}
    \cos \theta = \frac{\vec{n_1} \cdot \vec{n_2}}{|\vec{n_1}| |\vec{n_2}|}.
\end{equation*}

\end{multicols}



\begin{multicols}{2}

Calculamos el producto escalar:
\begin{equation*}
    \vec{n_1} \cdot \vec{n_2} = (3)(2) + (-1)(1) = 6 - 1 = 5.
\end{equation*}

Las normas de los vectores son:
\begin{align*}
    |\vec{n_1}| &= \sqrt{3^2 + (-1)^2} = \sqrt{10}, \\
    |\vec{n_2}| &= \sqrt{2^2 + 1^2} = \sqrt{5}.
\end{align*}

Entonces,
\begin{equation*}
    \cos \theta = \frac{5}{\sqrt{10} \cdot \sqrt{5}} = \frac{5}{\sqrt{50}} = \frac{5}{5\sqrt{2}} = \frac{1}{\sqrt{2}}.
\end{equation*}

Por lo tanto,
\begin{equation*}
    \theta = \arccos \left( \frac{1}{\sqrt{2}} \right) = 45^\circ.
\end{equation*}

\end{multicols}








\textbf{ (b)}

\begin{multicols}{2}
\begin{center}
    

\definecolor{gris}{rgb}{0.5,0.5,0.5}
\begin{tikzpicture}[scale=0.5]
% extremos para los ejes
\xdef\ext{4}
\xdef\yext{5}
\xdef\f{-4}
\xdef\g{-1}
\pgfmathtruncatemacro{\c}{\f+1}
\pgfmathtruncatemacro{\d}{\g+1}
\draw [color=gris!30!,dash pattern=on 1pt off 1pt, xstep=1cm,ystep=1cm] (\f-.3,\g-.3) grid (\ext+.3,\yext+.3);

\draw[->] (\f-.3,0) -- (\ext+.3,0) ;

\draw[->] (0,\g-.3) -- (0,\yext+.3) ;

\foreach \x in {\f,...,-1,1,2,...,\ext}
\draw[shift={(\x,0)},color=black] (0pt,2pt) -- (0pt,-2pt) node[below] {\tiny $\x$};

\foreach \y in {\g,...,-1,1,2,...,\yext }
\draw[shift={(0,\y)},color=black] (2pt,0pt) -- (-2pt,0pt) node[left] {\tiny $\y$};


    \draw[thick,blue,domain=-1.6:1.6,variable=\t,opacity=.5] plot ({3+\t}, {2+2*\t}) node[above right] {$ L_1 : \begin{cases}
        x = 3 + \lambda \\
        y = 2 + 2\lambda
        \end{cases}$};
    \draw[thick,teal,domain=1.53:4.3,opacity=.5] plot ({\x}, {2*\x-4}) node[below right] {$L_2 : \ \ -2x + y + 4 = 0.$};
    
\end{tikzpicture}

\end{center}




El vector director de $ L_1 $ es:
\begin{equation*}
    \vec{d_1} = (1, 2).
\end{equation*}

El vector normal de $L_2$ es:
\begin{equation*}
    \vec{n_2} = (-2, 1).
\end{equation*} 

El vector director de $ L_1 $ es:
\begin{equation*}
    \vec{d_1} = (1, 2).
\end{equation*}


\end{multicols}
Como las dos rectas tienen el mismo director, el ángulo es cero. Es más, en este caso, las rectas son coincidentes.





\item Dadas las rectas $L_1 : 3x + 5 = 0$ y $L_2 : ax + 2y - 1 = 0$. Determinar el valor de $a \in \mathbb{R}$ de modo que formen un ángulo de $45^\circ$. Para dicho valor de $a$, graficar ambas rectas en un mismo sistema de ejes cartesianos.





\textbf{Solución:}



El vector normal de una recta de la forma $Ax + By + C = 0$ es $(A, B)$. Por lo tanto:
\begin{itemize}
    \item Para $L_1: 3x + 5 = 0$, el vector normal es $\mathbf{n}_1 = (3, 0)$.
    \item Para $L_2: ax + 2y - 1 = 0$, el vector normal es $\mathbf{n}_2 = (a, 2)$.
\end{itemize}





El coseno del ángulo $\theta$ entre dos rectas viene dado por la expresión: \quad $
    \cos \theta = \dfrac{\mathbf{n}_1 \cdot \mathbf{n}_2}{\|\mathbf{n}_1\| \|\mathbf{n}_2\|},$

\medskip

\begin{itemize}
\begin{multicols}{3}
    \item  $    \mathbf{n}_1 \cdot \mathbf{n}_2 = (3,0) \cdot (a,2)$ 
     
    \, \qquad\  $ = 3a + 0 \cdot 2 = 3a.$
    \columnbreak

    \item $\|\mathbf{n}_1\| = \sqrt{3^2 + 0^2} = 3$
    
    \item $ \|\mathbf{n}_2\| = \sqrt{a^2 + 2^2} = \sqrt{a^2 + 4}$
\end{multicols}
\end{itemize}


\vspace{-6mm}
\begin{multicols}{2}



\begin{align*}
\text{Entonces, }\quad  \marka{a}{(-.1,.15)}  \cos 45^\circ 
    &=
    \frac{|3a|}{3\sqrt{a^2+4}} 
    \\[2mm]
    \marka{b}{(-.4,.35)} \tfrac{\sqrt{2}}{2} 
    &=
    \frac{3|a|}{3\sqrt{a^2+4}} 
    \\[2mm]
    \tfrac{\sqrt{2}}{2} 
    &=  \frac{|a|}{\sqrt{a^2+4}}
    \\[2mm]
    \tfrac{\sqrt{2}}{2} 
    &= 
    \frac{|a|}{\sqrt{a^2+4}} 
    \\[2mm]
    \left(\tfrac{\sqrt{2}}{2} \right)^2
    &= 
    \left(\frac{|a|}{\sqrt{a^2+4}} \right)^2
    \\[2mm]
     \frac{2}{4}
    &= 
    \frac{a^2}{a^2 + 4}
    \\[2mm]
    \frac{1}{2}
    &= 
    \frac{a^2}{a^2 + 4} 
    \\
    a^2 + 4 
    &= 
    2a^2
    \\
    4
    &=
     2a^2 - a^2
    \\
     4
    &= a^2 \\
    a=-2 \qquad &\text{o } \quad a= 2 
\end{align*}


\tikz[overlay,remember picture] {
\draw[thick,->,teal] (a) to[out=190,in=150] node [pos=0.65,text width=21mm , left  %right
] { \footnotesize  $\cos 45^\circ = \tfrac{\sqrt{2}}{2}$} (b) ;  
}






\textbf{ Grafico las rectas}

Para $a = 2$ y $a = -2$, graficamos las rectas en un mismo sistema de coordenadas.

\begin{center}
    \definecolor{gris}{rgb}{0.5,0.5,0.5}
\begin{tikzpicture}[scale=0.5]
% extremos para los ejes
\xdef\ext{5}
\xdef\yext{5}
\xdef\f{-5}
\xdef\g{-5}
\pgfmathtruncatemacro{\c}{\f+1}
\pgfmathtruncatemacro{\d}{\g+1}
\draw [color=gris!30!,dash pattern=on 1pt off 1pt, xstep=1cm,ystep=1cm] (\f-.3,\g-.3) grid (\ext+.3,\yext+.3);

\draw[->] (\f-.3,0) -- (\ext+.3,0) ;

\draw[->] (0,\g-.3) -- (0,\yext+.3) ;

\foreach \x in {\f,...,-1,1,2,...,\ext}
\draw[shift={(\x,0)},color=black] (0pt,2pt) -- (0pt,-2pt) node[below] {\tiny $\x$};

\foreach \y in {\g,...,-1,1,2,...,\yext }
\draw[shift={(0,\y)},color=black] (2pt,0pt) -- (-2pt,0pt) node[left] {\tiny $\y$};
    
    
            % Recta vertical x = -5/3 \draw[thick,blue,
            \draw[thick,domain=-5:5, thick, red] plot ({-5/3}, \x) node[left] {$L_1$};
            % Recta para a = 2
            \draw[thick,domain=-5:5, thick, blue] plot (\x,{-\x + 1/2}) node[text width=35mm,below right] {$L_2$ con $a=2$
            $L_2 : -2x + 2y - 1 = 0$
            };
            % Recta para a = -2
            \draw[thick,domain=-5:5, thick, violet] plot (\x,{\x + 1/2}) node[text width=35mm,right] {$L_2$ con $a=-2$
            $L_2 : -2x + 2y - 1 = 0$
            };
    \end{tikzpicture}
\end{center}

\end{multicols}





\item Hallar todas las rectas que contienen al punto $A = (3,-2)$ y forman el mismo ángulo con las rectas $L_1 : (x, y) = (1, 0) + \lambda(4, 2)$, $\lambda \in \mathbb{R}$ y $L_2 : 6x - 3y - 6 = 0$.



\textbf{Solución:}

Consideremos una recta genérica que pasa por $A$ y tiene vector director $(a, b)$:
\[ \text{(Ecuación vectorial) } \quad R :\quad  (x, y) = (3, -2) + \lambda (a, b), \quad \lambda \in \mathbb{R}. \]

El vector normal  es $\vec{n}_{R} =(-b, a)$. Utilizamos las fórmulas para calcular los ángulos entre rectas:
\[ \theta = \arccos\left( \frac{|\vec{d}_L \cdot \vec{d}_R|}{\|\vec{d}_L\| \|\vec{d}_R\|} \right) \quad \text{(para vectores directores)} \]
\[ \theta = \arccos\left( \frac{|\vec{n}_L \cdot \vec{n}_R|}{\|\vec{n}_L\| \|\vec{n}_R\|} \right) \quad \text{(para vectores normales)}. \]

Para la recta $L_1$, el vector director es $\vec{d}_{L_1} = (4, 2)$. Para la recta $L_2$, el vector normal es $\vec{n}_{L_2} = (6, -3)$. A continuación calcularemos los ángulos Igualamos los ángulos para determinar $a$ y $b$.



\begin{multicols}{2}
\textbf{Ángulo de la recta $\bf R$ con la recta $\bf L_1$:} 

\begin{itemize}
    \item $|\vec{d}_{L_1} \cdot \vec{d}|= \left| (4,2)  \cdot (a,b) \rule{0mm}{3.9mm}\right| =|4a + 2b|$

    \item $ \|\vec{d}_{L_1}\|   = \|(4,2) \| =  \sqrt{4^2 + 2^2} = \sqrt{20\,}$

    \item $\|\vec{d}_R  \| = \sqrt{a^2 + b^2} $

    
\end{itemize}



\begin{align*}
    \theta_1 &= \arccos\left( \frac{|\vec{d}_{L_1} \cdot \vec{d}|}{\|\vec{d}_{L_1}\| \|\vec{d}\|} \right) 
    \\[2mm]
    &= \arccos\left( \frac{|4a + 2b|}{\sqrt{20\,} \sqrt{a^2 + b^2}}  \right) 
\end{align*}


    \columnbreak % Forzar cambio de columna si es necesario

\textbf{Ángulo  de la recta $\bf R$ con la recta $L_2$:} 

\begin{itemize}
    \item $|\vec{n}_{L_2} \cdot \vec{n}_R|= | (6,-3)\cdot (-b,a) \rule{0mm}{3.9mm}|=|-6 \, b + (-3)a|$ \linebreak  \, \qquad \, $= |-6 \, b -3a|  $

    \item $\|\vec{n}_{L_2}\| = \| (6,- 3) \| = \sqrt{6^2 + (-3)^2} = \sqrt{45} $

    \item $\|\vec{n}_{R}\|= \| (-b,a) \| = \sqrt{(-b)^2 + a^2} = \sqrt{b^2+a^2}$

    
\end{itemize}
\begin{align*}
    \theta_2 &= \arccos\left( \frac{|\vec{n}_{L_2} \cdot \vec{n}|}{\|\vec{n}_{L_2}\| \|\vec{n}\|} \right)
    \\[2mm]
    &= \arccos\left( \frac{|-6 \, b -3a|}{\sqrt{45} \sqrt{b^2 + a^2}}  \right)
\end{align*}


\end{multicols}

Igualando estas dos ecuaciones, obtenemos valores para $a$ y $b$ que determinan el vector director de la recta buscada.

\begin{align*}
    \Cancelar[red]{\arccos} \left( \frac{|4a + 2b|}{\sqrt{20\,} \sqrt{a^2 + b^2\,}}  \right) 
    =
    \Cancelar[red]{\arccos}\left( \frac{|-6 \, b -3a|}{\sqrt{45} \sqrt{b^2 + a^2\,}}  \right)
    \qquad\Longleftrightarrow\quad
      \frac{|4a + 2b|}{\sqrt{20\,}\Cancelar[red]{\sqrt{a^2 + b^2\,}}}
    &=
    \frac{|-6 \, b -3a|}{\sqrt{45}\Cancelar[red]{ \sqrt{a^2 + b^2\,}}}
    \\[3mm]
       \left|4a + 2b \rule{0mm}{4mm} \right|    
    &=
    \sqrt{20 } \, \frac{\left|-6 \, b -3a \rule{0mm}{4mm} \right|}{\sqrt{45}}
    \\[3mm]
     \left|4a + 2b \rule{0mm}{4mm} \right|    
    &=
     \frac{2}{3} \left|-6 \, b -3 a\rule{0mm}{4mm} \right|
     \\[3mm]
     \left|4a + 2b \rule{0mm}{4mm} \right|    
    &=
      \left|\frac{2}{3}  \left( -6 \, b -3 a  \right)\rule{0mm}{4mm} \right|
    \\[3mm]
     \left|4a + 2b \rule{0mm}{4mm} \right|    
    &=
      \left|    -4 \, b -2 a  \rule{0mm}{4mm} \right|
\end{align*}

Si dos expresiones tienen el mismo valor absoluto, quiere decir que esas dos expresiones son iguales o tienen signo opuesto.
%
%

\vspace{-9mm}\begin{multicols}{2}  
\begin{align*}
\textbf{Caso 1:}\qquad   4a + 2b &= -4b - 2a  
  \\[3mm]
   4a +2a   &= -4b -  2b
   \\[3mm]
   6a &= -6 b
   \\[3mm]
   a &= -b
\end{align*} 
  
 \columnbreak % Forzar cambio de columna si es necesario 

\begin{align*}
\textbf{Caso 2:}\qquad    4a + 2b &= -(-4b - 2a ) 
  \\[1mm]
  4a + 2b &= 4b  + 2a 
  \\[1mm]
   4a - 2a   &= 4b -  2b
   \\[1mm]
   2a &= 2 b
   \\[1mm]
   a &= b
\end{align*}

\end{multicols}

 Por lo tanto, un vector director es $(a, a)$ con $a \in \mathbb{R}$, y otro vector director es $(a, -a)$ con $a \in \mathbb{R}$. En cada caso, tomamos $a=1$ para obtener un representante fijo para cada dirección.

  
\textbf{Resultados}
Las rectas que cumplen con las condiciones son:
\[ \text{Recta 1: } (x, y) = (3, -2) + \lambda (1, 1), \quad \lambda \in \mathbb{R}, \]
\[ \text{Recta 2: } (x, y) = (3, -2) + \lambda (1, -1), \quad \lambda \in \mathbb{R}. \]

\textbf{Gráfico de las rectas:} 


 \vspace{-15mm}  
\begin{center}
    \qquad\qquad\qquad\qquad\qquad
\definecolor{gris}{rgb}{0.5,0.5,0.5}
\begin{tikzpicture}[scale=0.6]
%%%% punto% 
\coordinate (A) at (3,-2);
% extremos para los ejes
\xdef\ext{9}
\xdef\yext{4}
\xdef\f{-3}
\xdef\g{-8}
\pgfmathtruncatemacro{\c}{\f+1}
\pgfmathtruncatemacro{\d}{\g+1}
\draw [color=gris!30!,dash pattern=on 1pt off 1pt, xstep=1cm,ystep=1cm] (\f-.3,\g-.3) grid (\ext+.3,\yext+.3);

\draw[->] (\f-.3,0) -- (\ext+.3,0) ;

\draw[->] (0,\g-.3) -- (0,\yext+.3) ;

\foreach \x in {\f,...,-1,1,2,...,\ext}
\draw[shift={(\x,0)},color=black] (0pt,2pt) -- (0pt,-2pt) node[below] {\tiny $\x$};

\foreach \y in {\g,...,-1,1,2,...,\yext }
\draw[shift={(0,\y)},color=black] (2pt,0pt) -- (-2pt,0pt) node[left] {\tiny $\y$};

        % Recta L1
        \draw[domain=2.15:-1.15, thick, blue] plot ({1+4*\x},{2*\x}) node[left]   {$L_1 : (x, y) = (1, 0) + \lambda(4, 2)$};

        % Recta L2
        \draw[domain=-3.1:3.1, thick, red] plot ({\x},{(6*\x - 6)/3}) node[above ]   {$L_2 : 6x - 3y - 6 = 0$};

        % Recta 1
        \draw[domain=-6.15:7.15, thick, green!70!black] plot ({3+\x},{-2+\x}) node[above]  {Recta 1};

        % Recta 2
        \draw[domain=-6.15:6.15, thick, orange] plot ({3+\x},{-2-\x}) node[right]  {Recta 2};

        % Punto 
\draw[blue,fill=blue] (A) circle (2pt) node[below]{$A = (3,-2)$} ; 


% Ángulos 
\coordinate (AA) at (-3,-8);
\coordinate (B) at ($(AA)+(1,2)$);
\coordinate (C) at ($(AA)+(1,1)$);

\pic[draw, color=magenta, fill=magenta!50!,   angle eccentricity=1.6, angle radius=.7cm, "$ \theta_2$",opacity=.8] {angle = C--AA--B};

% Ángulos 
\coordinate (AA) at (9,4);
\coordinate (B) at ($(AA)-(2,1)$);
\coordinate (C) at ($(AA)-(1,1)$);

\pic[draw, color=magenta, fill=magenta!50!,   angle eccentricity=1.6, angle radius=.7cm, "$ \theta_2$",opacity=.8] {angle = B--AA--C};

% Ángulos % Ángulos 
\coordinate (AA) at (1,0);
\coordinate (B) at ($(AA)+(2,1)$);
\coordinate (C) at ($(AA)+(1,-1)$);

\pic[draw, color=violet, fill=violet!50!,   angle eccentricity=1.4, angle radius=.6cm, "$ \theta_1$",opacity=.8] {angle = C--AA--B};

% Ángulos % Ángulos 
\coordinate (AA) at (1,0);
\coordinate (B) at ($(AA)-(1,2)$);
\coordinate (C) at ($(AA)+(1,-1)$);

\pic[draw, color=violet, fill=violet!50!,   angle eccentricity=1.4, angle radius=.4cm, "$ \theta_1$",opacity=.8] {angle = B--AA--C};


\end{tikzpicture}


\end{center}
























\item Analizar y calcular las siguientes distancias:

\begin{itemize}
    \item[(a)] Entre la recta $L : \frac{x - 2}{3} = \frac{y + 1}{-2}$ y el origen de coordenadas.
    \item[(b)] Entre la recta $M : \frac{x}{-3} + \frac{y}{4} = 1$ y el punto $A = \left(\frac{3}{2}, 6\right)$.
    \item[(c)] Entre las rectas $L_1 : \begin{cases} x &= -2 + t \\ y &= 3 - t \end{cases}, t \in \mathbb{R}$ y $L_2 : \begin{cases} x &= 5 + k \\ y &= -2 - k \end{cases}, k \in \mathbb{R}$.
    \item[(d)] Entre las rectas $T_1 : (x, y) = (1, 2) + \mu(4, -1), \mu \in \mathbb{R}$ y $T_2 : -x + 3y + 9 = 0$.
\end{itemize}





\textbf{Solución:}

\textbf{(a) Distancia entre la recta $L$ y el origen}

Primero pasamos de la ecuación  Cartesiana a la implícita: 
\[
    L:\  \frac{x - 2}{3} =  \frac{y + 1}{-2}
\quad
\Longleftrightarrow
\quad 
    -2(x - 2) = 3(y + 1)
\quad
\Longleftrightarrow
\quad
-2x + 4 = 3y + 3
\quad
\Longleftrightarrow
\quad
    2x + 3y - 1 = 0
\]

La distancia del origen $(0,0)$ a la recta es: \quad $ \displaystyle  d = \frac{|2(0) + 3(0) - 1|}{\sqrt{2^2 + 3^2}} = \frac{1}{\sqrt{13}}$. 



\begin{center}
    \definecolor{gris}{rgb}{0.5,0.5,0.5}
\begin{tikzpicture}[scale=1.3]
% extremos para los ejes
\xdef\ext{3}
\xdef\yext{2}
\xdef\f{-2}
\xdef\g{-2}
\pgfmathtruncatemacro{\c}{\f+1}
\pgfmathtruncatemacro{\d}{\g+1}
\draw [color=gris!30!,dash pattern=on 1pt off 1pt, xstep=1cm,ystep=1cm] (\f-.3,\g-.3) grid (\ext+.3,\yext+.3);

\draw[->] (\f-.3,0) -- (\ext+.3,0) ;

\draw[->] (0,\g-.3) -- (0,\yext+.3) ;

\foreach \x in {\f,...,-1,1,2,...,\ext}
\draw[shift={(\x,0)},color=black] (0pt,2pt) -- (0pt,-2pt) node[below] {\tiny $\x$};

\foreach \y in {\g,...,-1,1,2,...,\yext }
\draw[shift={(0,\y)},color=black] (2pt,0pt) -- (-2pt,0pt) node[left] {\tiny $\y$};

    \draw[thick, red, domain=-1.43:0.45] plot ({2+3*\x} , {-1-2*\x}) node[right] {$L:\  \frac{x - 2}{3} =  \frac{y + 1}{-2}$};
    \fill[violet] (0,0) circle (2pt) node[below left] {$O$};
\end{tikzpicture}
\end{center}

\textbf{(b) Distancia entre la recta $M$ y el punto $A=(\frac{3}{2},6)$}


Primero pasamos de la ecuación  Segmentaria a la implícita: 
\[
    \frac{x}{-3} + \frac{y}{4} = 1
\quad
\stackrel[\text{Multiplicamos por 12}]{ }{\Longleftrightarrow}
\quad  
    -4x + 3y = 12 
\quad
\Longleftrightarrow
\quad  4x - 3y + 12 = 0
\]

\begin{minipage}{.6\textwidth}


La distancia entre el punto $A=(\frac{3}{2},6)$ y la recta es:
\[
    d = \frac{|4(\frac{3}{2}) - 3(6) + 12|}{\sqrt{4^2 + (-3)^2}}
    = \frac{|6 - 18 + 12|}{5} = \frac{0}{5} = 0
\]
Es decir, el punto pertenece a la recta.

    
\end{minipage}\quad\begin{minipage}{.25\textwidth}
    
\begin{center}
    \definecolor{gris}{rgb}{0.5,0.5,0.5}
\begin{tikzpicture}[scale=.5]
% extremos para los ejes
\xdef\ext{3}
\xdef\yext{7}
\xdef\f{-5}
\xdef\g{-1}
\pgfmathtruncatemacro{\c}{\f+1}
\pgfmathtruncatemacro{\d}{\g+1}
\draw [color=gris!30!,dash pattern=on 1pt off 1pt, xstep=1cm,ystep=1cm] (\f-.3,\g-.3) grid (\ext+.3,\yext+.3);

\draw[->] (\f-.3,0) -- (\ext+.3,0) ;

\draw[->] (0,\g-.3) -- (0,\yext+.3) ;

\foreach \x in {\f,...,-1,1,2,...,\ext}
\draw[shift={(\x,0)},color=black] (0pt,2pt) -- (0pt,-2pt) node[below] {\tiny $\x$};

\foreach \y in {\g,...,-1,1,2,...,\yext }
\draw[shift={(0,\y)},color=black] (2pt,0pt) -- (-2pt,0pt) node[left] {\tiny $\y$};
% 

\coordinate (A) at (-3,0);%punto izquierda
\coordinate (B) at (0,4);%punto derecha
\coordinate (m) at ($1.3*(A) -.3*(B)$);%Punto por detras de A
\coordinate (M) at ($-.8*(A)+1.8*(B)$);%Punto por adelante  de B

\fill[shift={(A)},white] (0mm,-3mm) circle (3mm);
\draw[shift={(A)},color=violet] (0pt,2pt) -- (0pt,-2pt) node[below] {\footnotesize$-3$};

\fill[shift={(B)},white] (-3mm,0mm) circle (3mm);
\draw[shift={(B)},color=violet] (2pt,0pt) -- (-2pt,0pt) node[left] {\footnotesize $4$};

    \draw[thick, blue] (m) -- (M) node[right] {$M$};

    
    \coordinate (A) at (3/2,6);%punto 
    \fill[magenta] (A) circle (2pt) node[right] {$A$};
\end{tikzpicture}
\end{center}


\end{minipage}









\textbf{(c) Distancia entre las rectas $L_1 : \begin{cases} x &= -2 + t \\ y &= 3 - t \end{cases}, t \in \mathbb{R}$ y $L_2 : \begin{cases} x &= 5 + k \\ y &= -2 - k \end{cases}, k \in \mathbb{R}$}


Estas rectas son paralelas porque sus vectores directores son $(1, -1)$, por lo que calculamos la distancia entre un punto de una recta y la otra.

Primero pasaremos de la ecuación paramétrica a la implícita.


$\left\{
    \begin{array}{rl}
         x&= -2+t 
         \\[3mm]
        y & =3-t 
    \end{array}
    \right.
\ \ 
\Longleftrightarrow
\ \ 
\left\{
    \begin{array}{rl}
         x+2&=t  
         \\[3mm]
        y-3 & =-t  
    \end{array}
    \right.
\ \ 
\Longleftrightarrow
\ \ 
\left\{
    \begin{array}{rl}
         x+2&=t  \indicador{x}{(.1,.15)}{y}{(.8,-.35)}\marka{z}{(1.3,-.35)}
         \\[3mm]
        -y+3 & =t   \marka{t}{(.1,.15)} 
    \end{array}
    \right.
$

\tikz[overlay,remember picture] {
%%%%%%%%%%%%%%%%%%%%%%%% 
\draw [overlay,->,very thick,red,opacity=.5](x)-|(y)--(z);
\draw [overlay,very thick,red,opacity=.5](t)-|(y);
\node  at (z) [text width=60mm, below right , yshift=3mm ]{
$
x+2=-y+3
\quad
\Longleftrightarrow
\quad
x+y-1 = 0
$
};
%%%%%%%%%%%%%%%%%%%%%%%%
}

    


Entonces tenemos que la forma implícita de recta $L_1$ es  $L_1: \ x+y-1 = 0 $, el punto $P=(5,-2)$ pertenece a la recta $L_2$ y las dos rectas son paralelas, entonces la distancia entre las rectas es igual a la distancia entre $L_1$ y $P$:


\begin{equation*}
    d = \frac{| 5  + (-2) -1|}{\sqrt{1^2 + (-1)^2}} = \frac{|2|}{\sqrt{2}} = \frac{2}{\sqrt{2}} = \sqrt{2}
\end{equation*}





\begin{center}
    \definecolor{gris}{rgb}{0.5,0.5,0.5}
\begin{tikzpicture}[scale=.9]
% extremos para los ejes
\xdef\ext{5}
\xdef\yext{4}
\xdef\f{-3}
\xdef\g{-2}
\pgfmathtruncatemacro{\c}{\f+1}
\pgfmathtruncatemacro{\d}{\g+1}
\draw [color=gris!30!,dash pattern=on 1pt off 1pt, xstep=1cm,ystep=1cm] (\f-.3,\g-.3) grid (\ext+.3,\yext+.3);

\draw[->] (\f-.3,0) -- (\ext+.3,0) ;

\draw[->] (0,\g-.3) -- (0,\yext+.3) ;

\foreach \x in {\f,...,-1,1,2,...,\ext}
\draw[shift={(\x,0)},color=black] (0pt,2pt) -- (0pt,-2pt) node[below] {\tiny $\x$};

\foreach \y in {\g,...,-1,1,2,...,\yext }
\draw[shift={(0,\y)},color=black] (2pt,0pt) -- (-2pt,0pt) node[left] {\tiny $\y$};

    \draw[thick, red, domain=-1.15:5.15] plot ({-2+\x} , {3-\x}) node[below left] {$L_1 : \begin{cases} x &= -2 + t \\ y &= 3 - t \end{cases}, t \in \mathbb{R}$};

    \draw[thick, magenta, domain=-6.15:0.15] plot ({5+\x} , {-2-\x}) node[below right] {$L_2 : \begin{cases} x &= 5 + k \\ y &= -2 - k \end{cases}, k \in \mathbb{R}$};
    \fill[violet] (0,0) circle (2pt) node[below left] {$O$};


    \coordinate (P) at (5,-2);%punto 
    \fill[magenta] (P) circle (2pt) node[right] {$P$};
\end{tikzpicture}
\end{center}





\textbf{(d) Distancia entre $T_1$ y $T_2$}

{\bf Recordatorio:} En el plano, dos rectas pueden ser paralelas o incidentes. Si son incidentes, la distancia entre ellas es cero.

\medskip

La recta \( T_1 \) pasa por el punto \( (1,2) \) y tiene como vector director \( (4,-1) \), mientras que la recta \( T_2 \) tiene como vector normal \( \vec{n}_{T_2} = (-1,3) \).  Tenemos la siguiente:
\begin{align*} 
\vec{d}_{T_1} \cdot \vec{n}_{T_2} &= (1,2)\cdot(4,-1) 
\\
&= 1\cdot 4 + 2 \cdot (-1) = 2 \neq 0, \quad \text{ por lo tanto no son rectas paralelas}
\end{align*}
 Dado que las rectas no son paralelas, necesariamente son incidentes, por lo que la distancia entre ellas es cero.







\begin{center}
    \definecolor{gris}{rgb}{0.5,0.5,0.5}
\begin{tikzpicture}[scale=.9]
% extremos para los ejes
\xdef\ext{13}
\xdef\yext{4}
\xdef\f{-1}
\xdef\g{-2}
\pgfmathtruncatemacro{\c}{\f+1}
\pgfmathtruncatemacro{\d}{\g+1}
\draw [color=gris!30!,dash pattern=on 1pt off 1pt, xstep=1cm,ystep=1cm] (\f-.3,\g-.3) grid (\ext+.3,\yext+.3);

\draw[->] (\f-.3,0) -- (\ext+.3,0) ;

\draw[->] (0,\g-.3) -- (0,\yext+.3) ;

\foreach \x in {\f,...,-1,1,2,...,\ext}
\draw[shift={(\x,0)},color=black] (0pt,2pt) -- (0pt,-2pt) node[below] {\tiny $\x$};

\foreach \y in {\g,...,-1,1,2,...,\yext }
\draw[shift={(0,\y)},color=black] (2pt,0pt) -- (-2pt,0pt) node[left] {\tiny $\y$};

%recta T_1
    \draw[thick, blue, domain=-.5:3.15] plot ({1+4*\x} , {2-\x}) node[below ] {$T_1 : (x, y) = (1, 2) + \mu(4, -1), \mu \in \mathbb{R}$};
%recta T_2
    \draw[thick, magenta, domain=-2.15:1.555,variable=\y] plot ({3*\y+9} , {\y}) node[below right] {$T_2 :\  -x + 3y + 9 = 0$};

    
    

    \coordinate (P) at (1,2);%punto 
    %\fill[blue] (P) circle (2pt) node[right] {$P$};
\end{tikzpicture}
\end{center}



















  \item \,
\begin{itemize}
    \item[(a)] Hallar los puntos pertenecientes a la recta $L : 5x - y + 6 = 0$ que distan 2 unidades del punto $P = (-3, 1)$.
    \item[(b)] Determinar el o los valores que puede tomar $k \in \mathbb{R}$ para que la distancia de la recta $L : -8x + 6y = k$ al punto $P = (1, 1)$ sea $\frac{1}{2}$.
    \item[(c)] Encontrar el o los puntos de la recta $L_1 : y = 3x - 2$ que distan $\sqrt{5}$ unidades de la recta \linebreak $L_2 : -x + 2y - 3 = 0$.
\end{itemize}
   











\textbf{Solución:}

\textbf{(a) Puntos de la recta a distancia 2 del punto $P$}

\begin{itemize}
\begin{multicols}{2}  
\item La ecuación de la recta es:
\begin{align*}
 5x - y + 6 &= 0 
\intertext{despejando y, tenemos:}
 5x   + 6 &= y 
\end{align*}


 \columnbreak % Forzar cambio de columna si es necesario
\item La distancia entre dos puntos $A= (x_0,y_0)$ y $B=(x_1,y_1)$ está dada por:
\[
    d \left( A,B\right)= \sqrt{( x_0-x_1 )^2 + (y_0 -y_1)^2 \, }
\]

\end{multicols}
\end{itemize}


\begin{itemize}
    \item Un punto genérico de la recta es de la forma: $A=(x_0,\, 5x_0   + 6 )$ y tendría que verificar que su distancia al punto $P= (-3,1)$ sea $2$, entonces nos queda:
    \begin{align*}
       d \left( A,B\right)
       &= 2
       \\
       \sqrt{\left( x_0-(-3) \rule{0mm}{3.6mm} \right)^2 + \left( 5x_0   + 6 -1 \rule{0mm}{3.6mm}\right)^2 \, }
       &= 2
       \\
       \left( x_0+3 \rule{0mm}{3.6mm} \right)^2 + \left( 5x_0   + 5\rule{0mm}{3.6mm}\right)^2 \, 
       &= 4
       \\
       x_0\rule{0mm}{3.5mm}^2 + 6 x_0 +9 + 25\, x_0\rule{0mm}{3.5mm}^2   +10 x_0 +25  \, 
       &= 4
       \\
        26\, x_0\rule{0mm}{3.5mm}^2 + 56 x_0 + 34 \, 
       &= 4
       \\
        26\, x_0\rule{0mm}{3.5mm}^2 + 56 x_0 + 34-4 \, 
       &= 0
       \\
        26\, x_0\rule{0mm}{3.5mm}^2 + 56 x_0 + 30 \, 
       &= 0
    \end{align*}


    Aplicando Bhaskara tenemos: $x_0= -\tfrac{15}{13}, x_0= -1$.


\begin{itemize}
    \item Para $x_0= -\tfrac{15}{13}$ tenemos que $y_0= 5\cdot (-\tfrac{15}{13}) +6= \tfrac{3}{13}$

    \item Para $x_0= -1$ tenemos que $y_0= 5\cdot (-1) +6= 1 $
\end{itemize}

    
    
\end{itemize}



\begin{center}
        \definecolor{gris}{rgb}{0.5,0.5,0.5}
\begin{tikzpicture}[scale=2.1]
% extremos para los ejes
\xdef\ext{2}
\xdef\yext{2}
\xdef\f{-3}
\xdef\g{-1}
\pgfmathtruncatemacro{\c}{\f+1}
\pgfmathtruncatemacro{\d}{\g+1}
\draw [color=gris!30!,dash pattern=on 1pt off 1pt, xstep=1cm,ystep=1cm] (\f-.3,\g-.3) grid (\ext+.3,\yext+.3);

\draw[->] (\f-.3,0) -- (\ext+.3,0) ;

\draw[->] (0,\g-.3) -- (0,\yext+.3) ;

\foreach \x in {\f,...,-1,1,2,...,\ext}
\draw[shift={(\x,0)},color=black] (0pt,2pt) -- (0pt,-2pt) node[below] {\tiny $\x$};

\foreach \y in {\g,...,-1,1,2,...,\yext }
\draw[shift={(0,\y)},color=black] (2pt,0pt) -- (-2pt,0pt) node[left] {\tiny $\y$};

         \coordinate (P) at (-3,1);%punto 

        \draw[thick, magenta, domain=-1.47:-.79] plot ({\x} , {5*\x+6})  node[above right] {$L : 5x - y + 6 = 0$};
        \fill[red] (P) circle (1pt) node[below left] {$P=(-3,1)$};  

        \begin{scope}
        \clip (-3.1,1.2) --(-1,-1)--(0,2)--cycle;
        \draw[dashed] (P) circle (2cm);
        \end{scope}

        \coordinate (A) at (-1,1);%punto 
        \coordinate (B) at (-15/13,3/13);%punto 
        \fill[red] (A) circle (1pt) node[above right] {$A=(-1,1)$};
        \fill[red] (B) circle (1pt) node[right] {$P=(-\tfrac{15}{13}, \tfrac{3}{13})$};
    \end{tikzpicture}
\end{center}

\textbf{(b) Valores de $k$ para distancia $\frac{1}{2}$}


La ecuación de la recta es:
\[-8x + 6y = k
\quad
\Longleftrightarrow
\quad 
-8x + 6y - k = 0
\]
Utilizando la fórmula de la distancia de un punto a una recta y el punto $P = (1,1)$:
\begin{align*}
    \frac{| -8(1) + 6(1) - k |}{\sqrt{(-8)^2 + 6^2}} 
    &= 
    \frac{1}{2}
    \\
    \frac{| -2 - k |}{\sqrt{64 + 36}} 
    &= 
    \frac{1}{2}
    \\
    \frac{| -2 - k |}{\sqrt{100}} 
    &= 
    \frac{1}{2}
    \\ 
    \frac{| -2 - k |}{10} 
    &= 
    \frac{1}{2}
    \\ 
    | -2 - k | 
    &= 
    10 \cdot \frac{1}{2}
    \\ 
    | -2 - k | 
    &= 
    5  
\end{align*}  
\[
\begin{array}{rl@{\hspace{20mm}}c @{\hspace{20mm}} rl r}   
-2-k
\qquad
&= \qquad
5
& \mbox{ o }
&
-2-k
\qquad
&= \qquad
-5
\\[3mm] %%%%%%%%%%%%%%%% %%%%%%%%%%%%%%%% %%%%%%%%%%%%%%%% %%%%%%%%%%%%%%%% %%%%%%%%%%%%%%%% %%%%%%%%%%%%%%%%  
-2-5
&= 
k
& %\mbox{ o }
&
-2-(-5) 
&= 
k
\\[3mm] %%%%%%%%%%%%%%%% %%%%%%%%%%%%%%%% %%%%%%%%%%%%%%%% %%%%%%%%%%%%%%%% %%%%%%%%%%%%%%%% %%%%%%%%%%%%%%%%  
-7
&= 
k
& % \mbox{ o }
&
3 
&= 
k
\end{array}
\]
Obtenemos:
\[ k = -7 \quad \text{o} \quad k = 3 \]

\begin{center}
  \definecolor{gris}{rgb}{0.5,0.5,0.5}
\begin{tikzpicture}[scale=1.6]
% extremos para los ejes
\xdef\ext{2}
\xdef\yext{2}
\xdef\f{-2}
\xdef\g{-1}
\pgfmathtruncatemacro{\c}{\f+1}
\pgfmathtruncatemacro{\d}{\g+1}
\draw [color=gris!30!,dash pattern=on 1pt off 1pt, xstep=1cm,ystep=1cm] (\f-.3,\g-.3) grid (\ext+.3,\yext+.3);

\draw[->] (\f-.3,0) -- (\ext+.3,0) ;

\draw[->] (0,\g-.3) -- (0,\yext+.3) ;

\foreach \x in {\f,...,-1,1,2,...,\ext}
\draw[shift={(\x,0)},color=black] (0pt,2pt) -- (0pt,-2pt) node[below] {\tiny $\x$};

\foreach \y in {\g,...,-1,1,2,...,\yext }
\draw[shift={(0,\y)},color=black] (2pt,0pt) -- (-2pt,0pt) node[left] {\tiny $\y$};


        \draw[thick, magenta, domain=-.2:2.5] plot ({\x} , {(8*\x-7)/6})  node[above right] {$L$ para $k=-7$};
        \draw[thick, violet, domain=-1.4:1.4] plot ({\x} , {(8*\x+3)/6})  node[above right] {$L$ para $k=3$};
        \fill[red] (1,1) circle (1pt) node[above right] {$P(1,1)$};
    \end{tikzpicture}
\end{center}

\textbf{(c) Puntos en $L_1$ que distan $\sqrt{5}$ de $L_2$}
Las rectas son:
\[ L_1: y = 3x - 2, \quad L_2: -x + 2y - 3 = 0 \]
La distancia entre un punto $(x_0, y_0)$ de $L_1$ y la recta $L_2$ debe ser $\sqrt{5}$.

Tenemos: 

\begin{itemize}
    \item El punto $ P = (x_0, y_0) $ de $ L_1 $ satisface la ecuación $ y_0 = 3x_0 - 2 $, por lo que se puede expresar como $ P = (x_0, 3x_0 - 2) $.

    \item La distancia entre un punto genérico de $ L_1 $,  dado por $ P = (x_0, 3x_0 - 2) $, y la recta $ L_2: -x + 2y - 3 = 0 $ está dada por la expresión:

\[
\frac{\left| x_0 + 2(3x_0 - 2) - 3 \right|}{\sqrt{(-1)^2 + 2^2\,}}
\]
\end{itemize}
Como la distancia debe ser igual a $\sqrt{5}$ tenemos que resolver la siguiente ecuación: 
\begin{align*}
    \frac{| -x_0 + 2(3\, x_0 - 2) - 3 |}{\sqrt{(-1)^2 + 2^2}}  
    &= 
    \sqrt{5}
    \\
     \frac{| -x_0 +6\, x_0 - 4 - 3 |}{\sqrt{1 + 4}}  
    &= 
    \sqrt{5}
    \\
    \frac{| 5\, x_0 - 7 |}{\sqrt{ 5}}  
    &= 
    \sqrt{5}
    \\
    |7\, x_0 - 7 |
    &= 
    \sqrt{5} \cdot \sqrt{5}
    \\
    | 5\, x_0 - 7 |
    &= 
    5 
\end{align*}
\[
\begin{array}{rl@{\hspace{20mm}}c @{\hspace{20mm}} rl r}   
5\, x_0 - 7
\ 
&= \ 
5
& \mbox{ o }
&
5\, x_0 - 7
\
&= \
-5
\\[3mm] %%%%%%%%%%%%%%%% %%%%%%%%%%%%%%%% %%%%%%%%%%%%%%%% %%%%%%%%%%%%%%%% %%%%%%%%%%%%%%%% %%%%%%%%%%%%%%%%  
5\, x_0 
&= 
12
& %\mbox{ o }
&
5\, x_0  
&= 
2
\\[3mm] %%%%%%%%%%%%%%%% %%%%%%%%%%%%%%%% %%%%%%%%%%%%%%%% %%%%%%%%%%%%%%%% %%%%%%%%%%%%%%%% %%%%%%%%%%%%%%%%  
x_0
&= 
\tfrac{12}{5}
& % \mbox{ o }
&
x_0
&= 
\tfrac{2}{5}
\end{array}
\]
Obtenemos:
\begin{itemize}
    \item Si $x_0 = \tfrac{12}{5}$ tenemos que $y_0 = 3 \cdot \tfrac{12}{5} -2 =\tfrac{26}{5}$, con lo cual $P=(\tfrac{12}{5}, \tfrac{26}{5})$.

    \item Si $x_0 = \tfrac{2}{5}$ tenemos que $y_0 = 3 \cdot\tfrac{2}{5} -2= -\tfrac{4}{5}$, con lo cual $P=(\tfrac{2}{5},- \tfrac{4}{7} )$.
\end{itemize}



{\bf Grafico}
\[ (1,1) \quad \text{y} \quad (-1,-5) \]

 \begin{center}
\definecolor{gris}{rgb}{0.5,0.5,0.5}
\begin{tikzpicture}[scale=0.7]
% extremos para los ejes
\xdef\ext{4}
\xdef\yext{7}
\xdef\f{-4}
\xdef\g{-2}
\pgfmathtruncatemacro{\c}{\f+1}
\pgfmathtruncatemacro{\d}{\g+1}
\draw [color=gris!30!,dash pattern=on 1pt off 1pt, xstep=1cm,ystep=1cm] (\f-.3,\g-.3) grid (\ext+.3,\yext+.3);

\draw[->] (\f-.3,0) -- (\ext+.3,0) ;

\draw[->] (0,\g-.3) -- (0,\yext+.3) ;

\foreach \x in {\f,...,-1,1,2,...,\ext}
\draw[shift={(\x,0)},color=black] (0pt,2pt) -- (0pt,-2pt) node[below] {\tiny $\x$};

\foreach \y in {\g,...,-1,1,2,...,\yext }
\draw[shift={(0,\y)},color=black] (2pt,0pt) -- (-2pt,0pt) node[left] {\tiny $\y$}; 


        \draw[thick,blue,domain=-.1:3.1] plot ({\x}, {3*\x-2}) node[right] {$ L_1:\ y = 3x - 2$};
       \draw[thick,blue,domain=-.66:3.66,variable=\y] plot ({2*\y-3}, {\y}) node[right] {$L_2:\ -x + 2y - 3 = 0$};
        


        \coordinate (P1) at (12/5,26/5);%punto 
        \coordinate (P2) at (2/5,-4/5);%punto 

        \coordinate (P11) at ($(P1)-(-1,2)$);%punto 
        \coordinate (P22) at ($(P2)+(-1,2)$);%punto 

        \fill[red] (P1) circle (2pt) node[above left] {$(\tfrac{12}{5}, \tfrac{26}{5})$};
        \fill[red] (P2) circle (2pt) node[below right] {$(\tfrac{2}{5}, -\tfrac{4}{5})$};
        
         % distancia
          \begin{scope}
        %\clip (-2.1,1.2) --(3,2.5)--(3,5)-- (-2.1,5)--cycle;
        \draw[dashed] (P1) --(P11);
         \draw[dashed] (P2) --(P22);
        \end{scope}

        
    \end{tikzpicture}
\end{center}









   \item Sean el punto $A = (1,-2)$ y la recta $L :\ (x, y) = \lambda(3, 4)$, $\lambda \in \mathbb{R}$:

\begin{itemize}
    \item[(a)] ¿Cuál es el punto $P$ de $L$ que se encuentra a menor distancia de $A$? ¿Cuánto vale dicha distancia?
    \item[(b)] Calcular el vector $\overrightarrow{PA}$. ¿Qué se obtiene al hacer $\overrightarrow{PA} \cdot \overrightarrow{u_L}$? ¿Por qué? (Teniendo en cuenta que es el vector director de la recta).
    \item[(c)] Hallar, si es posible, una recta paralela a $L$ que pase por $A$ y que esté a 4 unidades de distancia de $L$.
    \item[(d)] Hallar, si es posible, una recta paralela a $L$ que pase por $A$ y que esté a 2 unidades de distancia de $L$.
    \item[(e)] ¿Qué conclusiones se pueden sacar después de resolver los incisos (c) y (d)?
\end{itemize}




\textbf{Recordemos:}

\textbf{\textbf{\textcolor{green!80!blue!50!}{Distancia Entre Rectas.}}}

En el plano, analizamos previamente las posibles intersecciones entre dos rectas \( L \) y \( R \). Consideremos los siguientes casos:

\begin{itemize}
    \item[\checkmark] Si \( L \cap R \neq \emptyset \), entonces \( d(L,R) = 0 \).
    
  {\footnotesize  Esto ocurre cuando las rectas tienen al menos un punto en común. En este caso, la  \textcolor{red}{distancia}  es cero, ya que existe un punto que pertenece simultáneamente a ambas rectas. Esto puede darse en dos situaciones: cuando las rectas se intersectan en un único punto (rectas secantes) o cuando coinciden completamente (rectas coincidentes).}
    
    \item[\checkmark] Si \( L \cap R = \emptyset \) y \( L \paralela R \) (rectas paralelas no coincidentes), entonces:
    
    \[
    d(L,R) = d(P_1, R) = d(L, P_2), \quad \text{donde } P_1 \in L \text{ y } P_2 \in R .
    \]

{\footnotesize En otras palabras, la \textcolor{red}{distancia} entre dos rectas paralelas no coincidentes es igual a la \textcolor{red}{distancia} entre un punto de una de las rectas y la otra recta. Dicho punto puede pertenecer a \( L \) y calcularse su distancia a \( R \), o pertenecer a \( R \) y calcularse su distancia a \( L \).}


\end{itemize}

\textbf{Solución:}



\textbf{(a) Punto más cercano a $A$ en $L$ y su distancia}


El punto más cercano $P$ es la proyección ortogonal de $A$ sobre la recta $L$. 

\begin{minipage}{.3\textwidth}

\begin{center}
    \definecolor{gris}{rgb}{0.5,0.5,0.5}
\begin{tikzpicture}[scale=.7]
% extremos para los ejes
\xdef\ext{2}
\xdef\yext{2}
\xdef\f{-3}
\xdef\g{-2}
\pgfmathtruncatemacro{\c}{\f+1}
\pgfmathtruncatemacro{\d}{\g+1}
\draw [color=gris!30!,dash pattern=on 1pt off 1pt, xstep=1cm,ystep=1cm] (\f-.3,\g-.3) grid (\ext+.3,\yext+.3);

\draw[->] (\f-.3,0) -- (\ext+.3,0) ;

\draw[->] (0,\g-.3) -- (0,\yext+.3) ;

\foreach \x in {\f,...,-1,1,2,...,\ext}
\draw[shift={(\x,0)},color=black] (0pt,2pt) -- (0pt,-2pt) node[below] {\tiny $\x$};

\foreach \y in {\g,...,-1,1,2,...,\yext }
\draw[shift={(0,\y)},color=black] (2pt,0pt) -- (-2pt,0pt) node[left] {\tiny $\y$};

    \draw[thick, red, domain=-.56:.58] plot ({3*\x} , {4*\x}) node[above ,xshift=3mm] {$L :\ (x, y) = \lambda(3, 4)$};


    \coordinate (A) at (1,-2);%punto 
    \fill[magenta] (A) circle (2pt) node[right] {$A$};

    \draw[dashed, magenta, domain=-.1:1.15] plot ({1-4*\x} , {-2+3*\x});

     \coordinate (P) at (-3/5,-4/5);%punto 
    \fill[magenta] (P) circle (2pt) node[below right] {$P$};
    
    %\draw[dashed]  (A)  circle (2);   
\end{tikzpicture}
\end{center}
\end{minipage}\qquad\begin{minipage}{.6\textwidth} 
Para hallar a $P$, trazamos $L_0$ que es la recta perpendicular a $L$, en la intersección entre esa perpendicular y la recta $L$ esta el punto buscado, es decir $L \cap L_0 = \{P\}$.

El vector director de la recta es $\vec{u}_L = (3,4)$, $\vec{u}_{L_0}= (-4,3)$  y la ecuación de la recta perpendicular a $L$ que pasa por $A$ es:
\[ L_0 : \quad (x, y) = (1,-2) + t(-4,3), \quad t \in \R  \] 
\end{minipage}


%Para calcular $  L \cap L_0$ Igualando  y resolviendo para $t$ y $\lambda$:

$
\begin{array}{lrl}
  L \cap L_0:\quad & \lambda(3,4) &= (1,-2) + t(-4,3)
   \\[4mm]
   &\left(3\lambda,4\lambda \rule{0mm}{3.9mm}\right) &= \indicador{A}{(-.15,-.1)}{B}{(.3,-1.7)} \left(1- 4t,-2 + 3 t \rule{0mm}{3.9mm} \right)  
\end{array}
$




\tikz[overlay,remember picture, rotate=3 ,every node/.style={rotate=3, transform shape}] {
\draw [overlay,->,very thick,red,opacity=.5](A)|-(B); 
\node  at (B) [text width=41mm, below right , yshift=11.5mm ]{
{\scriptsize igualando las coordenadas:}\\
$
\left\{
\begin{array}{rl}
3\lambda &= 1-4\, t     
\\[5mm]
4\lambda   &=  -2+3\, t
\end{array}
\right.  \indicador{X}{(.1,.35)}{Y}{(.6,.45)}  
$
};
%%%%%%%%%%------------------------------------------------------------------------------------------------------------
\draw [overlay,->,very thick,red,opacity=.5](X)--(Y); 
\node  at (Y) [text width=41mm, below right , yshift= 9mm ]{ 
$
\left\{
\begin{array}{rl}
\lambda &=\frac{ 1-4\, t }{3}\marka{a}{(.1,.55)}
\\[5mm]
\lambda   &=  \frac{ -2+3\, t}{4} \indicador{b}{(.1,.35)}{c}{(.6,.95)} \marka{d}{(1.4,1.05)}
\end{array}
\right.
$
};
%%%%%%%%%%------------------------------------------------------------------------------------------------------------
\draw [overlay,->,very thick,red,opacity=.5](b)-|(c)--(d);  
\draw [overlay,very thick,red,opacity=.5](a)++(0,-.1)-|(c);
\node  at (d) [text width=41, below right , yshift= 10mm ]{
\begin{align*}
    \frac{ 1-4\, t }{3}  &=  \frac{ -2+3\, t}{4} 
    \\
    4( 1-4\, t)  &= 3 ( -2+3\, t)
    \\
    4 -16\, t  &= -6+9\, t
    \\
    -16\, t -9\, t  &= -6-4
    \\
    -25 t  &= -10
    \\
    t&= \tfrac{2}{5}
\end{align*} 
};
}
%%%%%%%%%%%%%%%%%%%%%%%%%%%%%%%%%%%%%%%%%%%%%%%%%%%%%%%%%%%%%%%%%%%%%%%%%%%%%%%%%%%%%%%%%%%%%%%%%%%%%%%%%%%%%%%%%%%%%%%%%%%%%%%%%%%%%%%%%%%%%%%%%%%%%%%%%%%%%%%%%%%%%%%%%%%%%%%%%%%%




\vspace{33mm}


Luego reemplazando $t$ en la ecuación de $L$ tenemos el punto $P$:
\begin{multicols}{2}
\begin{align*}
  P=  (x, y) 
    &= (1,-2) + \tfrac{2}{5}  \left(-4,3\right)
    \\
    &= (1,-2) +   \left(-  \tfrac{8}{5}, \tfrac{6}{5}\right)
    \\
    &=\left(-\tfrac{3}{5},  -\tfrac{4}{5} \right)
\end{align*}

\begin{center}
    \definecolor{gris}{rgb}{0.5,0.5,0.5}
\begin{tikzpicture}[scale=.6]
% extremos para los ejes
\xdef\ext{2}
\xdef\yext{2}
\xdef\f{-3}
\xdef\g{-2}
\pgfmathtruncatemacro{\c}{\f+1}
\pgfmathtruncatemacro{\d}{\g+1}
\draw [color=gris!30!,dash pattern=on 1pt off 1pt, xstep=1cm,ystep=1cm] (\f-.3,\g-.3) grid (\ext+.3,\yext+.3);

\draw[->] (\f-.3,0) -- (\ext+.3,0) ;

\draw[->] (0,\g-.3) -- (0,\yext+.3) ;

\foreach \x in {\f,...,-1,1,2,...,\ext}
\draw[shift={(\x,0)},color=black] (0pt,2pt) -- (0pt,-2pt) node[below] {\tiny $\x$};

\foreach \y in {\g,...,-1,1,2,...,\yext }
\draw[shift={(0,\y)},color=black] (2pt,0pt) -- (-2pt,0pt) node[left] {\tiny $\y$};

    \draw[thick, red, domain=-.56:.58] plot ({3*\x} , {4*\x}) node[ right] {$L :\ (x, y) = \lambda(3, 4)$};


    \coordinate (A) at (1,-2);%punto 
    \fill[magenta] (A) circle (2pt) node[right] {$A$};

    \draw[dashed, magenta, domain=-.1:1.15] plot ({1-4*\x} , {-2+3*\x});

     \coordinate (P) at (-3/5,-4/5);%punto 
    \fill[magenta] (P) circle (2pt) node[right] {$P$};
    
    %\draw[dashed]  (A)  circle (2);   
\end{tikzpicture}
\end{center}
\end{multicols}

\hfill \textbf{Rta:} $P=\left( -\frac{3}{5},-\frac{4}{5} \right)$.











\begin{tcolorbox}[colback=white,colframe=red!60!black,title= \textbf{Observaciones:}]


Si te olvidas la formula de distancia entre recta y punto, o no te dan la recta en su forma implícita, podes calcular la distancia proyectando como lo hicimos en este ejercicio.


\bigskip

si te piden calcular el punto que efectiviza la distancia entre una recta y otro punto dado, bueno tenemos que resolver con el método anterior. 

\bigskip

Si recordamos de la formula, la resolución seria la siguiente, donde  
buscaremos la ecuación implícita de la recta:

$L :\ (x, y) =  \lambda(3, 4)$, $\lambda \in \mathbb{R}$

\qquad \quad$
\begin{array}{rl}
   L :\ (x, y) &= \lambda(3, 4),\  \lambda \in \mathbb{R}
   \\[2mm]
  (x, y)  &= \indicador{A}{(-.15,-.1)}{B}{(.7,-1.4)} \left(3\lambda,4\lambda \rule{0mm}{3.5mm}\right)  
\end{array}
$




\tikz[overlay,remember picture, rotate=2 ,every node/.style={rotate=1, transform shape}] {
\draw [overlay,->,very thick,red,opacity=.5](A)|-(B); 
\node  at (B) [text width=43mm, below right , yshift=12mm ]{
\hspace{-5mm}{\scriptsize Igualando las coordenadas:}\\
$
\left\{
\begin{array}{rl}
x  &= 3\lambda  
\\[4mm]
y    &=  4\lambda
\end{array}
\right.  \indicador{X}{(.1,.15)}{Y}{(1.4,.3)}  
$
};
%%%%%%%%%%------------------------------------------------------------------------------------------------------------
\draw [overlay,->,very thick,red,opacity=.5](X)--(Y); 
\node  at (Y) [text width=41mm, below right , yshift=11mm ]{
$
\left\{
\begin{array}{rl}
\dfrac{x}{3}  &= \lambda  \marka{a}{(.1,.35)}
\\[4mm]
\dfrac{y}{4}    &= \lambda \indicador{b}{(.1,.25)}{c}{(.6,.75)} \marka{d}{(1.9,.85)}
\end{array}
\right.   
$
}; 
%%%%%%%%%%------------------------------------------------------------------------------------------------------------
\draw [overlay,->,very thick,red,opacity=.5](b)-|(c)--(d);  
\draw [overlay,very thick,red,opacity=.5](a)++(0,-.1)-|(c);
\node  at (d) [text width=21, below right , yshift= 11mm ]{
\begin{align*}
    \frac{x}{3}  &=  \frac{y}{4} 
    \\
    4 x  &= 3 \, y
    \\
 \hspace{-5mm}   4 x -3y  &=  0 .
\end{align*} 
};
}

\vspace{21mm}

\hfill La ecuación implícita de la recta $L$ es : $ 4 x -3y  =  0 $

Usamos la fórmula de distancia punto-recta:
\[ \frac{ \left|4(1) - 3(-2) \right|  }{\| (4,-3) \|} = \frac{ \left|4-  (-6) \right|  }{\sqrt{4^2+(-3)^2\,}} =   \frac{ \left|10 \right|  }{\sqrt{25\,}} = \frac{10}{5} = 2. \]


\end{tcolorbox}












\bigskip 


\textbf{(b) Cálculo de $\overrightarrow{PA}$ y su producto escalar con el director de $L$.}

\[ \overrightarrow{PA} = (1,-2) - \left(-\tfrac{3}{5},  -\tfrac{4}{5} \right) = (\tfrac{8}{5},-\tfrac{6}{5}). \]

El producto escalar con el vector director es:
\[ \overrightarrow{PA} \cdot \overrightarrow{u_L} = (\tfrac{8}{5},-\tfrac{6}{5}) \cdot (3,4) =\tfrac{24}{5}- \tfrac{24}{5}= 0. \]

El resultado es cero, es que deben ser vectores perpendiculares, El punto de la recta $L$ que efectiviza la distancia al punto $A$ determina con $A$ una dirección perpendicular a recta. 





\bigskip


\textbf{(c) Recta paralela a $L$ que pase por $A$ y que esté a 4 unidades de distancia de $L$.} 


Debemos buscar la recta $L_0$ tal que verifique: 

$
\left\{
\begin{array}{l} 
L_0 \paralela  L  \indicador{A}{(.1,.15)}{B}{(1.5,.15)} 
\\[5mm] 
A \in L_0 \marka{a}{(.1,.25)} \indicador{c}{(1.6,.15)}{d}{(2.3,.15)}
\\[5mm]
 d(L, L_0 ) =  4  \marka{b}{(.1,.15)} 
\end{array}
\right.
$




\tikz[overlay,remember picture] {
%%%%%%%%%%------------------------------------------------------------------------------------------------------------
\draw [overlay,very thick,blue,opacity=.5] (A)-|(c);  
\draw [overlay,->,very thick,blue,opacity=.5](b)-|(c)--(d);  
\draw [overlay,very thick,blue,opacity=.5](a)++(0,-.1)-|(c);
\node  at (d) [text width=131mm, below right, yshift=9mm]{
\begin{minipage}{.2\textwidth}
\begin{align*}
    d (L,L_0) &= 4
    \\[3mm]
    d(L,A)&=4
    \\[3mm]
    2&= 4 \quad \text{Abs!}
\end{align*}
\end{minipage}\qquad\qquad\begin{minipage}{.65\textwidth}
Cosas a tener en cuenta: 
\begin{itemize}
    \item $L$ y $L_0$ al ser paralelas la distancia entre ellas es la distancia entre un punto de una y la otro. 

    \item Ya sabemos que la distancia entre $L$ y $A$ es 2, es decir $d(L,A)=2$. 
\end{itemize}
\end{minipage} 
};
}
%%%%%%%%%%%%%%%%%%%%%%%%%%%%%%%%%%%%%%%%%%%%%%%%%%%%%%%%%%%%%%%%%%%%%%%%%%%%%%%%%%%%%%%%%%%%%%%%%%%%%%%%%%%%%%%%%%%%%%%%%%%%%%%%%%%%%%%%%%%%%%%%%%%%%%%%%%%%%%%%%%%%%%%%%%%%%%%%%%%%


\vspace{9mm}







No es posible encontrar una recta paralela a $L$ que pase por $A$ y que esté exactamente a $4$ unidades de distancia de $L$.

\begin{comment}

Recordemos que la ecuación implícita de la recta $L$ está dada por: \quad$ L:\  4x - 3y = 0$.

Las rectas paralelas a $L$ siguen la forma general:: \quad $L_0:\ 4x - 3y + C = 0$.

Dado que las rectas son paralelas, podemos utilizar la fórmula de distancia punto-recta. Consideremos el punto $A = (1,-2)$ perteneciente a una posible recta paralela $L_0$. Sabemos que la distancia entre $A$ y $L$ ya fue calculada y es igual a $2$:

\[
\frac{\left| 4(1) - 3(-2) \right|}{\| (4,-3) \|} = 2.
\]

Dado que la distancia deseada es $4$ unidades y la distancia entre $A$ y $L$ es menor, no es posible encontrar una recta paralela a $L$ que pase por $A$ y que esté exactamente a $4$ unidades de distancia de $L$.
    
\end{comment}

\bigskip  

\textbf{(d) Recta paralela a $L$ que pase por $A$ y que esté a 2 unidades de distancia de $L$.}

Siguiendo el mismo razonamiento que en el inciso anterior, observamos que la distancia entre $A$ y la recta $L$ es efectivamente $2$. Por lo tanto, la ecuación de la recta buscada es simplemente la ecuación de la recta paralela a $L$ que pasa por $A$. 

Su ecuación paramétrica es: \quad$R:\ (x, y) = (1,-2) + \lambda (3, 4),\quad \lambda \in \mathbb{R}.$

 \begin{comment}
\textbf{Gráfico}

\vspace{-9mm}
\begin{center}  
\definecolor{gris}{rgb}{0.5,0.5,0.5}
\begin{tikzpicture}[scale=.9]
% extremos para los ejes
\xdef\ext{4}
\xdef\yext{1}
\xdef\f{-2}
\xdef\g{-3}
\pgfmathtruncatemacro{\c}{\f+1}
\pgfmathtruncatemacro{\d}{\g+1}
\draw [color=gris!30!,dash pattern=on 1pt off 1pt, xstep=1cm,ystep=1cm] (\f-.3,\g-.3) grid (\ext+.3,\yext+.3);

\draw[->] (\f-.3,0) -- (\ext+.3,0) ;

\draw[->] (0,\g-.3) -- (0,\yext+.3) ;

\foreach \x in {\f,...,-1,1,2,...,\ext}
\draw[shift={(\x,0)},color=black] (0pt,2pt) -- (0pt,-2pt) node[below] {\tiny $\x$};

\foreach \y in {\g,...,-1,1}
\draw[shift={(0,\y)},color=black] (2pt,0pt) -- (-2pt,0pt) node[left] {\tiny $\y$};


    
    % Recta L 
    \draw[thick, blue, domain=-.8:.3,variable=\t] plot ({3*\t} , {4*\t}) node[right] {$L$};
    
    % Punto 
    \filldraw[red] (1,-2) circle (3pt) node[below right] {$A=(1,-2)$};
    
    
    % Rectas paralelas 
    \draw[thick, magenta, domain=-.35:.8,variable=\t] plot ({1+3*\t} , {-2+4*\t})  node[right] {Paralela a 2 unidades};
\end{tikzpicture}
\end{center}      
 \end{comment}

\bigskip  

\textbf{(e) Conclusiones}

Si bien es posible encontrar rectas paralelas a $L$ a diferentes distancias, la existencia de la recta que nos piden depende de la distancia del punto $A$ a $L$. Para que una recta paralela cumpla la condición de estar a cierta distancia y pase por un punto, la distancia entre el punto y la recta original debe coincidir exactamente con la distancia buscada.

\bigskip







































\item  En cada uno de los siguientes incisos, hallar, si es posible, la ecuación de la o las rectas que satisfacen las condiciones dadas:

\begin{itemize}
    \item[(a)] Tienen pendiente $-\frac{2}{3}$ y forman con los ejes coordenados un triángulo de área 24.
    \item[(b)] Son paralelas a la recta $L : 3x - 4y - 5 = 0$ y distan 4 unidades del punto $P = (-2, 1)$.
    \item[(c)] Pasan por el punto $P = (3, 5)$ y forman un ángulo de $30^\circ$ con la recta $L : 3x - y - 6 = 0$.
    \item[(d)] Contienen al punto $Q = (3,-2)$ y forman el mismo ángulo con $L_1 : \frac{x + 3}{2} = y + 2$ y $L_2 : 2x - y - 2 = 0$.
\end{itemize}







    \textbf{Resolución del ejercicio}

\begin{enumerate}[(a)]
\begin{minipage}{.6\textwidth}
    \item La ecuación de una recta con pendiente $m = -\frac{2}{3}$ es de la forma: \quad $ y = -\frac{2}{3}x + b_0. $ 

\setlength{\columnsep}{9mm}
\begin{itemize}
\begin{multicols}{2}
    \item intersección eje $x$: ($y=0$)
    \begin{align*}
        0 &= -\frac{2}{3}x + b_0
        \\
       \frac{2}{3}x  &= b_0
       \\
       x&= \frac{3}{2} b_0
    \end{align*}
    \columnbreak % Forzar cambio de columna si es necesario
 \item intersección eje $y$ es la ordenada al origen de la ecuación, es decir la intersección se da en $b$.  
\end{multicols}
\end{itemize}

\vspace{-5mm}

\hfill    Intersecta los ejes en $(0,b_0)$ y $(\frac{3b_0}{2},0)$. 
 
    
\end{minipage} \begin{minipage}{.3\textwidth}

    \begin{center}  
\definecolor{gris}{rgb}{0.5,0.5,0.5}
\begin{tikzpicture}[scale=.45]
% extremos para los ejes
\xdef\ext{6}
\xdef\yext{5}
\xdef\f{-2}
\xdef\g{-2}


\draw[->] (\f-.3,0) -- (\ext+.3,0) ;

\draw[->] (0,\g-.3) -- (0,\yext+.3) ;



\pgfmathsetmacro{\b}{3}
    
    % Rectas 
    \draw[thick, blue, domain=-2.3:6.6,variable=\t] plot ({\t} , {(-2*\t)/3+\b}) node[below left] {$y=-\frac{2}{3}x+b$}; 
    
    %%%% intercecciones con los ejes 
    \draw[red,fill=red] (0,\b) circle (3pt) node[below left]{$(0,b_0)$} ;
    \draw[red,fill=red] ({3*\b/2},0) circle (3pt) node[below left]{$(\frac{3b_0}{2},0)$} ;
\end{tikzpicture}
\end{center}
   
\end{minipage} 

\begin{itemize}
    \item 
    El área del triángulo se calcula con la fórmula: \quad $A = \frac{b \cdot h}{2}$, \ donde $ b $ es la base y $ h $ la altura.  

    \item A partir del gráfico, observamos que:  \quad $b = \left| \frac{3b}{2} \right|, \quad h = |b_0|$ 

\end{itemize}
Sustituyendo estos valores en la fórmula, obtenemos una ecuación  y resolviendo:   
    \begin{align*}
        A = 24 \quad \Longleftrightarrow \quad 
        \frac{b \cdot h}{2} = 24 
        \quad \Longleftrightarrow \quad
        \frac{ \left| \frac{3b}{2} \right| \cdot\left| b_0 \right|}{2} 
        = 24 
        \quad \Longleftrightarrow \quad
        \frac{ \left| 3b_0 \right|    \cdot\left| b_0 \right|}{4} 
        = 24
        \quad \Longleftrightarrow \quad
        \frac{ 3 \left| b_0 \right|^2 }{4} 
        &= 24
        \\ 
        b_0^2\ 
        &= \frac{4\cdot 24 }{3}
        \\ 
        b_0^2
        &= 32
        \\ 
        |b_0|
        &= \sqrt{32}
    \end{align*} 
    obtenemos $b = \pm \sqrt{32} \approx \pm 5,65685$.
    
    
    \vspace{9mm}
    
    
    Así, las ecuaciones son: \qquad $\displaystyle L_1: \quad y = -\frac{2}{3}x + \sqrt{32}, \quad \text{ y } \quad L_2: \quad \quad y = -\frac{2}{3}x - \sqrt{32}. $ 

\vspace{9mm}

    \begin{center}  
\definecolor{gris}{rgb}{0.5,0.5,0.5}
\begin{tikzpicture}[scale=.33]
% extremos para los ejes
\xdef\ext{10}
\xdef\yext{7}
\xdef\f{-10}
\xdef\g{-6}
\pgfmathtruncatemacro{\c}{\f+2}
\pgfmathtruncatemacro{\d}{\g+2}
\draw [color=gris!30!,dash pattern=on 1pt off 1pt, xstep=1cm,ystep=1cm] (\f-.3,\g-.3) grid (\ext+.3,\yext+.3);

\draw[->] (\f-.3,0) -- (\ext+.3,0) ;

\draw[->] (0,\g-.3) -- (0,\yext+.3) ;

\foreach \x in {\f,\c,...,-2,2,4,...,\ext}
\draw[shift={(\x,0)},color=black] (0pt,2pt) -- (0pt,-2pt) node[below] {\tiny $\x$};

\foreach \y in {\g,\d,...,-2,2,4,...,\yext }
\draw[shift={(0,\y)},color=black] (2pt,0pt) -- (-2pt,0pt) node[left] {\tiny $\y$};

\pgfmathsetmacro{\yy}{ sqrt(32) }
    
    % Rectas 
    \draw[thick, blue, domain=-2.3:10.3,variable=\t] plot ({\t} , {(-2*\t)/3+\yy}) node[below] {$y=-\frac{2}{3}x+\sqrt{32}$}; 
    
    \draw[thick, magenta, domain=-10.3:.8,variable=\t] plot ({\t} , {(-2*\t)/3-\yy})  node[right] {$y=-\frac{2}{3}x-\sqrt{32}$};
\end{tikzpicture}
\end{center}
    


 






\bigskip
    \item La ecuación de la recta $L$ es $3x - 4y - 5 = 0$,  por lo que las rectas paralelas son de la forma:
    \[ 3x - 4y + C = 0, \quad \text{con } \ C\in \R \]
    
    La distancia entre $A=(-2,1)$ y $R$ es 4 obeniendo una ecuación y la resolvemos:
    \begin{align*}
        d (A,R)= 4 \quad \Longleftrightarrow \quad   \frac{|3\cdot (-2) - 4 \cdot 1 + C|}{\|\vec{n}_R \|} &= 4  \qquad\qquad\qquad
        \\
        \frac{|-6 - 4 + C|}{\sqrt{3^2 + (-4)^2}} 
        &= 4
        \\
        \frac{|-10 + C|}{\sqrt{9 + 16\, }\, } 
        &= 4
        \\
        \frac{|-10 + C|}{\sqrt{25\, }\, } 
        &= 4
        \\
        \frac{|-10+C|}{5\, } 
        &= 4 
        \\
         |-10+ C|  
        &= 5 \cdot 4  
        \\
         |-10 + C|  
        &= 20 
\end{align*} 
%%%%%%%%%%%%%%%%%%%%%%%%%%%%%%%%%%%%%%%%%%%%%%%%%%%%%%%%%%%%%%%%%%%%%%%%%%%%%%%%%%%%%%%%%%%%%%%%%% 
\[
\begin{array}{rl@{\hspace{20mm}}c @{\hspace{20mm}} rl r}   
-10 + C
\qquad
&= \qquad
20
& \mbox{ o }
&
-10  + C
\qquad
&= \qquad
-20 
\\[3mm] %%%%%%%%%%%%%%%% %%%%%%%%%%%%%%%% %%%%%%%%%%%%%%%% %%%%%%%%%%%%%%%% %%%%%%%%%%%%%%%% %%%%%%%%%%%%%%%%  
C
&= 
20-10
& %\mbox{ o }
&
C 
&= 
-20-10
\\[3mm] %%%%%%%%%%%%%%%% %%%%%%%%%%%%%%%% %%%%%%%%%%%%%%%% %%%%%%%%%%%%%%%% %%%%%%%%%%%%%%%% %%%%%%%%%%%%%%%%  
C
&= 
10
& % \mbox{ o }
&
C  
&= 
-30
\end{array}
\]
Por lo que las ecuaciones son:\qquad 
$ R_1: \quad 3x - 4y +30 = 0, \quad\qquad R_2:\quad 3x - 4y -10 = 0.$

    





  \begin{center}  
\definecolor{gris}{rgb}{0.5,0.5,0.5}
\begin{tikzpicture}[scale=.33]
% extremos para los ejes
\xdef\ext{10}
\xdef\yext{7}
\xdef\f{-10}
\xdef\g{-6}
\pgfmathtruncatemacro{\c}{\f+2}
\pgfmathtruncatemacro{\d}{\g+2}
\draw [color=gris!30!,dash pattern=on 1pt off 1pt, xstep=1cm,ystep=1cm] (\f-.3,\g-.3) grid (\ext+.3,\yext+.3);

\draw[->] (\f-.3,0) -- (\ext+.3,0) ;

\draw[->] (0,\g-.3) -- (0,\yext+.3) ;

\foreach \x in {\f,\c,...,-2,2,4,...,\ext}
\draw[shift={(\x,0)},color=black] (0pt,2pt) -- (0pt,-2pt) node[below] {\tiny $\x$};

\foreach \y in {\g,\d,...,-2,2,4,...,\yext }
\draw[shift={(0,\y)},color=black] (2pt,0pt) -- (-2pt,0pt) node[left] {\tiny $\y$};


\coordinate (A) at (-2,1); 
\draw[blue,fill=blue] (A) circle (2pt) node[below]{$A$} ;
 
 
    
    % Rectas 
    \draw[thick, violet, domain=-10.3:.65,variable=\t] plot ({\t}, {(3*\t+30)/4}) node[above right] {$R_1:\  3x - 4y +30 = 0$}; 

     \draw[thick, blue, domain=-6.9:10.3,variable=\t] plot ({\t}, {(3*\t-5)/4}) node[ right] {$L:\  3x - 4y -5 = 0$}; 
    
    \draw[thick, magenta, domain=-5.7:10.3,variable=\t] plot ({\t},{(3*\t-10)/4})  node[right] {$R_2:\ 3x - 4y -10 = 0.$};

    
\end{tikzpicture}
\end{center}
    






















    

    \item  Pasan por el punto $P = (3, 5)$ y forman un ángulo de $30^\circ$ con la recta $L : 3x - y - 6 = 0$.

    Consideremos que $R$ es la recta que vamos a buscar.  Para resolver, es conveniente saber si el punto dado pertenece o no a la recta.
    tenemos: $$3\cdot3  - 5 - 6 = -2 \neq 0$$
    Al ser rectas en el plano y $\text{áng}(L, R) = 30^\circ$ se concluye que las rectas se intersectan en un punto, al punto en la intersección lo llamaremos $A$. Entonces tenemos que $\overrightarrow{PA}$ es el director. Veamos esto en un gráfico: 




 \begin{center}  
\definecolor{gris}{rgb}{0.5,0.5,0.5}
\begin{tikzpicture}[scale=2.5]
    
    % Recta
     \draw[thick, blue, domain=3.1:4.2,variable=\t] plot ({\t}, {(3*\t-6)/1}) node[above,xshift=-9mm ] {$L:\  3x - y -6 = 0$};
     
     
     \coordinate (P) at (3,5);%punto 
     
        \pgfmathsetmacro{\xx}{(18-sqrt(3))/5}
        \pgfmathsetmacro{\yy}{3*\xx-6}
        \pgfmathsetmacro{\xxx}{(18+sqrt(3))/5}
        \pgfmathsetmacro{\yyy}{3*\xxx-6}
     \coordinate (A) at (\xx,\yy);
     \coordinate (AA) at (\xxx,\yyy);


      %rexta que pasa por P y A
\coordinate (m) at ($2.2*(P) -1.2*(A)$);%Punto por detras de p
\coordinate (M) at ($-.4*(P)+1.4*(A)$);%Punto por adelante  de A
\draw[dash pattern=on 1pt off 1pt, cyan] (M) -- (m) node[left ] {$R_1$};
    %
\coordinate (m) at ($3.1*(P) -2.1*(AA)$);%Punto por detras de P
\coordinate (M) at ($-.9*(P)+1.9*(AA)$);%Punto por adelante  de B
\draw[dash pattern=on 1pt off 1pt, cyan] (m) -- (M) node[right] {$R_2$};


    %vectores directores
    \draw[->,thick,violet] (P)--(A);
    \draw[->,thick,violet] (P)--(AA);


    \pic[draw, color=violet, fill=violet!50!,   angle eccentricity=1.4, angle radius=1.2cm, "$ \theta= 30^\circ$",opacity=.8] {angle = P--AA--A};

     \pic[draw, color=violet, fill=violet!50!,   angle eccentricity=1.4, angle radius=1.2cm, "$ \theta= 30^\circ$",opacity=.8] {angle = AA--A--P};

     

    
    \fill[red] (P) circle (2pt) node[below left] {$P=(3,5)$};
    
\end{tikzpicture}
\end{center}







    \textcolor{red}{\textbf{-}} \textbf{El punto} $ A = (x_A, y_A) $ es punto de la recta $ L $ $\Rightarrow$ Debe verificar su ecuación $ 3x_A -y_A - 6 = 0 $

\[
3x_A -y_A - 6 = 0 
\quad 
\Longrightarrow
\quad
3x_A  - 6 =  y_A   
\qquad
\qquad
\therefore A = \left( x_A, 3x_A  - 6  \right)
\]

\noindent
- Calculemos  $ \overrightarrow{PA}$:

\begin{align*}
    \overrightarrow{PA} = A - P 
    &= \left( x_A - 3, 3x_A  - 6 -5\right)
    \\
     &= \left( x_A - 3, 3x_A  - 11\right) &\text{Solo depende de una incógnita $ x_A $.}
\end{align*} 

\noindent
- Tomamos la fórmula de ángulo entre rectas y reemplazamos los datos:

\[
\text{áng}(R, L) = 30^\circ 
\quad \Rightarrow \quad 
30^\circ = \arccos \left( \frac{|\overrightarrow{u_{R}} \cdot \overrightarrow{u_L}|}{\|\overrightarrow{u_{R}}\| \cdot \|\overrightarrow{u_L}\|} \right)
\]

\begin{align*}
    30^\circ \qquad &= \arccos \left( \frac{|(x_A - 3, 3x_A  -11) \cdot (1,3)|}{\|(x_A - 3, 3x_A  - 11)\| \cdot \|(1,3)\|} \right) 
    \\[2mm] 
    \cos 30^\circ \qquad &= \frac{| x_A -3 + 3(3x_A  - 11) |}{\sqrt{(x_A - 3)^2 + \left( 3x_A  - 11\right)^2} \cdot \sqrt{1^2 + 3^2}} 
    \\[2mm]  
    \tfrac{\sqrt{3}}{2} \qquad &= \frac{| x_A -3 + 9x_A  - 33 |}{\sqrt{(x_A - 3)^2 + \left( 3x_A  - 11\right)^2} \cdot \sqrt{1 + 9\, }} 
    \\[2mm]  
    \tfrac{\sqrt{3}}{2} \qquad &= \frac{| 10x_A  -36 |}{\sqrt{10\, x_{A}\!\! \rule{0mm}{3.5mm}^2 -72 x_A +130\, }\cdot \sqrt{10}} 
    \\[2mm]      
    \tfrac{\sqrt{3}}{2} \sqrt{10\, x_{A}\!\! \rule{0mm}{3.5mm}^2 -72 x_A +130\,  }\cdot \sqrt{10}&= \qquad  | 10x_A  -36 |  
    \\[2mm]     
   \left( \tfrac{\sqrt{30}}{2} \sqrt{10\, x_{A}\!\! \rule{0mm}{3.5mm}^2 -72 x_A +130\,  }  \right) ^2 &= \qquad  | 10x_A  -36 | ^2
    \\[2mm]    
    \tfrac{30}{4} \left( 10\, x_{A}\!\! \rule{0mm}{3.5mm}^2 -72 x_A +130 \right) &= \qquad  100 x_{A}\!\! \rule{0mm}{3.5mm}^2 -720x_A + 1296 
    \\[2mm]     
   \left( \tfrac{\sqrt{30}}{2} \sqrt{10\, x_{A}\!\! \rule{0mm}{3.5mm}^2 -72 x_A +130\,  }  \right) ^2 &= \qquad  | 10x_A  -36 | ^2
    \\[2mm]    
    75\, x_{A}\!\! \rule{0mm}{3.5mm}^2 - 540 x_A + 975 &= \qquad  100 x_{A}\!\! \rule{0mm}{3.5mm}^2 -720x_A + 1296 
    \\[2mm]    
    0  &= \qquad  100 x_{A}\!\! \rule{0mm}{3.5mm}^2-75  x_{A}\!\! \rule{0mm}{3.5mm}^2 -720x_A  +  540 x_A  + 1296 -975
    \\[2mm]    
    0  &= \qquad  25 x_{A}\!\! \rule{0mm}{3.5mm}^2 -180 \, x_{A}   + 321
\end{align*}

\[
x_A = \frac{-(-180) \pm \sqrt{(-180)^2 - 4 \cdot 25 \cdot321 \, } }{2 \cdot 25 }
= \frac{180\pm 10\sqrt{3}}{50} = \frac{18\pm  \sqrt{3}}{5}\]
\[   
\Longrightarrow    
\qquad   
x_A = \frac{18 -  \sqrt{3}}{5} \qquad  \lor \qquad x_A = \frac{18 +  \sqrt{3}}{5}
\] 




\begin{itemize}
    \item Si $x_A=\frac{18 -  \sqrt{3}}{5}    $ tenemos que $y_A=  3\cdot\frac{18 -  \sqrt{3}}{5}  -6 = \frac{24 - 3 \sqrt{3}}{5}   $, así     $A=\left( \frac{18 -  \sqrt{3}}{5}   ,\, \frac{24 - 3 \sqrt{3}}{5}\right)$  y luego:




\[
\overrightarrow{u_{R_1}} =  \overrightarrow{PA} 
=
\left( \frac{18 -  \sqrt{3}}{5} -3  ,\, \frac{24 - 3 \sqrt{3}}{5} -5 \right)
= 
\left( \frac{3 -  \sqrt{3}}{5}   ,\,- \frac{1 + 3 \sqrt{3}}{5}  \right)
\]

Se puede elegir como vector director a:
\[
\overrightarrow{u_{R_1}} =  \left( \frac{3 -  \sqrt{3}}{5}   ,\,- \frac{1 + 3 \sqrt{3}}{5}  \right) \quad \text{ó cualquier vector paralelo a él.}
\]

Por ejemplo:
\[
5 \cdot \overrightarrow{u_{R_1}} = \left( 3 -  \sqrt{3}   ,\,-  1 -3 \sqrt{3}   \right)=\overrightarrow{d_{R_1}}
\]
y la ecuación de la recta nos queda: 

\[
\textcolor{red}{R_1: (x,y) = (3,5) + \lambda\left( 3 -  \sqrt{3}   ,\,-  1 -3 \sqrt{3}   \right), \quad \lambda \in \mathbb{R}}
\]


    

    \item Si $x_A=  \frac{18 + \sqrt{3}}{5}  $ tenemos que $y_A= 3\cdot  \frac{18 + \sqrt{3}}{5}  -6 = \frac{24 + 3 \sqrt{3}}{5}  $ así     $A=\left( \frac{18 +  \sqrt{3}}{5}   ,\, \frac{24 + 3 \sqrt{3}}{5}\right)$  y luego:




\[
\overrightarrow{u_{R_1}} =  \overrightarrow{PA} 
=
\left( \frac{18 + \sqrt{3}}{5} -3  ,\, \frac{24 + 3 \sqrt{3}}{5} -5 \right)
= 
\left( \frac{3 +  \sqrt{3}}{5}   ,\, \frac{-1 + 3 \sqrt{3}}{5}  \right)
\]

Se deja como vector director a:
\[
\overrightarrow{u_{R_1}} =  \left( \frac{3 +  \sqrt{3}}{5}   ,\, \frac{ -1 + 3 \sqrt{3}}{5}  \right) \quad \text{ó cualquier vector paralelo a él.}
\]

Por ejemplo:
\[
5 \cdot \overrightarrow{u_{R_1}} = \left( 3 +  \sqrt{3}   ,\,-  1 +3 \sqrt{3}   \right)=\overrightarrow{d_{R_1}}
\]
y la ecuación de la recta nos queda: 

\[
\textcolor{red}{R_2: (x,y) = (3,5) + \lambda\left( 3 +  \sqrt{3}   ,\,-  1 +3 \sqrt{3}   \right), \quad \lambda \in \mathbb{R}}
\]



\end{itemize}
















 \begin{center}  
\definecolor{gris}{rgb}{0.5,0.5,0.5}
\begin{tikzpicture}[scale=1.5]
% extremos para los ejes
\xdef\ext{5}
\xdef\yext{7}
\xdef\f{-0}
\xdef\g{-0}
\pgfmathtruncatemacro{\c}{\f+1}
\pgfmathtruncatemacro{\d}{\g+1}
\draw [color=gris!30!,dash pattern=on 1pt off 1pt, xstep=1cm,ystep=1cm] (\f-.3,\g-.3) grid (\ext+.3,\yext+.3);

\draw[->] (\f-.3,0) -- (\ext+.3,0) ;

\draw[->] (0,\g-.3) -- (0,\yext+.3) ;

\foreach \x in {\f,\c,...,-1,1,2,...,\ext}
\draw[shift={(\x,0)},color=black] (0pt,2pt) -- (0pt,-2pt) node[below] {\tiny $\x$};

\foreach \y in {\g,\d,...,-1,1,2,...,\yext }
\draw[shift={(0,\y)},color=black] (2pt,0pt) -- (-2pt,0pt) node[left] {\tiny $\y$};
 
    
    % Recta
     \draw[thick, blue, domain=1.9:4.3,variable=\t] plot ({\t}, {(3*\t-6)/1}) node[above left] {$L:\  3x - y -6 = 0$};
     
     
     \coordinate (P) at (3,5);%punto 
     
        \pgfmathsetmacro{\xx}{(18-sqrt(3))/5}
        \pgfmathsetmacro{\yy}{3*\xx-6}
        \pgfmathsetmacro{\xxx}{(18+sqrt(3))/5}
        \pgfmathsetmacro{\yyy}{3*\xxx-6}
     \coordinate (A) at (\xx,\yy);
     \coordinate (AA) at (\xxx,\yyy);


      %rexta que pasa por P y A
\coordinate (m) at ($2.8*(P) -1.8*(A)$);%Punto por detras de A
\coordinate (M) at ($-3.2*(P)+4.2*(A)$);%Punto por adelante  de B
\draw[dash pattern=on 1pt off 1pt, cyan] (M) -- (m) node[right] {$R_1$};
    %
\coordinate (m) at ($3.8*(P) -2.8*(AA)$);%Punto por detras de A
\coordinate (M) at ($-1.3*(P)+2.3*(AA)$);%Punto por adelante  de B
\draw[dash pattern=on 1pt off 1pt, cyan] (m) -- (M) node[right] {$R_2$};


    %vectores directores
    \draw[->,thick,violet] (P)--(A);
    \draw[->,thick,violet] (P)--(AA);


    \pic[draw, color=violet, fill=violet!50!,   angle eccentricity=1.4, angle radius=.8cm, "$ \theta= 30^\circ$",opacity=.8] {angle = P--AA--A};

     \pic[draw, color=violet, fill=violet!50!,   angle eccentricity=1.4, angle radius=.8cm, "$ \theta= 30^\circ$",opacity=.8] {angle = AA--A--P};

     

    
    \fill[red] (P) circle (2pt) node[below left] {$P=(3,5)$};
    
\end{tikzpicture}
\end{center}


    
\begin{comment}
Para analizar la situación, observemos el siguiente gráfico:
    
    
\begin{multicols}{2}
  \begin{center}  
\definecolor{gris}{rgb}{0.5,0.5,0.5}
\begin{tikzpicture}[scale=1.2]
% extremos para los ejes
\xdef\ext{5}
\xdef\yext{7}
\xdef\f{-0}
\xdef\g{-0}
\pgfmathtruncatemacro{\c}{\f+1}
\pgfmathtruncatemacro{\d}{\g+1}
\draw [color=gris!30!,dash pattern=on 1pt off 1pt, xstep=1cm,ystep=1cm] (\f-.3,\g-.3) grid (\ext+.3,\yext+.3);

\draw[->] (\f-.3,0) -- (\ext+.3,0) ;

\draw[->] (0,\g-.3) -- (0,\yext+.3) ;

\foreach \x in {\f,\c,...,-1,1,2,...,\ext}
\draw[shift={(\x,0)},color=black] (0pt,2pt) -- (0pt,-2pt) node[below] {\tiny $\x$};

\foreach \y in {\g,\d,...,-1,1,2,...,\yext }
\draw[shift={(0,\y)},color=black] (2pt,0pt) -- (-2pt,0pt) node[left] {\tiny $\y$};
 
    
    % Recta
     \draw[thick, blue, domain=1.9:4.33,variable=\t] plot ({\t}, {(3*\t-6)/1}) node[ left] {$L:\  3x - y -6 = 0$};
     
     
     \coordinate (P) at (3,5);%punto 
     
        
    \foreach \angle/\t in {0/.9,30/1.1,60/.8,90/1.3,120/1,150/.9,180/1.,210/.9 ,270/1.4,240/1.2 270/.9,300/1.3,330/1.1} {
        \draw[->,thick,violet] (P) -- ++({\t*cos(\angle)}, {\t*sin(\angle)});    }
    \fill[red] (P) circle (2pt) node[below left] {$P=(3,5)$};
    %punto en la rect L y vector normal
    \coordinate (A) at (8/3,2);
    \draw[->,thick,magenta] (A) -- ++( 3,-1) node[pos=.4,above right] {$\vec{n}_{L}=(3,-1) $}; 
\end{tikzpicture}
\end{center}

    \columnbreak % Forzar cambio de columna si es necesario


Para determinar la ecuación de la(s) recta(s), es necesario encontrar el vector normal de la recta $ R $, el cual tendrá la forma genérica $ \vec{n}_{R} = (a,b) $.  

Sabemos que la recta $ R $ pasa por el punto $ (3,5) $ y que forma un ángulo de $ 30^\circ $ con la recta $ L $, cuyo vector normal es $ \vec{n}_{L} = (3,-1) $. En el gráfico, hemos representado en violeta las posibles direcciones de este vector.  

El vector normal $ \vec{n}_{R} = (a,b) $ debe formar un angulo de $30^\circ$  con $ \vec{n}_{L} $, es decir:  

\[
\cos(30^\circ) = \frac{\vec{n}_L \cdot \vec{n}_R}{\|\vec{n}_L\| \cdot \|\vec{n}_R\|}
\]

Además, como solo nos importa la dirección del normal se impone que el vector $ \vec{n}_{R} $ tenga norma 1, es decir:  

\[
\|(a,b)\|=1
\]

Con estas condiciones, podemos resolver el problema y determinar la ecuación de la(s) recta(s) buscada(s).  



\end{multicols}    


\[
 \begin{array}{rl @{\hspace{59mm} } rl}
        \bullet \qquad \qquad 
        \cos(30^\circ) 
        &= \dfrac{ \vec{n}_L \cdot \vec{n}_{L_0} }{\|\vec{n}_L\| \cdot \|\vec{n}_{L_0}\|} 
        &
        \bullet  \qquad
        \|(a,b)\|&=1  \qquad
        \\[4mm]%--------------------------%--------------------------%--------------------------%--------------------------%--------------------------%--------------------------
        \tfrac{\sqrt{3}}{2}
        &= \dfrac{ (3,-1)\cdot (a,b) }{\|(3,-1) \| \cdot \underbrace{\|(a,b) \|}_{=1}}
        &
        \sqrt{a^2 + a^2\,} &=  1
        \\[-4mm]%--------------------------%--------------------------%--------------------------%--------------------------%--------------------------%--------------------------
        &
        & {a^2 + b^2\,} &=\marka{A}{(-.3,-.15)}  1
        \\[4mm]%--------------------------%--------------------------%--------------------------%--------------------------%--------------------------%--------------------------
        \tfrac{\sqrt{3}}{2}
        &= \dfrac{  3a -b }{\ \sqrt{3^2 + (-1)^2\, } \ }
        &  
        \\[4mm]%--------------------------%--------------------------%--------------------------%--------------------------%--------------------------%--------------------------
        \tfrac{\sqrt{3}}{2}
        &= \dfrac{ | 3a -b|}{\sqrt{10\, }  }
        & &
        \\[4mm]%--------------------------%--------------------------%--------------------------%--------------------------%--------------------------%--------------------------
        \tfrac{\sqrt{3}}{2} \sqrt{10\, }
        &=  3a - b 
        & & 
        \\[4mm]%--------------------------%--------------------------%--------------------------%--------------------------%--------------------------%--------------------------
        \tfrac{\sqrt{30}}{2} 
        &=  3 a - b 
        & &
        \\[4mm]%--------------------------%--------------------------%--------------------------%--------------------------%--------------------------%--------------------------
        b\ \ \ 
        &=  3 a - \tfrac{\sqrt{30}}{2} \indicador{B}{(.1,.1)}{C}{(1.5,2.5)} \marka{D}{(2.1,2.5)}
        & &
 \end{array}
\]

\tikz[overlay,remember picture] {
\draw [overlay,very thick,red,opacity=.5](B) to[out= 0, in= -120 ] (C) to[out= 60, in=-120](A);
\draw [overlay,->,very thick,red,opacity=.5] (C)--(D);
\node  at (D) [text width=75mm, below right , yshift= 7mm ]{
Sustituyendo $b= 3 a - \tfrac{\sqrt{30}}{2} $ en la ecuación de la norma: 
\begin{align*} 
a^2 + \left(3a - \frac{\sqrt{30}}{2} \right)^2 &= 1
\\
a^2 + 9a^2 - 3a \sqrt{30} + \frac{30}{4}&= 1
\\
10a^2 - 3a \sqrt{30} + \frac{30}{4} &= 1
\\
40a^2 - 12a \sqrt{30} + 30 &= 4
\\
40  \indicador{a}{(-.2,-.1)}{a1}{(-.4,-.4)}  a^2 - 12 \sqrt{30} \indicador{b}{(-.5,-.1)}{b1}{(-.5,-.4)} a  + 26\indicador{c}{(-.1,-.1)}{c1}{(.2,-.4)}  &= 0   
\end{align*}};
}



%%%
%%
%%
%%
%%
%%
%%
%%
%%%
\begin{tikzpicture}[overlay,remember picture] 
\node at (a1) [below, blue,xshift=1.5mm]{ 
$a= 40$
};
\draw [overlay,<-,very thick,blue,opacity=.5]
(a) --(a1);  %%%%%%%%%%%%%%%%%%%%%%%%%%%%%%%%%%%%%%%%%%%%%%%%%%%%%%%%%%%%%%%%%%%%%%%%%%%%%%%%%%%%%%%%%%%%%%%%%%%%%%%%%%%%%%%%%%%%%%%%%%%%%%%%%%%%%%%%%%%%%%%%
\node at (b1) [below, blue,xshift=1.5mm]{ 
$b= -12\sqrt{30}$
};
\draw [overlay,<-,very thick,blue,opacity=.5]
(b) --(b1); 
%%%%%%%%%%%%%%%%%%%%%%%%%%%%%%%%%%%%%%%%%%%%%%%%%%%%%%%%%%%%%%%%%%%%%%%%%%%%%%%%%%%%%%%%%%%%%%%%%%%%%%%%%%%%%%%%%%%%%%%%%%%%%%%%%%%%%%%%%%%%%%%%
\node at (c1) [below, blue,xshift=1.5mm]{ 
$c= 26$
};
\draw [overlay,<-,very thick,blue,opacity=.5]
(c) --(c1); 
\end{tikzpicture}





\vspace{24mm}



Aplicando la fórmula cuadrática o Baskara , tenemos:

\[a_1,\,a_2 = \frac{{-(-12\sqrt{30}) \pm \sqrt{{(-12\sqrt{30})^2 - 4 \cdot 40\cdot 26}}}}{{2 \cdot 40 }}\]


 y nos queda:

$ a_{1,\,2}
=
\frac{12\sqrt{30} \pm \sqrt{160}}{80 }
 =
\frac{3\sqrt{30} \pm \sqrt{10}}{20 } 
 =
\frac{3\sqrt{3} \sqrt{10} \pm \sqrt{10}}{20 }  =
\frac{3\sqrt{3} \pm 1}{20 } \sqrt{10} \marka{A}{(.2,.15)}\qquad
      \qquad
    \begin{array}{l}
    \marka{A1}{(-.1,.15)}  \frac{3\sqrt{3} + 1}{20 } \sqrt{10}  \\[6mm]
     \marka{A2}{(-.1,.15)}\frac{3\sqrt{3} - 1}{20 } \sqrt{10}
    \end{array} $


\tikz[overlay,remember picture] { 
\draw[thick,blue, ->,line width=1pt] (A)--(A1) ;
\draw[thick,blue, ->,line width=1pt] (A)--(A2) ; 
%%%%%%%%%%%%%%%%%%%%%%%%%%%%%%%%%%%%%%%%%%%%%%%%%%%%%%%%%%%%%%%%%%%%%%%%%%%%%%%%%%%%%%%%%%%%%%%%%%
}
\bigskip



\begin{itemize}
    \item Si $a= \frac{3\sqrt{3} + 1}{20 } \sqrt{10} $ tenemos que $b= \sqrt{1- \left( \frac{3\sqrt{3} + 1}{20 } \sqrt{10}\right)^2} = \frac{\sqrt{30-15\sqrt{3}}}{10}$

    \item Si $a=\frac{3\sqrt{3} - 1}{20 } \sqrt{10}$ tenemos que $b= \sqrt{1- \left( \frac{3\sqrt{3} - 1}{20 } \sqrt{10}\right)^2} = \frac{\sqrt{30+15\sqrt{3}}}{10} $
\end{itemize}


&%%%otra%%forma%%%%%%%%%%%%%%%%%%%%%%%%%%%%%%%%%%%%%%%%%%%%%%%%%%%%%%%%%%%%%%%%%%%%%%%%%%%%%%%%%%%%%%%%%%%%%%%%%%%%%%%%%%%%%%%%%%%%%%%%%%%%%%%%%%%%%%%%%%%%%%%%%%%%%%%%%%%%%%%%%%%%%%%%%%%%%%%%%%%%%%%%%%%%%%%%%%%%%%%%%%%%%%%%%%%%%%%%%%%%%%%%%%%%%%%%%%%%%%%%%%%%%%%%%%%%%%%%%%%%%%%%%%%%%%%%%%%%%%%%%%%%%%%%%%%%%%%%%%%%%%%%%%%%%%%%%%%%%%%%%%%%%%%%%%%%%%%%%%%%%%%%%%%%%%%%%%%%%%%%%%%%%%%%%%%%%%%%%%%%%%%%%%%%%%%%%%%%%%%%%%%%%%%%%%%%%%%%%%%%%%%%%%%%%%%%%%%%%%%%%%%%%%%%%%%%%%%%%%%%%%%%%%%%%%%%%%%%%%%%%%%%%%%%%%%%%%%%%%%%%%%%%%%%%%%%%%%%%%%%%%%%%%%%%%%%%%%%%%%%%%%%%%%%%%%%%%%%%%%%%%%%%%%%%%%%%%%%%%%%%%%%%%%%%%%%%%%%%%%%%%%%%%%%%%%%%%%%%%%%%%%%%%%%%%%%%%%%%%%%%%%%%%%%%%%%%%%%%%%%%%%%%%%%%%%%%%%%%%%%%%%%%%%%%%%%%%%%%%%%%%%%%%%%%%%%%%%%%%%%%%%%%%%%%%%%%%%%%%%%%%%%%%%%%%%%%%%%%%%%%%%%%%%%%%%%%%%%%%%%%%%%%%%%%%%%%%%%%%%%%%%%%%%%%%%%%%%%%%%%%%%%%%%%%%%%%%%%%%%%%%%%%%%%%%%%%%%%%%%%%%%%%%%%%%%%%%%%-%-%-%-%-%-%-%-%-%-%-%%-%--%-%-%-%-%-%-%-%-%-%-%-%-%-%-%%-%--%-%-%-%-%-%-%-%-%-%-%-%-%-%-%%-%--%-%-%-%-%-%-%-%-%-%-%-%-%-%-%%-%--%-%-%-%-%-%-%-%-%-%-%-%-%-%-%%-%--%-%-%-%-%-%-%-%-%-%-%-%-%-%-%%-%--%-%-%-%-%-%-%-%-%-%-%-%-%-%-%%-%--%-%-%-%-%-%-%-%-%-%-%-%-%-%-%%-%--%-%-%-%-%-%-%-%-%-%-%-%-%-%-%%-%--%-%-%-%-%-%-%-%-%-%-%-%-%-%-%%-%--%-%-%-%-%-%-%-%-%-%-%-%-%-%-%%-%--%-%-%-%-%-%-%-%-%-%-%-%-%-%-%%-%--%-%-%-%-%-%-%-%-%-%-%-%-%-%-%%-%--%-%-%-%-%-%-%-%-%-%-%-%-%-%-%%-%--%-%-%-%-%-%-%-%-%-%-%-%-%-%-%%-%--%-%-%-%-%-%-%-%-%-%-%-%-%-%-%%-%--%-%-%-%-%-%-%-%-%-%-%-%-%-%-%%-%--%-%-%-%-%-%-%-%-%-%-%-%-%-%-%%-%--%-%-%-%-%-%-%-%-%-%-%-%-%-%-%%-%--%-%-%-%-%-%-%-%-%-%-%-%-%-%-%%-%--%-%-%-%-%-%-%-%-%-%-%-%-%-%-%%-%--%-%-%-
    {\bf Recordemos:} 

    El ángulo entre dos rectas en el plano se puede calcular con alguna de las siguientes fórmulas:

\begin{multicols}{2}
\begin{tcolorbox}[colframe=blue, colback=blue!10, rounded corners=southwest]
\[
\theta = \text{áng}(L,R) = \arccos \left( \frac{|\vec{d_L} \cdot \vec{d_R}|}{\|\vec{d_L}\| \cdot \|\vec{d_R}\|} \right)
\]
\end{tcolorbox}
\columnbreak %salto de columna
\begin{tcolorbox}[colframe=blue, colback=blue!10, rounded corners=southwest]
\[
\theta = \text{áng}(L,R) = \arccos \left( \frac{|\vec{n_L} \cdot \vec{n_R}|}{\|\vec{n_L}\| \cdot \|\vec{n_R}\|} \right)
\]
\end{tcolorbox}
\end{multicols}
    
Consideremos la recta $L_0$ con ecuación general:
    \[ Ax + By + C = 0. \]
    
    El ángulo entre las rectas $L$ y $L_0$ se calcula usando sus vectores normales:
    \[
    \cos(30^\circ) = \frac{|\vec{n}_L \cdot \vec{n}_{L_0}|}{\|\vec{n}_L\| \cdot \|\vec{n}_{L_0}\|},
    \]
    donde $\vec{n}_L = (3,-1)$ y $\vec{n}_{L_0} = (A,B)$ y suponemos que $\vec{n}_{L_0} $ es de norma 1. despegando  $A$ y $B$:
\[
 \begin{array}{rl @{\hspace{30mm} } rl}
\cos(30^\circ) 
        &= \dfrac{|\vec{n}_L \cdot \vec{n}_{L_0}|}{\|\vec{n}_L\| \cdot \|\vec{n}_{L_0}\|} 
        &
        \|(A,B)\|&=1 
        \\[4mm]%--------------------------%--------------------------%--------------------------%--------------------------%--------------------------%--------------------------
        \tfrac{\sqrt{3}}{2}
        &= \dfrac{ |(3,-1)\cdot (A,B) |}{\|(3,-1) \| \cdot \underbrace{\|(A,B) \|}_{=1}}
        &
        \sqrt{A^2 +B^2\,} &=1
        \\[-4mm]%--------------------------%--------------------------%--------------------------%--------------------------%--------------------------%--------------------------
        &
        & {A^2 +B^2\,} &=1
        \\[4mm]%--------------------------%--------------------------%--------------------------%--------------------------%--------------------------%--------------------------
        \tfrac{\sqrt{3}}{2}
        &= \dfrac{ | 3A -B |}{\ \sqrt{3^2 + (-1)^2\, } \ }
        &  
        \\[4mm]%--------------------------%--------------------------%--------------------------%--------------------------%--------------------------%--------------------------
        \tfrac{\sqrt{3}}{2}
        &= \dfrac{ | 3A -B |}{\sqrt{10\, }  }
        & &
        \\[4mm]%--------------------------%--------------------------%--------------------------%--------------------------%--------------------------%--------------------------
        \tfrac{\sqrt{3}}{2} \sqrt{10\, }
        &= | 3A -B | 
        & & 
        \\[4mm]%--------------------------%--------------------------%--------------------------%--------------------------%--------------------------%--------------------------
        \tfrac{\sqrt{30}}{2} 
        &= | 3A -B | 
        & &
 \end{array}
\]
$
\qquad \begin{array}{rl @{\qquad } c @{\qquad } rl}
 3A -B = \tfrac{\sqrt{30}}{2}
 &\qquad\text{ o}&
 3A -B = -\tfrac{\sqrt{30}}{2}
 \\[3mm]
 3A - \tfrac{\sqrt{30}}{2}= B
 & &
 3A +\tfrac{\sqrt{30}}{2} = B
 \\
 \end{array}
$   

\begin{itemize}
    \item  \begin{align*}
        A^2 + \left( 3A - \tfrac{\sqrt{30}}{2}\right)^2 = 1
        \end{align*}
     \item  \begin{align*}
        A^2 + \left( 3A + \tfrac{\sqrt{30}}{2}\right)^2 = 1
        \end{align*}
\end{itemize}
    
    Despejamos $A$ y $B$ . Luego, como el punto $(3,5)$ pertenece a $L_0$, sustituimos en la ecuación para encontrar $C$.

    Así obtenemos las ecuaciones buscadas.

\end{comment}




































\item  Hallar, si es posible, la ecuación de la o las rectas que contienen al punto $Q = (3,-2)$ y forman el mismo ángulo con $L_1 : \frac{x + 3}{2} = y + 2$ y $L_2 : 2x - y - 2 = 0$.
























\textbf{Solución:}

Consideremos una recta genérica que pasa por $Q$ y tiene vector director $(a, b)$:
\[ \text{(Ecuación vectorial) } \quad R :\quad  (x, y) = (3, -2) + \lambda (a, b), \quad \lambda \in \mathbb{R}. \]

El vector normal  es $\vec{n}_{R} =(-b, a)$. Utilizamos las fórmulas para calcular los ángulos entre rectas:
\[ \theta = \arccos\left( \frac{|\vec{d}_L \cdot \vec{d}_R|}{\|\vec{d}_L\| \|\vec{d}_R\|} \right) \quad \text{(para vectores directores)} \]
\[ \theta = \arccos\left( \frac{|\vec{n}_L \cdot \vec{n}_R|}{\|\vec{n}_L\| \|\vec{n}_R\|} \right) \quad \text{(para vectores normales)}. \]

Para la recta $L_1$, el vector director es $\vec{d}_{L_1} = (2, 1)$. Para la recta $L_2$, el vector normal es $\vec{n}_{L_2} = (2, -1)$. A continuación calcularemos los ángulos Igualamos los ángulos para determinar $a$ y $b$.



\begin{minipage}{.4\textwidth}
    

\textbf{Ángulo de la recta $\bf R$ con la recta $\bf L_1$:} 

\begin{itemize}
    \item $|\vec{d}_{L_1} \cdot \vec{d}_R|= \left| (2,1)  \cdot (a,b) \rule{0mm}{3.9mm}\right| =|2a + b|$

    \item $ \|\vec{d}_{L_1}\|   = \|(2,1) \| =  \sqrt{2^2 + 1^2} = \sqrt{ 5\,}$

    \item $\|\vec{d}_R  \| = \sqrt{a^2 + b^2} $

    
\end{itemize}



\begin{align*}
    \theta &= \arccos\left( \frac{|\vec{d}_{L_1} \cdot \vec{d}|}{\|\vec{d}_{L_1}\| \|\vec{d}\|} \right) 
    \\[2mm]
    &= \arccos\left( \frac{|2a + b|}{\sqrt{5\,} \sqrt{a^2 + b^2}}  \right) 
\end{align*}


   \end{minipage}\qquad\qquad \begin{minipage}{.53\textwidth}
    

\textbf{Ángulo  de la recta $\bf R$ con la recta $L_2$:} 

\begin{itemize}
    \item $|\vec{n}_{L_2} \cdot \vec{n}_R|= | (2,-1)\cdot (-b,a) \rule{0mm}{3.9mm}|=|-2 \, b + (-1)a|$ \\  \, \qquad \, $= |-2 \, b -a|  $

    \item $\|\vec{n}_{L_2}\|= \| (2,-1) \| = \sqrt{2^2 + (-1)^2} = \sqrt{5} $

    \item $\|\vec{n}_{R}\|= \| (-b,a) \| = \sqrt{(-b)^2 + a^2} = \sqrt{b^2+a^2} $

    
\end{itemize}
\begin{align*}
    \theta &= \arccos\left( \frac{|\vec{n}_{L_2} \cdot \vec{n}_R|}{\|\vec{n}_{L_2}\| \|\vec{n}\|} \right)
    \\[2mm]
    &= \arccos\left( \frac{|-2 \, b -a|}{\sqrt{5} \sqrt{b^2 + a^2}}  \right)
\end{align*}


\end{minipage}




Igualando estas dos ecuaciones, obtenemos valores para $a$ y $b$ que determinan el vector director de la recta buscada.

\begin{align*}
    \Cancelar[blue]{\arccos }\left( \frac{|2a + b|}{\sqrt{5\,} \sqrt{a^2 + b^2\,}}  \right) 
    &=
   \Cancelar[blue]{ \arccos }\left( \frac{|-2 \, b - a|}{\sqrt{5} \sqrt{b^2 + a^2\,}}  \right)
    \\[3mm]
      \frac{|2a + b|}{\Cancelar[blue]{\sqrt{5\,}}\Cancelar[red]{\sqrt{a^2 + b^2\,}}}
    &=
    \frac{|-2\, b - a|}{\Cancelar[blue]{\sqrt{5}}\Cancelar[red]{ \sqrt{a^2 + b^2\,}}}
    \\[3mm]
       \left|2a + b \rule{0mm}{4mm} \right|    
    &=
      \left|    -2 \, b - a  \rule{0mm}{4mm} \right|
\end{align*}

Si dos expresiones tienen el mismo valor absoluto, quiere decir que esas dos expresiones son iguales o tienen signo opuesto.
%
%

\vspace{-9mm}\begin{multicols}{2}  
\begin{align*}
\textbf{Caso 1:}\qquad   2 a + b &= -2b - a  
  \\[3mm]
   2a +a   &= -2b -  b
   \\[3mm]
   3a &= -3 b
   \\[3mm]
   a &= -b
\end{align*} 
  
 \columnbreak % Forzar cambio de columna si es necesario 

\begin{align*}
\textbf{Caso 2:}\qquad    2a + b &= -(-2b - a ) 
  \\[3mm]
  2a + b &= 2b  + a 
  \\[3mm]
   2a - a   &= 2b -  b
   \\[3mm]   
   a &= b
\end{align*}

\end{multicols}

 Por lo tanto, un vector director es $(a, a)$ con $a \in \mathbb{R}$, y otro vector director es $(a, -a)$ con $a \in \mathbb{R}$. En cada caso, tomamos $a=1$ para obtener un representante fijo para cada dirección.



\textbf{Resultados}
Las rectas que cumplen con las condiciones son:
\[ \text{Recta 1: } (x, y) = (3, -2) + \lambda (1, 1), \quad \lambda \in \mathbb{R}, \]
\[ \text{Recta 2: } (x, y) = (3, -2) + \lambda (1, -1), \quad \lambda \in \mathbb{R}. \]

\textbf{Gráfico de las rectas}



\begin{center}
    
\definecolor{gris}{rgb}{0.5,0.5,0.5}
\begin{tikzpicture}[scale=0.6]
%%%% punto% 
\coordinate (A) at (3,-2);
% extremos para los ejes
\xdef\ext{9}
\xdef\yext{4}
\xdef\f{-5}
\xdef\g{-8}
\pgfmathtruncatemacro{\c}{\f+1}
\pgfmathtruncatemacro{\d}{\g+1}
\draw [color=gris!30!,dash pattern=on 1pt off 1pt, xstep=1cm,ystep=1cm] (\f-.3,\g-.3) grid (\ext+.3,\yext+.3);

\draw[->] (\f-.3,0) -- (\ext+.3,0) ;

\draw[->] (0,\g-.3) -- (0,\yext+.3) ;

\foreach \x in {\f,...,-1,1,2,...,\ext}
\draw[shift={(\x,0)},color=black] (0pt,2pt) -- (0pt,-2pt) node[below] {\tiny $\x$};

\foreach \y in {\g,...,-1,1,2,...,\yext }
\draw[shift={(0,\y)},color=black] (2pt,0pt) -- (-2pt,0pt) node[left] {\tiny $\y$};

        % Recta L1
        \draw[domain=-1.15:6.6, thick, blue, variable=\t] plot ({-3+2*\t},{-2+\t}) node[right]   {$L_1 : \frac{x + 3}{2} = y + 2$};

        % Recta L2
        \draw[domain=-3.1:3.1, thick, red] plot ({\x},{(2*\x - 2)}) node[above right]   {$L_2 : 2x - y - 2 = 0$};

        % Recta 1
        \draw[domain=-6.15:7.15, thick, green!70!black] plot ({3+\x},{-2+\x}) node[above]  {Recta 1};

        % Recta 2
        \draw[domain=-6.15:6.15, thick, orange] plot ({3+\x},{-2-\x}) node[right]  {Recta 2};

        % Punto 
\draw[blue,fill=blue] (A) circle (2pt) node[below]{$A = (3,-2)$} ; 


% Ángulos 
\coordinate (AA) at (-3,-8);
\coordinate (B) at ($(AA)+(1,2)$);
\coordinate (C) at ($(AA)+(1,1)$);

\pic[draw, color=magenta, fill=magenta!50!,   angle eccentricity=1.6, angle radius=.7cm, "$ \theta_2$",opacity=.8] {angle = C--AA--B};

% Ángulos 
\coordinate (AA) at (9,4);
\coordinate (B) at ($(AA)-(2,1)$);
\coordinate (C) at ($(AA)-(1,1)$);

\pic[draw, color=magenta, fill=magenta!50!,   angle eccentricity=1.6, angle radius=.7cm, "$ \theta_2$",opacity=.8] {angle = B--AA--C};

% Ángulos % Ángulos 
\coordinate (AA) at (1,0);
\coordinate (B) at ($(AA)+(2,1)$);
\coordinate (C) at ($(AA)+(1,-1)$);

\pic[draw, color=violet, fill=violet!50!,   angle eccentricity=1.4, angle radius=.6cm, "$ \theta_1$",opacity=.8] {angle = C--AA--B};

% Ángulos % Ángulos 
\coordinate (AA) at (1,0);
\coordinate (B) at ($(AA)-(1,2)$);
\coordinate (C) at ($(AA)+(1,-1)$);

\pic[draw, color=violet, fill=violet!50!,   angle eccentricity=1.4, angle radius=.4cm, "$ \theta_1$",opacity=.8] {angle = B--AA--C};


\end{tikzpicture}


\end{center}




    
\end{enumerate}














    

\item 
\begin{itemize}
    \item[(a)] Calcular el área del triángulo cuyos vértices son $A = (1, 2)$, $B = (4, -1)$ y $C = (5, 2)$.
    \item[(b)] Hallar los vértices de un triángulo isósceles sabiendo que la base está sobre la recta $L : x - y + 1 = 0$, un vértice es $V = (-3, 4)$ y el área del triángulo es 6.
\end{itemize}




    \textbf{Solución:}

\begin{enumerate}[(a) ] 
\item \textbf{  Cálculo del área del triángulo.}\\





\begin{minipage}{.43\textwidth}

Dado el triángulo con vértices \( A = (1,2) \), \( B = (4,-1) \) y \( C = (5,2) \), calculamos su área utilizando como base el segmento \( \overline{AC} \) y determinando la altura correspondiente desde el punto \( B \) hasta la recta que contiene a \( AC \). En el gráfico, tanto la base como la altura están claramente señaladas.




\end{minipage}\qquad\begin{minipage}{.47\textwidth}

\begin{center}
    \begin{tikzpicture}[scale=.5]
        \draw[->] (-1,0) -- (9,0) node[right] {$x$};
        \draw[->] (0,-2) -- (0,3.9) node[above] {$y$};
        
        \coordinate (A) at (1,2);
        \coordinate (B) at (4,-1);
        \coordinate (C) at (5,2);
        \coordinate (D) at (4,2);
        
        \draw[thick] (A) -- (B) -- (C) --node[above,pos=.65] {$b$} cycle; 
        \filldraw[black] (A) circle (2pt) node[above left] {$A=(1,2)$};
        \filldraw[black] (B) circle (2pt) node[below right] {$B=(4,-1)$};
        \filldraw[black] (C) circle (2pt) node[above right] {$C=(5,2)$};
        \filldraw[black] (D) circle (2pt);
        \draw[dashed] (B) -- (D) node[left,pos=.65] {$h$};
    \end{tikzpicture}
\end{center}
\end{minipage}




\textbf{Ecuación de la recta $\overleftrightarrow{ AC}$}
\begin{multicols}{2}
La pendiente de la recta $\overleftrightarrow{ AC}$ es:
\[
m = \frac{2 - 2}{5 - 1} = \frac{0}{4} = 0.
\]
Dado que la pendiente es cero, la ecuación de la recta $\overleftrightarrow{ AC}$ es:
\[
y = 2.
\] 
\end{multicols}

\textbf{Cálculo de la altura desde $B$ a la recta $AC$}

\bigskip

\begin{minipage}{.4\textwidth}
    

La distancia de un punto $(x_0, y_0)$ a una recta de la forma $ax + by + c = 0$ se calcula como:
\[
d = \frac{|ax_0 + by_0 + c|}{\sqrt{a^2 + b^2}}.
\]
\end{minipage}\quad\qquad\begin{minipage}{.47\textwidth}
    
Sustituyendo $B(4,-1)$ y la ecuación \\ $y = 2 \qquad\Longrightarrow\qquad 0x + 1y - 2 = 0$:
\[
d = \frac{|0(4) + 1(-1) - 2|}{\sqrt{0^2 + 1^2}} = \frac{|-1 - 2|}{1} = \frac{|-3|}{1} = 3.
\]

\end{minipage}

\bigskip 

\textbf{Cálculo del área}

\bigskip



\begin{minipage}{.6\textwidth}

La base es la distancia $AC$:
\[
AC = \sqrt{(5 - 1)^2 + (2 - 2)^2} = \sqrt{4^2 + 0} = 4.
\]
La altura es la distancia de $B$ a la recta $AC$, que es $3$.
\end{minipage}\quad\qquad\begin{minipage}{.3\textwidth}


El área se calcula como:
\begin{align*}
\text{Area} &= \frac{1}{2} \times AC \times h 
\\
&= \frac{1}{2} \times 4 \times 3 
\\
&= \frac{12}{2} = 6.
\end{align*} 
\end{minipage}

Por lo tanto, el área del triángulo es $6$.




\begin{comment}
    
\bigskip 

{\bf Otra forma:}


\begin{minipage}{.45\textwidth}

Dado el triángulo con vértices $ A = (1,2) $, $ B = (4,-1) $ y $ C = (5,2) $, calculamos su área utilizando el segmento $ \overline{AB} $ como base y determinando la altura desde $ C $ hasta la recta que contiene a $ AB $. Como se observa en el gráfico, la altura $ h $ está sobre la perpendicular al segmento $ \overline{AB} $ y pasa por $C $.



\end{minipage}\qquad\begin{minipage}{.45\textwidth}

\begin{center}
    \begin{tikzpicture}[scale=.5]
        \draw[->] (-1,0) -- (9,0) node[right] {$x$};
        \draw[->] (0,-2) -- (0,3.9) node[above] {$y$};
        
        \coordinate (A) at (1,2);
        \coordinate (B) at (4,-1);
        \coordinate (C) at (5,2);
        \coordinate (D) at (3,0);
        
        \draw[thick] (A) -- (B) -- (C) -- cycle; 
        \filldraw[black] (A) circle (2pt) node[above left] {$A=(1,2)$};
        \filldraw[black] (B) circle (2pt) node[below right] {$B=(4,-1)$};
        \filldraw[black] (C) circle (2pt) node[above right] {$C=(5,2)$};
        \filldraw[black] (D) circle (2pt);
        \draw[dashed] (C) -- (D) node[left,pos=.45] {$h$};
    \end{tikzpicture}
\end{center}
\end{minipage}


\medskip 


\setlength{\columnsep}{9mm}
\begin{multicols}{2}
\noindent $\bullet$
\textbf{Ecuación de la recta que pasa  $A$ y $B$.}


La pendiente de la recta $AB$ es:
\[
m = \frac{-1 - 2}{4 - 1} = \frac{-3}{3} = -1.
\]

    
Usando la ecuación punto-pendiente con $A=(1,2)$:
\[
y - 2 = -1(x - 1)
\qquad
\Rightarrow 
\qquad 
y = -x + 3.
\] 

 \columnbreak % Forzar cambio de columna si es necesario

 
\noindent $\bullet$ \textbf{Ecuación de la recta perpendicular a $AB$ que pasa por $C$}
La pendiente perpendicular es $m_\perp = 1$. Usamos la ecuación punto-pendiente con $C=(5,2)$:
\[
y - 2 = 1(x - 5) 
\quad 
\Rightarrow 
\quad 
y = x - 3.
\]
\end{multicols}    



\subsubsection*{Intersección de ambas rectas}
Igualamos las ecuaciones:
\[
-x + 3 = x - 3.
\]
Resolviendo para $x$:
\[
3 + 3 = x + x \Rightarrow 6 = 2x \Rightarrow x = 3.
\]
Sustituyendo en $y = -x + 3$:
\[
y = -3 + 3 = 0.
\]
Por lo tanto, la intersección es $D = (3,0)$ y este punto es el pie de la altura.

\subsubsection*{Cálculo del área}
\begin{multicols}{3}

La base es la distancia $AB$:
\begin{align*}
  d(A,B) &= \sqrt{(4 - 1)^2 + (-1 - 2)^2} 
    \\
    &= \sqrt{9 + 9} = \sqrt{18} = 3\sqrt{2}.
\end{align*}

 
La altura es la distancia $CD$:
\begin{align*}
d(C,D) &= \sqrt{(5 - 3)^2 + (2 - 0)^2} 
\\
&= \sqrt{4 + 4} = \sqrt{8} = 2\sqrt{2}.
\end{align*} 
El área se calcula como:
\begin{align*}
     \text{Area} 
     &= \frac{1}{2} \times CD \times AB 
    \\
    &= \frac{1}{2} \times 2\sqrt{2} \times 3\sqrt{2} 
    \\
    &= \frac{1}{2} \times 12 = 6.
\end{align*}

Por lo tanto, el área del triángulo es $6$.

\begin{center}
    \begin{tikzpicture}[scale=.5]
        \draw[->] (-1,0) -- (6,0) node[right] {$x$};
        \draw[->] (0,-2) -- (0,3) node[above] {$y$};
        \coordinate (A) at (1,2);
        \coordinate (B) at (4,-1);
        \coordinate (C) at (5,2);
        \coordinate (D) at (3,0);
        
        \draw[thick] (A) -- node[below] {$b$}(B) -- (C) --  cycle;
        \filldraw[black] (A) circle (2pt) node[above left] {$A=(1,2)$};
        \filldraw[black] (B) circle (2pt) node[below] {$B=(4,-1)$};
        \filldraw[black] (C) circle (2pt) node[above right] {$C=(5,2)$};
        \filldraw[black] (D) circle (2pt);
        \draw[dashed] (C) --node[below] {$h$} (D);
    \end{tikzpicture}
\end{center}
\end{multicols}


\end{comment}









\begin{tcolorbox}[colback=white,colframe=red!60!black,title= \textbf{Observación:}]

Para calcular el área de un triángulo dados tres vértices $
A(x_1, y_1), \quad B(x_2, y_2), \quad C(x_3, y_3)$:
\begin{enumerate}
    \item Elegir un lado (por ejemplo \( AC \)) como base.
    \item Calcular su longitud con la fórmula de distancia:
    \[
    \text{base} = \sqrt{(x_3 - x_1)^2 + (y_3 - y_1)^2}.
    \]
    \item Obtener la ecuación de la recta que contiene esa base, es decir considerar $L$ la recta que pasa por los puntos $A$ y $C$. 
    \item Calcular la distancia perpendicular desde el vértice opuesto (por ejemplo \( B \)) a dicha recta:
    \[
    d = \frac{|ax_0 + by_0 + c|}{\sqrt{a^2 + b^2}},
    \]
    donde \( ax + by + c = 0 \) es la ecuación de la recta que pasa por los puntos $A$ y $C$. 
    \item El área se obtiene como:
    \[
    S = \frac{1}{2} \times \text{base} \times \text{altura}.
    \]
\end{enumerate}

\end{tcolorbox}


\item \textbf{ Hallar los vértices de un triángulo isósceles} \\
 



\begin{minipage}{.53\textwidth} 



Dado que la base del triángulo se encuentra sobre la recta \( L: x - y + 1 = 0 \) y que uno de sus vértices es \( V = (-3,4) \), además de conocer que su área es 6, podemos analizar su altura. Al trazarla, observamos que su pie divide la base en dos partes iguales, como se muestra en el gráfico adjunto. Por otro lado, sabemos que la altura del triángulo es igual a la distancia entre \( V \) y la recta \( L \) y juntando con la fórmula del área de un triángulo, podemos plantear la siguiente ecuación:

\begin{align*}
h= d(V,L) , \qquad \text{Area} &= \frac{b\cdot h}{2}
    \\
\text{despejando: \qquad }    \frac{6\cdot 2}{h} &=  b
\\
\frac{12}{h} &=  b
\end{align*}





\end{minipage}\qquad\begin{minipage}{.37\textwidth}

\begin{center}
    \begin{tikzpicture}[scale=.7]
        %\draw[->] (-3,0) -- (3,0) node[right] {$x$};
        %\draw[->] (0,-2) -- (0,5) node[above] {$y$};
        
        \coordinate (V) at (-3,4);
        \coordinate (B) at (1,2);
        \coordinate (C) at (-1,0);
        \coordinate (D) at (0,1);

        % Recta L2
        \draw[domain=-3.1:3.1, thick, red] plot ({\x},{\x + 1}) node[above ]   {$L : x - y + 1 = 0$};
        
        \draw[thick] (V) -- (B) --node[below right,pos=.35] {$b$} (C) --cycle; 
        \filldraw[black] (V) circle (2pt) node[above left] {$V=(-3,4)$};
        %\filldraw[black] (B) circle (1pt) node[above right] {\scriptsize$B=\left(\frac{1}{2},\frac{3}{2}\right)$};
        %\filldraw[black] (C) circle (1pt) node[below right] {\scriptsize$C=\left(-\frac{1}{2},\frac{1}{2}\right)$};
        \filldraw[black] (D) circle (1pt); 


        
        \draw[dashed] (V) -- (D) node[left,pos=.65] {$h$};


        \pgfmathsetmacro{\rr}{sqrt(2)}
        \begin{scope}
        \clip (-1,-1) --(-2,0)--(1,3)--(2,2)--cycle;
        \draw[dashed] (D) circle (\rr cm);
        \end{scope}
        
    \end{tikzpicture}
\end{center}
\end{minipage}



Así, podemos determinar el valor de la base y ubicar los vértices del triángulo en el plano. Es que de los vértices sabemos distan del pie de la altura la mitad de la base y estan a una distancia $\frac{b}{2}$.


\bigskip




\subsubsection*{Ecuación de la recta perpendicular a $L$ por $V$}
La pendiente de la recta $L$ es $m = 1$, por lo que la pendiente de la recta perpendicular es $m'= -1$.
La ecuación de la recta perpendicular que pasa por $V$ es:
\[
y - 4 = -1(x + 3) \Rightarrow y = -x + 1.
\]

\subsubsection*{Intersección de la recta perpendicular con $L$}
\begin{multicols}{3}
Igualamos las ecuaciones:
\[
x + 1 = -x + 1.
\]
Resolviendo:
\[
2x + 1 = 1 \Rightarrow 2x = 0 \Rightarrow x = 0.
\]
Sustituyendo en $y = x + 1$:
\[
y = 0 + 1 = 1.
\] 

\end{multicols}

Por lo tanto, la intersección es $H = (0,1)$ y este es el pie de la altura.

\subsubsection*{Cálculo de la distancia de $V$ a $L$}
Usamos la fórmula de distancia punto-recta:
\[
d = \frac{|(-3) - 4 + 1|}{\sqrt{1^2 + (-1)^2}} = \frac{|-6|}{\sqrt{2}} = \frac{6}{\sqrt{2}}.
\]
Entonces tenemos que la altura   $h = \frac{6}{\sqrt{2}}$, entonces la altura vale: $b= \frac{12}{h} = \frac{12}{\frac{6}{\sqrt{2}}} = \frac{12\sqrt{2}}{6}= 2 \sqrt{2} $



\bigskip


\textbf{Los vértices restantes  distan del pie de la altura $\frac{b}{2}$:}



\bigskip



Los vértices están sobre la recta $L : x - y + 1 = 0$ entonces verifica su ecuación, consideremos el vértice $ A = (x_A, y_A) $ tenemos: 

\[
x_A -y_A +1 = 0 
\quad 
\Longrightarrow
\quad
x_A  + 1=  y_A   
\qquad
\qquad
\therefore A = \left( x_A, x_A  + 1 \right)
\]

\begin{align*}
    d(A , H) 
    &= d \left(\rule{0mm}{4mm} (x_A, x_A  + 1 ) , (0,1)\right)
    \\[4mm]
    \frac{b}{2}
    &= \sqrt{x_A^2 + (x_A+1 -1)^2 \rule{0mm}{3.5mm}}
    \\[4mm]
    \frac{\Cancelar[red]{2}\sqrt{2}}{\Cancelar[red]{2}}
    &= \sqrt{ 2x_A^2  \rule{0mm}{3.5mm}}
    \\[4mm]
    2
    &=
    2x_A^2
    \\[4mm]
    1
    &=
    x_A^2
    \\[4mm]
    x_A = -1 \qquad
    & \text{o}\qquad
    x_A= 1
\end{align*}

Sustituyendo $y = x + 1$ tenemos:


\begin{itemize}
    \item Si $x_A=-1$ tenemos que $y_A=  0 $  y el vértice es $A_1= (-1,0)$

    \item Si $x_A= 1$ tenemos que $y_A= 2$ y el vértice es $A_2= (1,2)$
\end{itemize}

\begin{center}
    \definecolor{gris}{rgb}{0.5,0.5,0.5}
\begin{tikzpicture}[scale=1.1]
% extremos para los ejes
\xdef\ext{3}
\xdef\yext{5}
\xdef\f{-3}
\xdef\g{-2}
\pgfmathtruncatemacro{\c}{\f+1}
\pgfmathtruncatemacro{\d}{\g+1}
\draw [color=gris!30!,dash pattern=on 1pt off 1pt, xstep=1cm,ystep=1cm] (\f-.3,\g-.3) grid (\ext+.3,\yext+.3);

\draw[->] (\f-.3,0) -- (\ext+.3,0) ;

\draw[->] (0,\g-.3) -- (0,\yext+.3) ;

\foreach \x in {\f,\c,...,-2,2,4,...,\ext}
\draw[shift={(\x,0)},color=black] (0pt,2pt) -- (0pt,-2pt) node[below] {\tiny $\x$};

\foreach \y in {\g,\d,...,-2,2,4,...,\yext }
\draw[shift={(0,\y)},color=black] (2pt,0pt) -- (-2pt,0pt) node[left] {\tiny $\y$};
        %\draw[->] (-3,0) -- (3,0) node[right] {$x$};
        %\draw[->] (0,-2) -- (0,5) node[above] {$y$};
        
        \coordinate (V) at (-3,4);
        \coordinate (B) at (1,2);
        \coordinate (C) at (-1,0);
        \coordinate (D) at (0,1);

        % Recta L2
        \draw[domain=-3.1:3.1, thick, red] plot ({\x},{\x + 1}) node[above ]   {$L : x - y + 1 = 0$};
        
        \draw[thick] (V) -- (B) --node[below right,pos=.35] {$b$} (C) --cycle; 
        \filldraw[black] (V) circle (2pt) node[above left] {$V=(-3,4)$};
        \filldraw[black] (B) circle (1pt) node[ right] {\scriptsize$B=\left(1,2\right)$};
        \filldraw[black] (C) circle (1pt) node[below right] {\scriptsize$C=\left(-1,0\right)$};
        \filldraw[black] (D) circle (1pt); 


        
        \draw[dashed] (V) -- (D) node[left,pos=.65] {$h$};


        \pgfmathsetmacro{\rr}{sqrt(2)}
        \begin{scope}
        \clip (-1,-1) --(-2,0)--(1,3)--(2,2)--cycle;
        \draw[dashed] (D) circle (\rr cm);
        \end{scope}
        
    \end{tikzpicture}
\end{center}


\end{enumerate}












\item En cada inciso, encontrar la ecuación del haz de rectas determinado por las condiciones pedidas y graficar tres elementos de la familia indicando el valor del parámetro:

\begin{itemize}
\begin{multicols}{2}
    \item[(a)] Paralelas a la recta $L : 3y - x - 3 = 0$.
    \item[(b)] Perpendicular a la recta $L : 3x + 2y - 7 = 0$.
    
    \item[(c)] Pasan por el punto $P = (4, -2)$.
    \item[(d)] Paralelas al eje $x$.
\end{multicols}
\end{itemize}





\textbf{Solución}

\textbf{(a) Rectas paralelas a $L : 3y - x - 3 = 0$}

La ecuación dada se reescribe en la forma general:
\[ 3y - x - 3 = 0 \Rightarrow y = \frac{1}{3}x + 1. \]

Las rectas paralelas a esta tienen la misma pendiente, es decir:
\[ y = \frac{1}{3}x + k, \quad k \in \mathbb{R}. \]

Gráfico:
\begin{center}
    \begin{tikzpicture}[scale=0.8]
        \draw[->] (-5,0) -- (5,0) node[right] {$x$};
        \draw[->] (0,-5) -- (0,5) node[above] {$y$};
        \draw[domain=-.9:2.9, smooth, variable=\y, color=red] plot ({3*\y-3},{\y}) node[right] {$L : 3y -x- 3 = 0$};

        \foreach \k in {-2,0,2} {
            \draw[domain=-5:5, smooth, variable=\x, color=blue] plot ({\x},{(1/3)*\x+\k}) node[right] {$k=\k$};
        }
    \end{tikzpicture}
\end{center}

\textbf{(b) Rectas perpendiculares a $L : 3x + 2y - 7 = 0$}

La pendiente de la recta dada es:
\[ 3x + 2y - 7 = 0 \qquad \Rightarrow \qquad  y = -\frac{3}{2}x + \frac{7}{2}. \]

Las rectas perpendiculares tienen pendiente inversa y opuesta:
\[ y = \frac{2}{3}x + k, \quad k \in \mathbb{R}. \]

Gráfico:
\begin{center}
    \begin{tikzpicture}[scale=0.8]
        \draw[->] (-5,0) -- (5,0) node[right] {$x$};
        \draw[->] (0,-5) -- (0,5) node[above] {$y$};

        \draw[domain=-1:5, smooth, variable=\x, color=blue] plot ({\x},{(7 - 3*\x)/2}) node[right] {$L : 3x + 2y - 7 = 0$};

        
        \foreach \k in {-2,0,2} {
            \draw[domain=-5:5, smooth, variable=\x, color=red] plot ({\x},{(2/3)*\x+\k}) node[right] {$k=\k$};
        }
    \end{tikzpicture}
\end{center}

\textbf{(c) Rectas que pasan por $P = (4,-2)$}

La ecuación de una recta que pasa por $(4,-2)$ es:
\[ y + 2 = m(x - 4), \quad m \in \mathbb{R},\qquad \text{también la recta: } x=4\]

Gráfico:
\begin{center}
    \begin{tikzpicture}[scale=0.8]
        \draw[->] (-3,0) -- (7,0) node[right] {$x$};
        \draw[->] (0,-5) -- (0,5) node[above] {$y$};
        \begin{scope}
        \clip (-5,-5.5)rectangle(11.5,5);
        \foreach \m/\t in {-3/{(-3)},-2/{(-2)},-1/{(-1)},-.5/{-\frac{1}{5}} ,.5/{\frac{1}{5}},1/1,2/2,3/3} {
            \draw[domain=-3:7, smooth, variable=\x, color=violet] plot ({\x},{\m*(\x-4)-2}) node[right] {$y= \t (x-4) -2$};
        }
        \end{scope}
        
        \draw[domain=-5:5, smooth, variable=\y, color=violet] plot (4,{\y}) node[right] {$x=4$};
        \draw[domain=-3:7, smooth, variable=\x, color=violet] plot ({\x},-2) node[right] {$y=-2$};
        \filldraw[black] (4,-2) circle (2pt) node[below left] {$P(4,-2)$};
        
    \end{tikzpicture}
\end{center}

\textbf{(d) Rectas paralelas al eje $x$}

Las rectas paralelas al eje $x$ tienen ecuación:
\[ y = k, \quad k \in \mathbb{R}. \]

Gráfico:
\begin{center}
    \begin{tikzpicture}[scale=0.8]
        \draw[->] (-5,0) -- (5,0) node[right] {$x$};
        \draw[->] (0,-5) -- (0,5) node[above] {$y$};
        \foreach \k in {-4,-3,-2,-1,1,2,3} {
            \draw[domain=-5:5, smooth, variable=\x, color=purple] plot ({\x},{\k}) node[right] {$k=\k$};
        }
    \end{tikzpicture}
\end{center}







\item Determinar la propiedad que cumplen todas las rectas que pertenecen a cada uno de los siguientes haces de rectas:

\begin{itemize}
\begin{multicols}{2}
    \item[(a)] $-3x + 3y + k = 0$.
    \item[(b)] $y = kx - 8$.

    
    \item[(c)] $y + 2 = k(x - 2)$.
    \item[(d)] $\frac{x}{-2} + \frac{y}{k} = 1$, con $k \neq 0$.
\end{multicols}
\end{itemize}
 

   
   \textbf{Solución del ejercicio}

\textbf{Análisis de cada haz de rectas}

\begin{enumerate}
    \item[(a)] \textbf{Ecuación del haz:} $-3x + 3y + k = 0$ 

\begin{minipage}{.45\textwidth}

    \begin{center}
    \begin{tikzpicture}[scale=01,>=triangle 45]
        \draw[->] (-4,0) -- (3,0) node[below] {$x$};
        \draw[->] (0,-2) -- (0,4) node[left] {$y$};
        \coordinate (A) at (0,0);%punto a la izquierda
        \coordinate (B) at (-3,3);%punto a la derecha
        % vector director
        \draw[->, thick, purple,line width=1pt] (A) -- (B) node[midway, above right,pos=.65] {$\vec{n}$};

    
        \begin{scope}
        \clip (-3,-2)rectangle(7,4);
        \foreach \k in {-2,-1,0,1,2} {
            \draw[domain=-3:3,smooth,variable=\x,blue] plot ({\x},{\x-\k/3}) node[right] {$k=\k$};
        }
        \end{scope}
        
    \end{tikzpicture}
    \end{center}
    \end{minipage}\qquad\begin{minipage}{.45\textwidth}

         \textbf{Propiedad:} Todas las rectas del haz tienen el mismo vector normal, por lo que son \textbf{paralelas entre sí}.    
\end{minipage}     
     
    \item[(b)] \textbf{Ecuación del haz:} $y = kx - 8$
\begin{multicols}{2}
    \begin{itemize}
        \item si $k=-2$ tenemos $y = -2 x - 8$
        \item si $k=-1$ tenemos $y = -x - 8$
        \item si $k=0$ tenemos $y =  - 8$
        \item si $k=1$ tenemos $y = x - 8$        
        \item si $k=1$ tenemos $y = x - 8$
    \end{itemize}
todas las expreciones tienen en común que tienen ordenada al origen $-8$.

    
    
    \begin{center}
    \begin{tikzpicture}[scale=0.5]
        \draw[->] (-5,0) -- (5,0) node[below] {$x$};
        \draw[->] (0,-10) -- (0,2) node[left] {$y$};
        \begin{scope}
        \clip (-5,-10)rectangle(7,4);
        \foreach \k in {-2,-1,0,1,2} {
            \draw[domain=-5:5,smooth,variable=\x,red] plot ({\x},{\k*\x-8});
        }
        \end{scope}
        
        \fill[black] (0,-8) circle (2pt) node[below left] {$(0,-8)$};
    \end{tikzpicture}
    \end{center}

\end{multicols}

\textbf{Propiedad:} \textbf{Todas las rectas pasan por el punto $(0,-8)$}.


\bigskip

\begin{minipage}{.45\textwidth}
    

    \item[(c)] \textbf{Ecuación del haz:} $y + 2 = k(x - 2)$ 
    
    \textbf{Propiedad:} Todas las rectas pasan por el punto $(2,-2)$, excepto la recta vertical que pasa por $(2,-2)$.

    \end{minipage}\qquad\begin{minipage}{.45\textwidth}

    \begin{center}
    \begin{tikzpicture}[scale=0.5]
        \draw[->] (-5,0) -- (5,0) node[below] {$x$};
        \draw[->] (0,-5) -- (0,5) node[left] {$y$};
        \begin{scope}
        \clip (-5,-5)rectangle(7,5);
        \foreach \k in {-2,-1,0,1,2} {
            \draw[domain=-5:5,smooth,variable=\x,green] plot ({\x},{\k*\x-2*\k-2});
        }
        \end{scope}
        \fill[black] (2,-2) circle (2pt) node[below left] {$(2,-2)$};
    \end{tikzpicture}
    \end{center}
    
\end{minipage}

\bigskip


\begin{minipage}{.45\textwidth}

    \item[(d)] \textbf{Ecuación del haz:} $\frac{x}{-2} + \frac{y}{k} = 1$ % \qquad \longrightarrow \qquad \frac{y}{k} = \frac{x}{2} + 1 \qquad\longrightarrow \qquad y = \frac{k}{2}x + k$
    %\\\textbf{Forma explicita:} $y = \frac{k}{2}x + k$
    esta en su forma segmentaría,entonces sabemos que cortan el eje $y$ en $(0,k)$ y el eje $x$ en $(-2,0)$


    \end{minipage}\qquad\begin{minipage}{.45\textwidth}    
    
    \begin{center}
    \begin{tikzpicture}[scale=0.5]
        \draw[->] (-5,0) -- (5,0) node[below] {$x$};
        \draw[->] (0,-5) -- (0,5) node[left] {$y$};
        \begin{scope}
        \clip (-5,-5)rectangle(5,5);
        \foreach \k in {-2,-1,1,2,3} {
            \draw[domain=-5:5,smooth,variable=\x,purple] plot ({\x},{\k/2*\x+\k});
        }
        \end{scope}
        \fill[black] (-2,0) circle (2pt) node[below left] {$(-2,0)$};
    \end{tikzpicture}
    \end{center}
\end{minipage}


\end{enumerate}
   




\end{enumerate}



 




\end{document}